\input{preamble}

% OK, start here.
%
\begin{document}

\title{Exemples}


\maketitle

\phantomsection
\label{section-phantom}

\tableofcontents

\section{Introduction}
\label{section-introduction}

\noindent
Ce chapitre contiendra des exemples qui éclairent la théorie.





\section{Une limite vide}
\label{section-empty-limit}

\noindent
Cet exemple est dû à Waterhouse, voir \cite{Waterhouse}.
Soit $S$ un ensemble non dénombrable. Pour toute partie finie
$T \subset S$, considérons l'ensemble $M_T$ des applications injectives $T \to \mathbf{N}$.
Si $T \subset T' \subset S$ sont finies, la restriction $M_{T'} \to M_T$
est surjective. Nous obtenons ainsi un système projectif indexé par
l'ensemble ordonné filtrant des parties finies de $S$,
dont les applications de transition sont surjectives.
Mais $\lim M_T = \emptyset$, car un élément de la limite
définirait une application injective $S \to \mathbf{N}$.



\section{Une limite nulle}
\label{section-zero-limit}

\noindent
Soit $(S_i)_{i \in I}$ un système projectif d'ensembles non vides
indexé par un ensemble ordonné filtrant, à applications de transition surjectives,
et tel que $\lim S_i = \emptyset$, voir la section \ref{section-empty-limit}.
Soit $K$ un corps et posons
$$
V_i = \bigoplus\nolimits_{s \in S_i} K
$$
Alors les applications de transition $V_i \to V_j$ sont surjectives pour $i \geq j$.
Cependant, $\lim V_i = 0$. En effet, si $v = (v_i)$ est un élément de la
limite, le support de $v_i$ serait une partie finie $T_i \subset S_i$
telle que $\lim T_i \not = \emptyset$, voir
Catégories, lemme \ref{categories-lemma-nonempty-limit}.

\medskip\noindent
Pour tout $i$, considérons l'unique application $K$-linéaire $V_i \to K$ qui envoie
chaque vecteur de base $s \in S_i$ sur $1$. Soit $W_i \subset V_i$
son noyau. Alors
$$
0 \to (W_i) \to (V_i) \to (K) \to 0
$$
est une suite exacte courte non scindée de systèmes projectifs d'espaces vectoriels
sur l'ensemble ordonné filtrant $I$. Ainsi,
$W_i$ forme un système projectif de $K$-espaces vectoriels
à applications de transition surjectives, de limite nulle et dont
$R^1\lim$ ne s'annule pas.



\section{Une limite projective non quasi-compacte d'espaces quasi-compacts}
\label{section-lim-not-quasi-compact}

\noindent
Désignons par $\mathbf{N}$ l'ensemble des entiers naturels.
Pour tout entier $n$, désignons par $I_n$ l'ensemble des entiers naturels $> n$.
Définissons $T_n$ comme l'unique topologie sur $\mathbf{N}$ ayant pour base
$\{1\}, \ldots , \{n\}, I_n$.  Notons $X_n$ l'espace topologique
$(\mathbf{N}, T_n)$.  Pour tous $m < n$, l'application identité
$$
f_{n, m} : X_n \longrightarrow X_m
$$
est continue.  Il est clair que, pour $m < n < p$, la composée
$f_{p, n} \circ f_{n, m}$ est égale à $f_{p, m}$.  Ainsi, $((X_n), (f_{n,m}))$
est un système projectif d'espaces topologiques quasi-compacts.

\medskip\noindent
Soit $T$ la topologie discrète sur $\mathbf{N}$, et soit $X$
l'espace $(\mathbf{N}, T)$. Pour tout entier $n$, l'application identité
$$
f_n : X \longrightarrow X_n
$$
est continue. Nous affirmons qu'il s'agit de la limite projective du système
ci-dessus. Soit $(Y, S)$ un espace topologique quelconque. Pour tout entier $n$,
soit
$$
g_n : (Y, S) \longrightarrow (\mathbf{N}, T_n)
$$
une application continue. Supposons que, pour tous $m < n$, nous ayons
$f_{n,m} \circ g_n = g_m$, c'est-à-dire que le système $(g_n)$ soit compatible
avec le système projectif ci-dessus. En particulier, toutes les applications ensemblistes
$g_n$ sont égales à une même application ensembliste
$$
g : Y \longrightarrow \mathbf{N}.
$$
De plus, pour tout entier $n$, puisque $\{n\}$ est ouvert dans $X_n$,
$g^{-1}(\{n\}) = g_n^{-1}(\{n\})$ est également ouvert dans $Y$.
Par conséquent, l'application ensembliste $g$ est continue pour la topologie $S$ sur $Y$
et la topologie $T$ sur $\mathbf{N}$. Ainsi, $(X, (f_n))$ est la limite projective
du système projectif ci-dessus.

\medskip\noindent
Cependant, $X$ n'est manifestement pas quasi-compact, puisque le recouvrement ouvert
infini par les parties réduites à un point n'admet aucun sous-recouvrement fini.

\begin{lemma}
\label{lemma-lim-not-quasi-compact}
Il existe un système projectif d'espaces topologiques quasi-compacts
indexé par $\mathbf{N}$ dont la limite projective n'est pas quasi-compacte.
\end{lemma}

\begin{proof}
Voir la discussion qui précède.
\end{proof}






\section{Le faisceau structural sur le produit fibré}
\label{section-silly}

\noindent
Soient $X, Y, S, a, b, p, q, f$ comme dans l'introduction de
Catégories dérivées de schémas,
section \ref{perfect-section-kunneth}. Diagramme :
$$
\xymatrix{
& X \times_S Y \ar[ld]^p \ar[rd]_q \ar[dd]^f \\
X \ar[rd]_a & & Y \ar[ld]^b \\
& S
}
$$
Nous disposons alors d'un homomorphisme canonique
$$
can :
p^{-1}\mathcal{O}_X \otimes_{f^{-1}\mathcal{O}_S} q^{-1}\mathcal{O}_Y
\longrightarrow
\mathcal{O}_{X \times_S Y}
$$
qui n'est pas un isomorphisme en général.

\medskip\noindent
Par exemple, soient $S = \Spec(\mathbf{R})$, $X = \Spec(\mathbf{C})$ et
$Y = \Spec(\mathbf{C})$. Alors
$X \times_S Y = \Spec(\mathbf{C}) \amalg \Spec(\mathbf{C})$
est un espace discret à deux points,
et les faisceaux $p^{-1}\mathcal{O}_X$, $q^{-1}\mathcal{O}_Y$
et $f^{-1}\mathcal{O}_S$ sont les faisceaux constants de valeurs
$\mathbf{C}$, $\mathbf{C}$ et $\mathbf{R}$.
Ainsi, la source de $can$ est le faisceau constant
de valeur $\mathbf{C} \otimes_\mathbf{R} \mathbf{C}$
sur l'espace discret à deux points. Ses sections
globales sont donc de dimension $8$ comme $\mathbf{R}$-espace vectoriel,
tandis que les sections globales du but de $can$
forment $\mathbf{C} \times \mathbf{C}$,
qui est de dimension $4$ comme $\mathbf{R}$-espace vectoriel.

\medskip\noindent
Un autre exemple est le suivant. Soit $k$ un corps algébriquement
clos. Considérons
$S = \Spec(k)$, $X = \mathbf{A}^1_k$ et $Y = \mathbf{A}^1_k$.
Alors, pour tout ouvert non vide $U \subset X \times_S Y = \mathbf{A}^2_k$,
les images $p(U) \subset X = \mathbf{A}^1_k$ et $q(U) \subset \mathbf{A}^1_k$
sont ouvertes, et le lecteur vérifiera que
$$
\left(
p^{-1}\mathcal{O}_X \otimes_{f^{-1}\mathcal{O}_S} q^{-1}\mathcal{O}_Y
\right)(U) = \mathcal{O}_X(p(U)) \otimes_k \mathcal{O}_Y(q(U))
$$
Ce dernier n'est pas égal à $\mathcal{O}_{X \times_S Y}(U)$
si $U$ est le complémentaire d'une courbe irréductible $C$ dans
$X \times_S Y = \mathbf{A}^2_k$ telle que $p|_C$ et $q|_C$
soient toutes deux non constantes.

\medskip\noindent
Revenons au cas général et soit $z = (x, y, s, \mathfrak p)$
un point de $X \times_S Y$ comme dans
Schémas, lemme \ref{schemes-lemma-points-fibre-product}.
Alors, sur les fibres en $z$, l'application $can$ donne l'application
$$
can_z :
\mathcal{O}_{X, x} \otimes_{\mathcal{O}_{S, s}} \mathcal{O}_{Y, y}
\longrightarrow
\mathcal{O}_{X \times_S Y, z}
$$
Cet homomorphisme d'anneaux est plat, car le but est un anneau localisé de la
source (détails omis; indication : se ramener au cas où $X$, $Y$ et $S$
sont affines).
Notons que la source n'est pas en général un anneau local, ce qui fournit
une autre manière de voir que $can$ n'est pas un isomorphisme en général.

\medskip\noindent
Plus généralement, supposons donnés un $\mathcal{O}_X$-module $\mathcal{F}$ et un
$\mathcal{O}_Y$-module $\mathcal{G}$. Il existe alors un homomorphisme canonique
\begin{align*}
& p^{-1}\mathcal{F} \otimes_{f^{-1}\mathcal{O}_S} q^{-1}\mathcal{G} \\
& =
p^{-1}(\mathcal{F} \otimes_{\mathcal{O}_X} \mathcal{O}_X)
\otimes_{f^{-1}\mathcal{O}_S}
q^{-1}(\mathcal{O}_Y \otimes_{\mathcal{O}_Y} \mathcal{G}) \\
& =
p^{-1}\mathcal{F} \otimes_{p^{-1}\mathcal{O}_X} p^{-1}\mathcal{O}_X
\otimes_{f^{-1}\mathcal{O}_S}
q^{-1}\mathcal{O}_Y \otimes_{q^{-1}\mathcal{O}_Y} q^{-1}\mathcal{G} \\
& \xrightarrow{can}
p^{-1}\mathcal{F} \otimes_{q^{-1}\mathcal{O}_X}
\mathcal{O}_{X \times_S Y}
\otimes_{q^{-1}\mathcal{O}_Y} q^{-1}\mathcal{G}  \\
& =
p^{-1}\mathcal{F} \otimes_{q^{-1}\mathcal{O}_X}
\mathcal{O}_{X \times_S Y}
\otimes_{\mathcal{O}_{X \times_S Y}}
\mathcal{O}_{X \times_S Y}
\otimes_{q^{-1}\mathcal{O}_Y} q^{-1}\mathcal{G}  \\
& =
p^*\mathcal{F} \otimes_{\mathcal{O}_{X \times_S Y}} q^*\mathcal{G}
\end{align*}
qui est rarement un isomorphisme.





\section[Schéma connexe non intègre à anneaux locaux intègres]{Un schéma connexe non intègre dont les anneaux locaux sont intègres}
\label{section-connected-locally-integral-not-integral}

\noindent
Nous donnons un exemple de schéma affine $X = \Spec(A)$ qui est
connexe, dont tous les anneaux locaux sont intègres, mais qui n'est pas intègre.
La connexité de $X$ signifie que $A$ ne possède aucun idempotent non trivial, voir
Algèbre, lemme \ref{algebra-lemma-disjoint-decomposition}.
Les anneaux locaux de $X$ sont intègres si, lorsque $fg = 0$ dans $A$, tout
point de $X$ possède un voisinage sur lequel $f$ ou $g$ s'annule.
Si $A$ n'est pas un anneau intègre, alors $X$ n'est pas intègre
(Propriétés, définition \ref{properties-definition-integral}).

\medskip\noindent
En termes imagés, la construction est la suivante : soit $X_0$ la croix
(la réunion des axes de coordonnées) dans le plan affine. Soit ensuite $X_1$
l'image inverse totale (réduite) de $X_0$ dans l'éclatement du plan ($X_1$
a trois composantes rationnelles formant une chaîne).  Éclatons ensuite la
surface obtenue aux deux points singuliers de $X_1$, et soit $X_2$
l'image inverse réduite de $X_1$ (qui a cinq composantes rationnelles), etc.
Prenons pour $X$ la limite projective. Le seul problème de cette construction
est que les éclatements ajoutent une droite projective, de sorte que $X_1$ n'est pas affine.
Corrigeons cela en recollant plutôt une droite affine (notre schéma sera donc une
partie ouverte de celui décrit ci-dessus).

\medskip\noindent
Voici une construction entièrement algébrique : pour tout $k \ge 0$, soit $A_k$
l'anneau suivant : ses éléments sont les familles de
polynômes $p_i \in \mathbf{C}[x]$, où $i = 0, \ldots, 2^k$, telles que
$p_i(1) = p_{i + 1}(0)$. Posons $X_k = \Spec(A_k)$. Remarquons que $X_k$ est
la réunion de $2^k + 1$ droites affines qui se coupent transversalement en chaîne.
Définissons un homomorphisme d'anneaux $A_k \to A_{k + 1}$ par
$$
(p_0, \ldots, p_{2^k})
\longmapsto
(p_0, p_0(1), p_1, p_1(1), \ldots, p_{2^k}),
$$
autrement dit, un polynôme sur deux est constant. Ceci identifie
$A_k$ à un sous-anneau de $A_{k + 1}$. Soit $A$ la limite inductive des $A_k$
(essentiellement, leur réunion). Posons $X = \Spec(A)$. Pour tout $k$, nous avons
une injection canonique $A_k \to A$, c'est-à-dire une application $X\to X_k$.
Chaque $A_k$ est connexe mais non intègre; il s'ensuit que $A$ est
connexe mais non intègre. Il reste à montrer que les anneaux locaux de
$A$ sont intègres.

\medskip\noindent
Soient $f, g \in A$ tels que $fg = 0$ et $x \in X$. Construisons un
voisinage de $x$ sur lequel l'un de $f$ et $g$ s'annule. Choisissons $k$
tel que $f, g \in A_{k - 1}$ (noter l'indice $k - 1$).
Soit $y$ l'image de $x$ dans $X_k$. Il suffit de montrer que $y$ possède
un voisinage sur lequel $f$ ou $g$ s'annule en tant que section de
$\mathcal{O}_{X_k}$.
Si $y$ est un point lisse de $X_k$, c'est-à-dire s'il n'appartient qu'à l'une des
$2^k + 1$ droites, c'est clair. Nous pouvons donc supposer que $y$ est l'un
des $2^k$ points singuliers; deux composantes de $X_k$ passent donc par
$y$. Cependant, sur l'une de ces deux composantes (celle d'indice impair),
$f$ et $g$ sont tous deux constants, puisque ce sont des images inverses de fonctions sur
$X_{k - 1}$. Puisque $fg = 0$ partout, $f$ ou $g$ (disons $f$)
s'annule sur l'autre composante.
Il s'ensuit que $f$ s'annule sur les deux composantes, comme voulu.






\section{Un complété non complet}
\label{section-noncomplete-completion}

\noindent
Soit $R$ un anneau et soit $\mathfrak m$ un idéal maximal. Considérons le
complété
$$
R^\wedge = \lim R/\mathfrak m^n.
$$
Remarquons que $R^\wedge$ est un anneau local d'idéal maximal
$\mathfrak m' = \Ker(R^\wedge \to R/\mathfrak m)$.
En effet, si $x = (x_n) \in R^\wedge$ n'appartient pas à $\mathfrak m'$, alors
$y = (x_n^{-1}) \in R^\wedge$ vérifie $xy = 1$; ainsi, $R^\wedge$ est local d'après
Algèbre, lemme \ref{algebra-lemma-characterize-local-ring}. Or il est
toujours vrai que $R^\wedge$ est complet pour sa topologie de limite projective (voir la
discussion dans
Compléments d'algèbre, section \ref{more-algebra-section-topological-ring}).
Mais au-delà se posent les questions suivantes :
\begin{enumerate}
\item A-t-on $\mathfrak m R^\wedge = \mathfrak m'$ ?
\item $R^\wedge$, considéré comme un $R^\wedge$-module, est-il
complet pour la topologie $\mathfrak m'$-adique ?
\item $R^\wedge$, considéré comme un $R$-module, est-il complet pour la topologie $\mathfrak m$-adique ?
\end{enumerate}
Il se trouve que toutes ces questions ont une réponse négative.
L'exemple ci-dessous est tiré d'une note non publiée de
Bart de Smit et Hendrik Lenstra. Voir aussi
\cite[Exercise III.2.12]{Bourbaki-CA} et
\cite[exemple 1.8]{Yekutieli}

\medskip\noindent
Soient $k$ un corps, $R = k[x_1, x_2, x_3, \ldots]$ et
$\mathfrak m = (x_1, x_2, x_3, \ldots)$.
Nous considérerons un élément $f$ de $R^\wedge$ comme une somme éventuellement infinie
$$
f = \sum a_I x^I
$$
(en notation multi-indice) telle que, pour tout $d \geq 0$, il
n'y ait qu'un nombre fini de coefficients $a_I$ non nuls pour $|I| = d$. L'idéal
maximal $\mathfrak m' \subset R^\wedge$ est l'ensemble des $f$ dont le
terme constant est nul. En particulier, l'élément
$$
f = x_1 + x_2^2 + x_3^3 + \ldots
$$
appartient à $\mathfrak m'$ mais non à $\mathfrak m R^\wedge$, ce qui
montre que (1) est fausse dans cet exemple. Cependant, si (1) est
fausse, alors (3) l'est nécessairement car
$\mathfrak m' = \Ker(R^\wedge \to R/\mathfrak m)$
et nous pouvons appliquer
Algèbre, lemme \ref{algebra-lemma-hathat} avec $n = 1$.

\medskip\noindent
Pour terminer, démontrons que $R^\wedge$ n'est pas complet pour la topologie $\mathfrak m'$-adique.
Pour $n \geq 1$, posons $K_n = \Ker(R^\wedge \to R/\mathfrak m^n)$. Alors
nous avons des suites exactes courtes
$$
0 \to K_n/(\mathfrak m')^n \to R^\wedge/(\mathfrak m')^n \to
R/\mathfrak m^n \to 0
$$
L'application de projection $R^\wedge \to R/\mathfrak m^{n + 1}$ envoie
$(\mathfrak m')^n$ sur $\mathfrak m^n/\mathfrak m^{n + 1}$.
Il s'ensuit que $K_{n + 1} \to K_n/(\mathfrak m')^n$ est surjective.
Par conséquent, le système projectif $\left(K_n/(\mathfrak m')^n\right)$
a des applications de transition surjectives et,
en passant aux limites projectives, nous obtenons une suite exacte
$$
0 \to \lim K_n/(\mathfrak m')^n \to
\lim R^\wedge/(\mathfrak m')^n \to
\lim R/\mathfrak m^n \to 0
$$
d'après Algèbre, lemme \ref{algebra-lemma-Mittag-Leffler}.
Nous voyons ainsi que $R^\wedge$ est complet pour la topologie $\mathfrak m'$-adique
si et seulement si $K_n = (\mathfrak m')^n$ pour tout $n \geq 1$.

\medskip\noindent
Pour montrer que $R^\wedge$ n'est pas complet pour la topologie $\mathfrak m'$-adique
dans notre exemple, montrons que $K_2 = \Ker(R^\wedge \to R/\mathfrak m^2)$
n'est pas égal à $(\mathfrak m')^2$.
Remarquons qu'un élément de $(\mathfrak m')^2$
peut s'écrire comme une somme finie
\begin{equation}
\label{equation-sum}
\sum\nolimits_{i = 1, \ldots, t} f_i g_i
\end{equation}
où $f_i, g_i \in R^\wedge$ ont un terme constant nul.
Pour obtenir un exemple, nous allons choisir un $z \in K_2$
de la forme
$$
z = z_1 + z_2 + z_3 + \ldots
$$
ayant les propriétés suivantes
\begin{enumerate}
\item il existe des suites $1 < d_1 < d_2 < d_3 < \ldots $ et
$0 < n_1 < n_2 < n_3 < \ldots$ telles que
$z_i \in k[x_{n_i}, x_{n_i + 1}, \ldots, x_{n_{i + 1} - 1}]$
soit homogène de degré $d_i$, et
\item dans l'anneau $k[[x_{n_i}, x_{n_i + 1}, \ldots, x_{n_{i + 1} - 1}]]$,
l'élément $z_i$ ne peut s'écrire comme une somme (\ref{equation-sum})
avec $t \leq i$.
\end{enumerate}
Cela implique manifestement que $z$ n'appartient pas à $(\mathfrak m')^2$,
car l'image de la relation (\ref{equation-sum}) dans
l'anneau $k[[x_{n_i}, x_{n_i + 1}, \ldots, x_{n_{i + 1} - 1}]]$
pour $i$ assez grand conduirait à une contradiction. Il suffit donc
de montrer que, pour tout $t > 0$, il existe un $d \gg 0$ et un entier
$n$ tels que nous puissions trouver un élément homogène
$z \in k[x_1, \ldots, x_n]$ de degré $d$ qui ne puisse s'écrire comme
une somme (\ref{equation-sum}) pour ce $t$ dans $k[[x_1, \ldots, x_n]]$.
Prenons $n > 2t$ et un $d > 1$ quelconque, premier à la caractéristique de $k$, et
posons $z = \sum_{i = 1, \ldots, n} x_i^d$. Alors l'ensemble des zéros
de l'idéal
$$
(\frac{\partial z}{\partial x_1}, \ldots, \frac{\partial z}{\partial x_n})
=
(dx_1^{d - 1}, \ldots, dx_n^{d - 1})
$$
se réduit à un point. D'autre part,
$$
\frac{\partial ( \sum\nolimits_{i = 1, \ldots, t} f_i g_i ) }{\partial x_j}
\in (f_1, \ldots, f_t, g_1, \ldots, g_t)
$$
par la règle de Leibniz; par conséquent, l'ensemble des zéros de ces dérivées
contient au moins
$$
V(f_1, \ldots, f_t, g_1, \ldots, g_t) \subset
\Spec(k[[x_1, \ldots, x_n]]).
$$
On obtient ainsi une contradiction, puisque la dimension de
$V(f_1, \ldots, f_t, g_1, \ldots, g_t)$ est au moins $n - 2t \geq 1$.

\begin{lemma}
\label{lemma-noncomplete-completion}
Il existe un anneau local $R$ et un idéal maximal $\mathfrak m$ tels que
le complété $R^\wedge$ de $R$ pour la topologie $\mathfrak m$-adique ait les
propriétés suivantes :
\begin{enumerate}
\item $R^\wedge$ est local, mais son idéal maximal n'est pas égal à
$\mathfrak m R^\wedge$,
\item $R^\wedge$ n'est pas un anneau local complet, et
\item $R^\wedge$ n'est pas complet pour la topologie $\mathfrak m$-adique comme $R$-module.
\end{enumerate}
\end{lemma}

\begin{proof}
Cela résulte de la discussion qui précède, puisque (avec $R = k[x_1, x_2, x_3, \ldots]$)
le complété du localisé $R_{\mathfrak m}$ est égal au
complété de $R$.
\end{proof}





\section{Un quotient non complet}
\label{section-noncomplete-quotient}

\noindent
Soit $k$ un corps. Posons
$$
R = k[t, z_1, z_2, z_3, \ldots, w_1, w_2, w_3, \ldots, x]/
(z_it - x^iw_i, z_i w_j)
$$
Remarquons qu'en particulier $z_iz_jt = 0$ dans cet anneau. Tout élément $f$ de $R$
s'écrit de manière unique comme une somme finie
$$
f = \sum\nolimits_{i = 0, \ldots, d} f_i x^i
$$
où chaque $f_i \in k[t, z_i, w_j]$ ne comporte aucun terme contenant les produits
$z_it$ ou $z_iw_j$. De plus, si $f$ s'écrit de cette manière, alors
$f \in (x^n)$ si et seulement si $f_i = 0$ pour $i < n$.
Ainsi, $x$ est un non-diviseur de zéro et $\bigcap (x^n) = 0$.
Soit $R^\wedge$ le complété de $R$ pour la topologie définie par l'idéal $(x)$.
Remarquons que $R^\wedge$ est complet pour la topologie $(x)$-adique; voir
Algèbre, lemme \ref{algebra-lemma-hathat-finitely-generated}.
D'après ce qui précède, un élément de $R^\wedge$ s'écrit de manière unique
comme une somme infinie
$$
f = \sum\nolimits_{i = 0}^\infty f_i x^i
$$
où chaque $f_i \in k[t, z_i, w_j]$ ne comporte aucun terme contenant les produits
$z_it$ ou $z_iw_j$. Considérons l'élément
$$
f = \sum\nolimits_{i = 1}^\infty x^i w_i =
xw_1 + x^2w_2 + x^3w_3 + \ldots
$$
c'est-à-dire que $f_n = w_n$. Remarquons que $f \in (t , x^n)$ pour tout $n$
car $x^mw_m \in (t)$ pour tout $m$.
Nous affirmons que $f \not \in (t)$. Pour le montrer, supposons que
$tg = f$, où $g = \sum g_lx^l$ est sous la forme canonique ci-dessus.
Comme $tz_iz_j = 0$, nous pouvons supposer que les $g_l$ ne
comportent aucun terme contenant les produits $z_iz_j$. En examinant comment
$tg$ est mis sous forme canonique, nous obtenons les règles suivantes :
Considérons un terme $c m$ quelconque de $g_l$, où $c \in k$ et où $m$ est un
monôme en $t, z_i, w_j$; effectuons la distinction suivante :
\begin{enumerate}
\item si le monôme $m$ ne fait intervenir aucun $z_i$, alors $ctm$ est
un terme de $f_l$, et
\item si le monôme $m$ fait intervenir un $z_i$, alors il est égal à
$m = z_i$ et nous voyons que $cw_i$ est un terme de $f_{l + i}$.
\end{enumerate}
Comme $g_0$ est un polynôme, les variables $z_i$ qui y apparaissent sont en nombre fini.
Choisissons $n$ tel que $z_n$ n'apparaisse pas dans $g_0$.
Alors les règles ci-dessus montrent que $w_n$ n'apparaît pas dans $f_n$, ce qui est
une contradiction. Il s'ensuit que $R^\wedge/(t)$ n'est pas complet; voir
Algèbre, lemme \ref{algebra-lemma-quotient-complete}.

\begin{lemma}
\label{lemma-noncomplete-quotient}
Il existe un anneau $R$ complet pour la topologie définie par un idéal principal
$I$ et un idéal principal $J$ tels que $R/J$ ne soit pas complet pour la topologie
$I$-adique.
\end{lemma}

\begin{proof}
Voir la discussion qui précède.
\end{proof}



\section{La complétion n'est pas exacte}
\label{section-completion-not-exact}

\noindent
Voici un exemple simple. Supposons que $R = k[t]$. Notons $M^\wedge$
le complété d'un $R$-module pour la topologie définie par $I = (t)$. Posons
$P = K = \bigoplus_{n \in \mathbf{N}} R$ et
$M = \bigoplus_{n \in \mathbf{N}} R/(t^n)$. Il existe alors une suite exacte courte
$0 \to K \to P \to M \to 0$ dont la première application est donnée par la
multiplication par $t^n$ sur le $n$-ième facteur direct. Nous affirmons que
$0 \to K^\wedge \to P^\wedge \to M^\wedge \to 0$ n'est pas exacte au terme médian.
En effet, $\xi = (t^2, t^3, t^4, \ldots) \in P^\wedge$ a une image nulle dans
$M^\wedge$, mais n'appartient pas à l'image de $K^\wedge \to P^\wedge$, car
il devrait être l'image de $(t, t, t, \ldots)$, qui n'est pas un élément de
$K^\wedge$.

\medskip\noindent
Voici un exemple « plus petit ». Dans la situation du
lemme \ref{lemma-noncomplete-quotient},
la suite exacte courte $0 \to J \to R \to R/J \to 0$ ne reste pas
exacte après complétion pour la topologie $I$-adique.
En effet, si $f \in J$ est un générateur, alors
$f : R \to J$ est surjective, donc $R \to J^\wedge$ est surjective, et donc
l'image de $J^\wedge \to R$ est $(f) = J$; mais le fait que
$R/J$ ne soit pas complet signifie que le noyau de la surjection
$R \to (R/J)^\wedge$ contient strictement $J$; voir
Algèbre, lemmes \ref{algebra-lemma-completion-generalities} et
\ref{algebra-lemma-quotient-complete}.
De même, la suite
$R \to R \to R/(f) \to 0$ ne reste pas exacte après complétion.

\begin{lemma}
\label{lemma-completion-not-exact}
\begin{slogan}
En général, la complétion n'est exacte ni à gauche ni à droite.
\end{slogan}
En général, la complétion n'est pas un foncteur exact; elle n'est même pas
exacte à droite. Cela vaut encore lorsque $I$ est de type fini,
sur la catégorie des modules de présentation finie.
\end{lemma}

\begin{proof}
Voir la discussion qui précède.
\end{proof}



\section{La catégorie des modules complets n'est pas abélienne}
\label{section-non-abelian}

\noindent
Soient $R$ un anneau et $I \subset R$ un idéal de type fini.
Considérons la catégorie $\mathcal{A}$ des modules complets pour la topologie $I$-adique
qui sont des $R$-modules; voir
Algèbre, définition \ref{algebra-definition-complete}.
Soit $\varphi : M \to N$ un morphisme de $\mathcal{A}$.
Le conoyau de $\varphi$ dans $\mathcal{A}$ est le complété $I$-adique
$(\Coker(\varphi))^\wedge$ du conoyau usuel
(comme $I$ est de type fini, ce complété est complet; voir
Algèbre, lemme \ref{algebra-lemma-hathat-finitely-generated}).
Posons $K = \Ker(\varphi)$; comme $\varphi$ est continue, c'est
un sous-module fermé de $M$. Le lemme \ref{lemma-closed-is-complete}
ci-dessous montre que $K$ est complet pour la topologie $I$-adique et est donc le
noyau de $\varphi$ dans $\mathcal{A}$.

\begin{lemma}
\label{lemma-closed-is-complete}
Soit $I \subset R$ un idéal d'un anneau. Soit $M$ un module complet
pour la topologie $I$-adique sur $R$ et soit $N \subset M$ un sous-module. Les conditions suivantes sont équivalentes :
\begin{enumerate}
\item $N$ est fermé dans $M$,
\item $N = \bigcap_{n \geq 1} (N + I^n M)$,
\item $M/N$ est complet pour la topologie $I$-adique.
\end{enumerate}
Si $I$ est de type fini, ces conditions impliquent que $N$
est complet pour la topologie $I$-adique.
\end{lemma}

\begin{proof}
L'équivalence de (1) et (2) résulte de
Espaces formels, lemme \ref{formal-spaces-lemma-closed}.
L'équivalence de (2) et (3) est le
lemme \ref{algebra-lemma-quotient-complete} d'Algèbre.
Supposons $I$ de type fini et (1) satisfaite.
Notons $N^\wedge$ le complété $I$-adique de $N$. Comme $M$
est complet pour la topologie $I$-adique, l'application d'inclusion se factorise en
$$
N \to N^\wedge \to M
$$
Comme $N^\wedge \to M$ est continue, que $N$ est dense dans $N^\wedge$
et que $N$ est fermé dans $M$, nous voyons que
$N^\wedge \to M$ se factorise par $N$. Ainsi $N^\wedge = N \oplus C$
pour un certain $R$-module $C$. Ainsi $N$ est un facteur direct du
module complet pour la topologie $I$-adique qu'est $N^\wedge$; voir
Algèbre, lemme \ref{algebra-lemma-hathat-finitely-generated}
(c'est ici que nous utilisons le fait que $I$ est de type fini),
et donc $N$ est complet pour la topologie $I$-adique.
\end{proof}

\noindent
Nous allons donner un exemple montrant qu'en général $\Im \not = \Coim$ dans la
catégorie $\mathcal{A}$. Prenons
$R = \mathbf{Z}_p = \lim_n \mathbf{Z}/p^n\mathbf{Z}$
pour l'anneau des entiers $p$-adiques et prenons $I = (p)$.
Considérons l'application
$$
\text{diag}(1, p, p^2, \ldots) :
\left(\bigoplus\nolimits_{n \geq 1} \mathbf{Z}_p\right)^\wedge
\longrightarrow
\prod\nolimits_{n \geq 1} \mathbf{Z}_p
$$
où le membre de gauche est le complété $p$-adique de la somme directe.
Un élément du membre de gauche est donc un vecteur $(x_1, x_2, x_3, \ldots)$
tel que $x_i \in \mathbf{Z}_p$ et que sa valuation $p$-adique vérifie $v_p(x_i) \to \infty$ lorsque
$i \to \infty$. Son image est $(x_1, px_2, p^2x_3, \ldots)$. Nous voyons donc
que $(1, p, p^2, \ldots)$ appartient à l'adhérence de l'image, mais pas à
l'image. D'après la description des noyaux et conoyaux donnée ci-dessus, il est
clair que $\Im \not = \Coim$ pour cette application.

\begin{lemma}
\label{lemma-complete-modules-not-abelian}
Soient $R$ un anneau et $I \subset R$ un idéal de type fini.
La catégorie des modules complets pour la topologie $I$-adique qui sont des $R$-modules possède des noyaux et des
conoyaux, mais elle n'est pas abélienne en général, même lorsque $R$ est noethérien.
\end{lemma}

\begin{proof}
Voir ci-dessus.
\end{proof}





\section{La catégorie des modules complets au sens dérivé}
\label{section-derived-complete-modules}

\noindent
Lire d'abord Compléments d'algèbre, section
\ref{more-algebra-section-derived-complete-modules}
avant d'aborder la présente section.

\medskip\noindent
Soient $A$ un anneau et $I$ un idéal de $A$, et notons $\mathcal{C}$
la catégorie des modules complets au sens dérivé définie dans
Compléments d'algèbre, définition \ref{more-algebra-definition-derived-complete}.

\medskip\noindent
Soient $T$ un ensemble et $M_t$, $t \in T$, une famille de modules complets
au sens dérivé. Nous affirmons qu'en général $\bigoplus M_t$ n'est pas un module complet
au sens dérivé. Pour un exemple précis, prenons $A = \mathbf{Z}_p$ et $I = (p)$ et
considérons $\bigoplus_{n \in \mathbf{N}} \mathbf{Z}_p$. L'application de
$\bigoplus_{n \in \mathbf{N}} \mathbf{Z}_p$ dans son complété $p$-adique
n'est pas surjective. Cela signifie que $\bigoplus_{n \in \mathbf{N}} \mathbf{Z}_p$
ne peut pas être complet au sens dérivé, car le contraire entraînerait sa surjectivité; voir
Compléments d'algèbre, lemme \ref{more-algebra-lemma-complete-derived-complete}.
Par conséquent, le foncteur d'inclusion $\mathcal{C} \to \text{Mod}_A$ ne commute
ni aux sommes directes ni aux limites inductives (filtrantes).

\medskip\noindent
Supposons $I$ de type fini. D'après la discussion de Compléments d'algèbre, section
\ref{more-algebra-section-derived-complete-modules}, la catégorie
$\mathcal{C}$ possède des limites inductives quelconques. Cependant, nous affirmons que les limites inductives filtrantes
ne sont pas exactes dans la catégorie $\mathcal{C}$. En effet, supposons que
$A = \mathbf{Z}_p$ et $I = (p)$. On a des inclusions
$f_n : \mathbf{Z}_p/p\mathbf{Z}_p \to \mathbf{Z}_p/p^n\mathbf{Z}_p$
de modules complets pour la topologie $p$-adique sur $A$, données par la multiplication par
$p^{n - 1}$. On a des diagrammes commutatifs
$$
\xymatrix{
\mathbf{Z}_p/p\mathbf{Z}_p \ar[r]_{f_n} \ar[d]^1 &
\mathbf{Z}_p/p^n\mathbf{Z}_p \ar[d]_p \\
\mathbf{Z}_p/p\mathbf{Z}_p \ar[r]^{f_{n + 1}} &
\mathbf{Z}_p/p^{n + 1}\mathbf{Z}_p
}
$$
Nous affirmons que la limite inductive de ces inclusions dans la catégorie $\mathcal{C}$
donne l'application $\mathbf{Z}_p/p\mathbf{Z}_p \to 0$. En effet, la limite inductive dans
$\text{Mod}_A$ du système de droite est $\mathbf{Q}_p/\mathbf{Z}_p$.
La limite inductive dans $\mathcal{C}$ est donc
$$
H^0((\mathbf{Q}_p/\mathbf{Z}_p)^\wedge) =
H^0(\mathbf{Z}_p[1]) = 0
$$
d'après Compléments d'algèbre, section
\ref{more-algebra-section-derived-complete-modules}
où ${}^\wedge$ désigne la complétion dérivée.
Cela démontre notre affirmation.

\begin{lemma}
\label{lemma-derived-complete-modules}
Soient $A$ un anneau et $I \subset A$ un idéal.
La catégorie $\mathcal{C}$ des modules complets au sens dérivé
est abélienne, et le foncteur d'inclusion $F : \mathcal{C} \to \text{Mod}_A$
est exact et commute aux limites projectives quelconques.
Si $I$ est de type fini, alors $\mathcal{C}$ possède
des sommes directes quelconques et des limites inductives quelconques, mais $F$ ne commute pas à celles-ci
en général. Enfin, les limites inductives filtrantes ne sont pas exactes dans $\mathcal{C}$
en général; ainsi, $\mathcal{C}$ n'est pas une catégorie abélienne de Grothendieck.
\end{lemma}

\begin{proof}
Voir Compléments d'algèbre, lemme \ref{more-algebra-lemma-derived-complete-modules},
et la discussion qui précède.
\end{proof}





\section{Des complétés non plats}
\label{section-nonflat}

\noindent
Le complété d'un anneau pour un idéal n'est pas toujours plat,
contrairement au cas noethérien. Nous avons vu deux exemples de ce
phénomène dans
Compléments d'algèbre, exemple \ref{more-algebra-example-not-glueing-pair}.
Nous donnons dans cette section deux exemples supplémentaires.

\begin{lemma}
\label{lemma-countable-fg-tensor}
Soit $R$ un anneau. Soit $M$ un $R$-module dénombrable.
Alors $M$ est un $R$-module de type fini si et seulement si
$M \otimes_R R^\mathbf{N} \to M^\mathbf{N}$ est surjective.
\end{lemma}

\begin{proof}
Si $M$ est un module de type fini, l'application est surjective d'après Algèbre, proposition
\ref{algebra-proposition-fg-tensor}. Réciproquement, supposons l'application surjective.
Soit $m_1, m_2, m_3, \ldots$ une énumération des éléments de $M$.
Soit $\sum_{j = 1, \ldots, m} x_j \otimes a_j$ un élément du
produit tensoriel dont l'image est $(m_n) \in M^\mathbf{N}$. Alors
nous voyons que $x_1, \ldots, x_m$ engendrent $M$ sur $R$, comme dans la démonstration de
Algèbre, proposition \ref{algebra-proposition-fg-tensor}.
\end{proof}

\begin{lemma}
\label{lemma-countable-fp-tensor}
Soit $R$ un anneau dénombrable. Soit $M$ un $R$-module dénombrable. Alors $M$
est de présentation finie si et seulement si l'homomorphisme canonique
$M \otimes_R R^\mathbf{N} \to M^\mathbf{N}$ est un isomorphisme.
\end{lemma}

\begin{proof}
Si $M$ est un module de présentation finie, l'application est un isomorphisme
d'après Algèbre, proposition \ref{algebra-proposition-fp-tensor}. Réciproquement,
supposons que l'application soit un isomorphisme. Le lemme \ref{lemma-countable-fg-tensor}
montre que le module $M$ est de type fini. Choisissons une surjection $R^{\oplus m} \to M$ de
noyau $K$. Alors $K$ est dénombrable comme sous-module de $R^{\oplus m}$.
En raisonnant comme dans la démonstration d'Algèbre, proposition
\ref{algebra-proposition-fp-tensor}, nous voyons que
$K \otimes_R R^\mathbf{N} \to K^\mathbf{N}$ est surjective.
Nous concluons donc que $K$ est un $R$-module de type fini par le
lemme \ref{lemma-countable-fg-tensor}.
Ainsi $M$ est de présentation finie.
\end{proof}

\begin{lemma}
\label{lemma-countable-coherent}
Soit $R$ un anneau dénombrable. Alors $R$ est cohérent si et seulement si
$R^\mathbf{N}$ est un $R$-module plat.
\end{lemma}

\begin{proof}
Si $R$ est cohérent, alors $R^\mathbf{N}$ est un module plat d'après
Algèbre, proposition \ref{algebra-proposition-characterize-coherent}.
Supposons $R^\mathbf{N}$ plat. Soit $I \subset R$ un idéal de type
fini. Pour démontrer le lemme, nous montrons que $I$ est de
présentation finie comme $R$-module. En effet, l'application
$I \otimes_R R^\mathbf{N} \to R^\mathbf{N}$ est
injective puisque $R^\mathbf{N}$ est plat et que son image est
$I^\mathbf{N}$ d'après le lemme \ref{lemma-countable-fg-tensor}.
Nous concluons donc par le lemme \ref{lemma-countable-fp-tensor}.
\end{proof}

\noindent
Soit $R$ un anneau dénombrable. Remarquons que $R[[x]]$ est isomorphe à
$R^\mathbf{N}$ comme $R$-module. Le lemme \ref{lemma-countable-coherent}
montre que $R \to R[[x]]$ est plat si et seulement si $R$ est cohérent.
Il existe de nombreux anneaux dénombrables non cohérents, par exemple
$$
R = k[y, z, a_1, b_1, a_2, b_2, a_3, b_3, \ldots]/
(a_1 y + b_1 z, a_2 y + b_2 z, a_3 y + b_3 z, \ldots)
$$
où $k$ est un corps dénombrable. Cet anneau n'est pas cohérent, car
l'idéal $(y, z)$ de $R$ n'est pas un $R$-module de présentation finie.
Notons que $R[[x]]$ est le complété de $R[x]$ pour l'idéal
principal $(x)$.

\begin{lemma}
\label{lemma-completion-polynomial-ring-not-flat}
Il existe un anneau tel que le complété $R[[x]]$ de $R[x]$
en $(x)$ ne soit pas plat sur $R$ et, a fortiori, pas plat sur $R[x]$.
\end{lemma}

\begin{proof}
Voir la discussion qui précède.
\end{proof}

\noindent
Il existe en fait un anneau $R$ tel que $R[[x]]$ soit plat
sur $R$, mais que $R[[x]]$ ne soit pas plat sur $R[x]$. Voir
\href{https://math.stackexchange.com/users/164860/badam-baplan}{ce billet}
de Badam Baplan. En effet, soit
$R$ un anneau de valuation. Alors $R$ est cohérent (Algèbre, exemple
\ref{algebra-example-valuation-ring-coherent}) et donc
$R[[x]]$ est plat sur $R$ d'après
Algèbre, proposition \ref{algebra-proposition-characterize-coherent}.
En revanche, nous avons le lemme suivant.

\begin{lemma}
\label{lemma-almost-integral-when-powerseries-flat}
Soit $R$ un anneau intègre de corps des fractions $K$.
Si $R[[x]]$ est plat sur $R[x]$, alors $R$ est normal si et seulement
si $R$ est complètement intégralement clos
(Algèbre, définition \ref{algebra-definition-almost-integral}).
\end{lemma}

\begin{proof}
Supposons donnés $\alpha \in K$ et un élément non nul $r \in R$ tels que
$r \alpha^n \in R$ pour tout $n \geq 1$. Considérons alors
$f = \sum r \alpha^n x^n$ dans $R[[x]]$. Écrivons $\alpha = a/b$
avec $a, b \in R$ et $b$ non nul. Nous voyons alors que $(a x - b)f = -rb$.
Il s'ensuit que $rb$ appartient à l'idéal $(ax - b)R[[x]]$.
Soit $S = \{h \in R[x] : h(0) = 1\}$. C'est une partie multiplicativement stable,
et la platitude de $R[x] \to R[[x]]$ implique que $S^{-1}R[x] \to R[[x]]$
est fidèlement plat (détails omis; indication : utiliser Algèbre, lemme
\ref{algebra-lemma-ff-rings}). Par conséquent,
$$
S^{-1}R[x]/(ax - b)S^{-1}R[x] \to R[[x]]/(ax - b)R[[x]]
$$
est injective. Nous en concluons que
$h rb = (ax - b) g$ pour certains $h \in S$ et $g \in R[x]$.
En écrivant $h = 1 + h_1 x + \ldots + h_d x^d$, nous obtenons
$$
1 + h_1 x + \ldots + h_d x^d = (1/r)(\alpha x - 1)g
$$
Cette factorisation dans $K[x]$ donne une factorisation correspondante
dans $K[x^{-1}]$, qui montre que $\alpha$ est racine d'un polynôme unitaire
à coefficients dans $R$, comme voulu.
\end{proof}

\begin{lemma}
\label{lemma-completion-polynomial-ring-not-flat-bis}
Si $R$ est un anneau de valuation de dimension $> 1$, alors $R[[x]]$
est plat sur $R$, mais pas sur $R[x]$.
\end{lemma}

\begin{proof}
Les arguments ci-dessus montrent que, pour cela, il suffit d'établir que
$R$ n'est pas complètement intégralement clos (les anneaux de valuation sont
normaux; voir Algèbre, lemme \ref{algebra-lemma-valuation-ring-normal}).
Soit $\mathfrak p \subset \mathfrak m \subset R$ une chaîne d'idéaux premiers.
Choisissons $x \in \mathfrak p$ non nul et $y \in \mathfrak m \setminus \mathfrak p$.
Alors $x y^{-n} \in R$ pour tout $n \geq 1$ (sinon $y^n/x \in R$,
ce qui est absurde puisque $y \not \in \mathfrak p$). Ainsi $1/y$ est
presque entier sur $R$, mais n'appartient pas à $R$.
\end{proof}

\noindent
Construisons maintenant un exemple où le complété d'un localisé
n'est pas plat. Pour cela, considérons l'anneau
$$
R = k[y, z, a_1, a_2, a_3, \ldots]/(ya_i, a_i a_j)
$$
Notons $f \in R$ la classe de $z$. Nous affirmons que l'homomorphisme d'anneaux
\begin{equation}
\label{equation-nonflat}
R[[x]] \longrightarrow R_f[[x]]
\end{equation}
n'est pas plat. Soit $I$ le noyau de $y : R[[x]] \to R[[x]]$. Un
élément $g$ de $I$ est typiquement de la forme $g = \sum g_{n, m} a_mx^n$
où $g_{n, m} \in k[z]$ et, pour $n$ donné, seul un nombre fini de
coefficients $g_{n, m}$ sont non nuls. Soit $J$ le noyau de $y : R_f[[x]] \to R_f[[x]]$.
Nous affirmons que $J \not = I R_f[[x]]$. En effet, si c'était le cas, nous
aurions
$$
\sum z^{-n} a_n x^n = \sum\nolimits_{i = 1, \ldots, m} h_i g_i
$$
pour un certain $m \geq 1$, avec $g_i \in I$ et $h_i \in R_f[[x]]$. Écrivons
$h_i = \bar h_i \bmod (y, a_1, a_2, a_3, \ldots)$
avec $\bar h_i \in k[z, 1/z][[x]]$. En considérant le coefficient de
$a_n$ et en utilisant la description ci-dessus des éléments $g_i$, nous obtiendrions
$$
z^{-n} x^n = \sum \bar h_i \bar g_{i, n}
$$
pour certains $\bar g_{i, n} \in k[z][[x]]$. Cela signifierait que
tous les $z^{-n}x^n$ appartiennent au $k[z][[x]]$-module de type fini
engendré par les éléments $\bar h_i$. Puisque $k[z][[x]]$ est noethérien,
il en résulterait que le sous-$k[z][[x]]$-module de $k[z, 1/z][[x]]$
engendré par $1, z^{-1}x, z^{-2}x^2, \ldots$ est de type fini. D'après
Algèbre, lemme \ref{algebra-lemma-characterize-integral-element},
nous conclurions que $z^{-1}x$ est entier sur $k[z][[x]]$,
ce qui est absurde. En revanche,
si (\ref{equation-nonflat}) était plat, alors nous
aurions $J = IR_f[[x]]$ en tensorisant la suite exacte
$0 \to I \to R[[x]] \xrightarrow{y} R[[x]]$ par $R_f[[x]]$.

\begin{lemma}
\label{lemma-nonflat-completion-localization}
Il existe un anneau $A$ complet pour la topologie définie par un idéal principal $I$
et un élément $f \in A$ tels que le complété $I$-adique
$A_f^\wedge$ de $A_f$ ne soit pas plat sur $A$.
\end{lemma}

\begin{proof}
Posons $A = R[[x]]$ et $I = (x)$, et remarquons que $R_f[[x]]$
est le complété de $R[[x]]_f$.
\end{proof}





\section{Catégorie non abélienne des modules quasi-cohérents}
\label{section-nonabelian-QCoh}

\noindent
Dans Faisceaux sur les champs, section \ref{stacks-sheaves-section-quasi-coherent},
nous avons défini la catégorie des modules quasi-cohérents sur une catégorie fibrée en
groupoïdes au-dessus de $\Sch$. Bien que nous montrions dans
Faisceaux sur les champs, section
\ref{stacks-sheaves-section-quasi-coherent-algebraic-stacks}
que cette catégorie est abélienne pour les champs algébriques, nous montrons dans cette
section que ce n'est pas le cas pour les espaces algébriques formels.

\medskip\noindent
Plus précisément, considérons $\mathbf{Z}_p$ comme un anneau topologique, muni de
la topologie $p$-adique. Posons $X = \text{Spf}(\mathbf{Z}_p)$ ; voir
Espaces formels, définition
\ref{formal-spaces-definition-affine-formal-spectrum}.
Alors $X$ est un faisceau d'ensembles sur $(\Sch/\mathbf{Z})_{fppf}$
et donne lieu à un champ en catégories discrètes $\mathcal{X}$ ; voir
Champs, lemme \ref{stacks-lemma-when-stack-in-sets}.
La discussion de Faisceaux sur les champs, section
\ref{stacks-sheaves-section-quasi-coherent-algebraic-stacks}
s'applique donc.

\medskip\noindent
Soit $\mathcal{F}$ un module quasi-cohérent sur $\mathcal{X}$.
Puisque $X = \colim \Spec(\mathbf{Z}/p^n\mathbf{Z})$, il résulte aussitôt de
Faisceaux sur les champs, lemme \ref{stacks-sheaves-lemma-quasi-coherent},
que $\mathcal{F}$ est déterminé par une suite $(\mathcal{F}_n)$ vérifiant :
\begin{enumerate}
\item $\mathcal{F}_n$ est un module quasi-cohérent sur
$\Spec(\mathbf{Z}/p^n\mathbf{Z})$, et
\item les applications de transition induisent des isomorphismes
$\mathcal{F}_n = \mathcal{F}_{n + 1}/p^n\mathcal{F}_{n + 1}$.
\end{enumerate}
En passant aux modules, nous voyons que $\mathcal{F}$ correspond à un
système $(M_n)$ où chaque $M_n$ est un groupe abélien annulé
par $p^n$ et où les applications de transition induisent des isomorphismes
$M_n = M_{n + 1}/p^n M_{n + 1}$. Dans cette situation, le module
$M = \lim M_n$ est complet pour la topologie $p$-adique et $M_n = M/p^n M$ ; voir
Algèbre, lemme \ref{algebra-lemma-limit-complete}.
Nous en déduisons que la catégorie des modules quasi-cohérents sur $X$
est équivalente à la catégorie des groupes abéliens complets pour
la topologie $p$-adique. Cette catégorie n'est pas abélienne ; voir
la section \ref{section-non-abelian}.

\begin{lemma}
\label{lemma-quasi-coherent-not-abelian}
La catégorie des modules quasi-cohérents\footnote{Les modules quasi-cohérents
sont ici entendus au sens défini ci-dessus. Compte tenu de la manière dont le projet Champs
est organisé, il s'agit bien de la bonne définition ; comme nous l'avons vu ci-dessus,
notre définition coïncide avec celle que l'on aurait naïvement donnée des modules quasi-cohérents
sur $\text{Spf}(A)$, à savoir les $A$-modules complets.}
sur un espace algébrique formel
$X$ n'est pas abélienne en général, même lorsque $X$ est un espace algébrique
formel affine noethérien.
\end{lemma}

\begin{proof}
Voir la discussion qui précède.
\end{proof}





\section{Suite non scindée, mais localement scindée}
\label{section-nonsplit-locally-split}

\noindent
Considérons la suite exacte courte
$$
0 \to M \to
\bigoplus\nolimits_{p\text{ premier}} \mathbf{Z}_{(p)} \to \mathbf{Q} \to 0
$$
où $M$ est le noyau de l'homomorphisme de droite, défini par l'inclusion
$\mathbf{Z}_{(p)} \to \mathbf{Q}$ sur chaque facteur. Cette suite de
$\mathbf{Z}$-modules n'est pas scindée, car il n'existe aucun
homomorphisme non nul $\mathbf{Q} \to \mathbf{Z}_{(p)}$. En revanche, si
nous localisons en un idéal premier quelconque $p$, la suite devient scindée.

\begin{lemma}
\label{lemma-nonsplit-locally-split}
Il existe un anneau $R$ et une suite non scindée de modules
qui devient scindée localement pour la topologie de Zariski.
\end{lemma}

\begin{proof}
Voir la discussion qui précède.
\end{proof}





\section{Suites régulières et changement de base}
\label{section-regular-base-change}

\noindent
Nous allons construire un anneau $R$ muni d'une suite régulière
$(x, y, z)$ telle qu'il existe un élément non nul $\delta \in R/zR$
vérifiant $x\delta = y\delta = 0$.

\medskip\noindent
Pour construire notre exemple, nous commençons par
construire un module particulier $E$ sur l'anneau $k[x, y, z]$,
où $k$ est un corps quelconque. Plus précisément, $E$ sera une somme amalgamée, comme
dans le diagramme suivant :
$$
\xymatrix{
\frac{xk[x, y, z, y^{-1}]}{xyk[x, y, z]} \ar[r] \ar[d]^{z/x} &
\frac{k[x, y, z, x^{-1}, y^{-1}]}{yk[x, y, z, x^{-1}]} \ar[r] \ar[d] &
\frac{k[x, y, z, x^{-1}, y^{-1}]}{yk[x, y, z, x^{-1}] + xk[x, y, z, y^{-1}]}
\ar[d] \\
\frac{k[x, y, z, y^{-1}]}{yzk[x, y, z]} \ar[r] &
E \ar[r] &
\frac{k[x, y, z, x^{-1}, y^{-1}]}{yk[x, y, z, x^{-1}] + xk[x, y, z, y^{-1}]}
}
$$
dont les lignes sont des suites exactes courtes (nous avons omis les zéros extérieurs pour
des raisons typographiques). On peut aussi décrire $E$ par
$$
E = \{(f, g) \mid f \in k[x, y, z, x^{-1}, y^{-1}],
g \in k[x, y, z, y^{-1}] \}/\sim
$$
où $(f, g) \sim (f', g')$ si et seulement s'il existe un
$h \in k[x, y, z, y^{-1}]$ tel que
$$
f = f' + xh \bmod yk[x, y, z, x^{-1}], \quad
g = g' - zh \bmod yzk[x, y, z]
$$
Nous affirmons que : (a) $x : E \to E$ est injective, (b)
$y : E/xE \to E/xE$ est injective, (c) $E/(x, y)E = 0$, (d) il
existe un élément non nul $\delta \in E/zE$ tel que
$x\delta = y\delta = 0$.

\medskip\noindent
Pour prouver (a), supposons que $(f, g)$ soit un couple qui représente un
élément de $E$ et que $(xf, xg) \sim 0$. Il existe alors un
$h \in k[x, y, z, y^{-1}]$ tel que $xf + xh \in yk[x, y, z, x^{-1}]$
et $xg - zh \in yzk[x, y, z]$. Nous pouvons supposer que
$h = \sum a_{i, j, k}x^iy^jz^k$ est une somme de monômes dans laquelle seuls
les $j \leq 0$ interviennent. Alors $xg - zh \in yzk[x, y, z]$ entraîne que
seuls les $i > 0$ interviennent, c'est-à-dire que $h = xh'$ pour un certain $h' \in k[x, y, z, y^{-1}]$.
Alors $(f, g) \sim (f + xh', g - zh')$ et nous voyons que nous pouvons supposer
que $g = 0$ et $h = 0$. Dans ce cas, $xf \in yk[x, y, z, x^{-1}]$
entraîne $f \in yk[x, y, z, x^{-1}]$ et nous voyons que $(f, g) \sim 0$.
Par conséquent, $x : E \to E$ est injective.

\medskip\noindent
Puisque la multiplication par $x$ est un isomorphisme de
$\frac{k[x, y, z, x^{-1}, y^{-1}]}{yk[x, y, z, x^{-1}]}$, nous voyons que
$E/xE$ est isomorphe à
$$
\frac{k[x, y, z, y^{-1}]}{
yzk[x, y, z] + xk[x, y, z, y^{-1}] + zk[x, y, z, y^{-1}]}
=
\frac{k[x, y, z, y^{-1}]}{xk[x, y, z, y^{-1}] + zk[x, y, z, y^{-1}]}
$$
et la multiplication par $y$ est donc un isomorphisme de $E/xE$. Cela
entraîne manifestement (b) et (c).

\medskip\noindent
Soit $e \in E$ la classe d'équivalence de $(1, 0)$.
Supposons que $e \in zE$. Il existe alors $f \in k[x, y, z, x^{-1}, y^{-1}]$,
$g \in k[x, y, z, y^{-1}]$ et $h \in k[x, y, z, y^{-1}]$ tels que
$$
1 + zf + xh \in yk[x, y, z, x^{-1}], \quad
0 + zg - zh \in yzk[x, y, z].
$$
C'est impossible : le monôme $1$ ne peut apparaître dans
$zf$, ni dans $xh$. En revanche, nous avons $ye = 0$ et
$xe = (x, 0) \sim (0, z) = z(0, 1)$. Par conséquent, en prenant pour $\delta$
la classe de $e$ dans $E/zE$, nous obtenons (d).

\begin{lemma}
\label{lemma-strange-regular-sequence}
Il existe un anneau local $R$ et une suite régulière $x, y, z$
(dans l'idéal maximal) tels qu'il existe un élément non nul
$\delta \in R/zR$ vérifiant $x\delta = y\delta = 0$.
\end{lemma}

\begin{proof}
Soit $R = k[x, y, z] \oplus E$, où $E$ est le module ci-dessus, considéré
comme un idéal de carré nul. Il est alors clair que $x, y, z$ est une suite régulière
dans $R$ et que l'élément $\delta \in E/zE \subset R/zR$
a les propriétés voulues. Pour obtenir un exemple local,
nous pouvons localiser $R$ en l'idéal maximal $\mathfrak m = (x, y, z, E)$.
La suite $x, y, z$ reste régulière (car la localisation est
exacte), et l'élément $\delta$ reste non nul puisque son support est réduit
à $\mathfrak m$.
\end{proof}

\begin{lemma}
\label{lemma-base-change-regular-sequence}
Il existe un homomorphisme local d'anneaux locaux $A \to B$
et une suite régulière $x, y$ dans l'idéal maximal de $B$ tels que
$B/(x, y)$ soit plat sur $A$, mais que les images
$\overline{x}, \overline{y}$ de $x, y$ dans $B/\mathfrak m_AB$ ne
forment pas une suite régulière, ni même une suite régulière au sens de Koszul.
\end{lemma}

\begin{proof}
Posons $A = k[z]_{(z)}$ et $B = (k[x, y, z] \oplus E)_{(x, y, z, E)}$.
Puisque $x, y, z$ est une suite régulière dans $B$ (voir la démonstration du
lemme \ref{lemma-strange-regular-sequence}),
nous voyons que $x, y$ est une suite régulière dans $B$ et que
$B/(x, y)$ est un $A$-module sans torsion, donc plat.
En revanche, il existe un élément non nul
$\delta \in B/\mathfrak m_AB = B/zB$ qui est annulé
par $\overline{x}, \overline{y}$. Par conséquent,
$H_2(K_\bullet(B/\mathfrak m_AB, \overline{x}, \overline{y})) \not = 0$.
Ainsi, la suite $\overline{x}, \overline{y}$ n'est pas régulière au sens de Koszul ; en particulier,
elle n'est pas une suite régulière ; voir
Compléments d'algèbre, lemme \ref{more-algebra-lemma-regular-koszul-regular}.
\end{proof}





\section{Un anneau noethérien de dimension infinie}
\label{section-Noetherian-infinite-dimension}

\noindent
Un anneau local noethérien est de dimension finie, comme nous l'avons vu dans
Algèbre, proposition \ref{algebra-proposition-dimension}.
Il existe cependant des anneaux noethériens de dimension infinie.
Voir \cite[Appendix, exemple 1]{Nagata}.

\medskip\noindent
Plus précisément, soit $k$ un corps, et considérons l'anneau
$$
R = k[x_1, x_2, x_3, \ldots ].
$$
Soit $\mathfrak p_i = (x_{2^{i - 1}}, x_{2^{i - 1} + 1}, \ldots, x_{2^i - 1})$
pour $i = 1, 2, \ldots$ ; ce sont des idéaux premiers de $R$. Soit $S$
la partie multiplicativement stable
$$
S = \bigcap\nolimits_{i \geq 1} (R \setminus \mathfrak p_i).
$$
Considérons l'anneau $A = S^{-1}R$.
Nous affirmons que
\begin{enumerate}
\item Les idéaux maximaux de l'anneau $A$ sont les idéaux
$\mathfrak m_i = \mathfrak p_iA$.
\item Nous avons $A_{\mathfrak m_i} = R_{\mathfrak p_i}$, qui est
un anneau local noethérien de dimension $2^{i - 1}$.
\item L'anneau $A$ est noethérien.
\end{enumerate}
Il est donc clair qu'il s'agit de l'exemple recherché.
Nous omettons les détails.




\section{Anneaux locaux dont le complété n'est pas réduit}
\label{section-local-completion-nonreduced}

\noindent
Dans Algèbre, exemple \ref{algebra-example-bad-dvr-char-p}, nous avons donné l'exemple
d'un anneau local noethérien intègre de caractéristique $p$, noté $R$, et de dimension $1$,
dont le complété n'est pas réduit. Nous présentons ici l'exemple
de \cite[Proposition 3.1]{Ferrand-Raynaud}, qui fournit un anneau semblable
en caractéristique zéro.

\medskip\noindent
Soit $\mathbf{C}\{x\}$ l'anneau des séries entières convergentes sur
le corps $\mathbf{C}$ des nombres complexes. L'anneau de toutes les séries entières
$\mathbf{C}[[x]]$ est son complété. Soit $K = \mathbf{C}\{x\}[1/x]$
le corps des séries de Laurent convergentes. Le $K$-module
$\Omega_{K/\mathbf{C}}$ des différentielles algébriques
de $K$ sur $\mathbf{C}$ est un $K$-espace vectoriel de dimension infinie
(démonstration omise). Nous pouvons choisir $f_n \in x\mathbf{C}\{x\}$,
$n \geq 1$, de telle sorte que
$
\text{d}x, \text{d}f_1, \text{d}f_2, \ldots
$
fassent partie d'une base de $\Omega_{K/\mathbf{C}}$. Nous pouvons donc
trouver une $\mathbf{C}$-dérivation
$$
D : \mathbf{C}\{x\} \longrightarrow \mathbf{C}((x))
$$
telle que $D(x) = 0$ et $D(f_n) = x^{-n}$. Posons
$$
A = \{f \in \mathbf{C}\{x\} \mid D(f) \in \mathbf{C}[[x]]\}
$$
Nous affirmons que
\begin{enumerate}
\item $\mathbf{C}\{x\}$ est entier sur $A$,
\item $A$ est un anneau local intègre,
\item $\dim(A) = 1$,
\item l'idéal maximal de $A$ est engendré par $x$ et $xf_1$,
\item $A$ est noethérien, et
\item le complété de $A$ est égal à l'algèbre des nombres duaux
sur $\mathbf{C}[[x]]$.
\end{enumerate}
Puisque l'algèbre des nombres duaux n'est pas réduite, l'anneau $A$ fournit l'exemple recherché.

\medskip\noindent
Remarquons que, si $0 \not = f \in x\mathbf{C}\{x\}$, alors
nous pouvons écrire $D(f) = h/f^n$ pour un certain $n \geq 0$ et $h \in \mathbf{C}[[x]]$.
Par conséquent, $D(f^{n + 1}/(n + 1)) \in \mathbf{C}[[x]]$
et $D(f^{n + 2}/(n + 2)) \in \mathbf{C}[[x]]$. Nous voyons donc que
$f^{n + 1}, f^{n + 2} \in A$ !
En particulier, l'assertion (1) est vérifiée. Nous en déduisons aussi que
le corps des fractions de $A$ est égal au corps des fractions de
$\mathbf{C}\{x\}$. Il en résulte aussi immédiatement que
$A \cap x\mathbf{C}\{x\}$ est l'ensemble des éléments non inversibles de $A$ ; ainsi,
$A$ est un anneau local intègre de dimension $1$. Si nous pouvons prouver (4),
il s'ensuivra que $A$ est noethérien (démonstration omise).
Supposons que $f \in A \cap x\mathbf{C}\{x\}$. Écrivons
$D(f) = h$, avec $h \in \mathbf{C}[[x]]$. Écrivons $h = c + xh'$
avec $c \in \mathbf{C}$ et $h' \in \mathbf{C}[[x]]$. Alors
$D(f - cxf_1) = c + xh' - c = xh'$. D'autre part,
$f - cxf_1 = xg$ avec $g \in \mathbf{C}\{x\}$, mais le
calcul ci-dessus donne $D(g) = h' \in \mathbf{C}[[x]]$
et donc $g \in A$. Ainsi, $f = cxf_1 + xg \in (x, xf_1)$, comme voulu.

\medskip\noindent
Enfin, pourquoi le complété de $A$ n'est-il pas réduit ? Notons $\hat A$ le
complété de $A$. Celui-ci admet évidemment un homomorphisme surjectif vers le complété
$\mathbf{C}[[x]]$ de $\mathbf{C}\{x\}$, puisque $x \in A$. Notons
cet homomorphisme $\psi : \hat A \to \mathbf{C}[[x]]$.
Nous avons vu ci-dessus que $\mathfrak m_A = (x, xf_1)$,
et donc $D(\mathfrak m_A^n) \subset (x^{n - 1})$ par un calcul
facile. Ainsi, $D : A \to \mathbf{C}[[x]]$ est continue et
donne lieu à une dérivation continue $\hat D : \hat A \to \mathbf{C}[[x]]$
au-dessus de $\psi$. Nous obtenons donc un homomorphisme d'anneaux
$$
\psi + \epsilon \hat D :
\hat A
\longrightarrow
\mathbf{C}[[x]][\epsilon].
$$
Comme $\hat A$ est un anneau local noethérien complet de dimension un, si nous
montrons que cette flèche est surjective, il s'ensuivra que $\hat A$
n'est pas réduit. En fait, cet homomorphisme est un isomorphisme, mais nous en omettons
la vérification. Le sous-anneau $\mathbf{C}[x]_{(x)} \subset A$
donne lieu, par complétion, à un homomorphisme $i : \mathbf{C}[[x]] \to \hat A$ tel
que $\psi \circ i = \text{id}$ et que $\hat D \circ i = 0$
(puisque $D(x) = 0$ par construction). Considérons les éléments $x^nf_n \in A$.
Nous avons
$$
(\psi + \epsilon D)(x^nf_n) = x^n f_n + \epsilon
$$
pour tout $n \geq 1$. La surjectivité résulte aussitôt de ces remarques.





\section{Un autre anneau local dont le complété n'est pas réduit}
\label{section-another-local-completion-nonreduced}

\noindent
Dans cette section, nous construisons un exemple d'anneau local noethérien intègre
de dimension $2$, complet pour la topologie définie par un idéal principal,
tel que le complété d'un localisé ne soit pas réduit.

\medskip\noindent
Soit $p$ un nombre premier. Soit $k$ un corps de caractéristique $p$
tel que $k$ soit de degré infini sur son sous-corps $k^p$ des puissances $p$-ièmes.
Par exemple, $k = \mathbf{F}_p(t_1, t_2, t_3, \ldots)$.
Considérons l'anneau
$$
A =
\left\{
\begin{matrix}
\sum a_{i, j} x^iy^j \in k[[x, y]] \text{ tel que, pour tout }n \geq 0
\text{ on ait } \\
[k^p(a_{n, n}, a_{n, n + 1}, a_{n + 1, n},
a_{n, n + 2}, a_{n + 2, n}, \ldots) : k^p] < \infty
\end{matrix}
\right\}
$$
Ensemblistement, on a
$$
k^p[[x, y]] \subset A  \subset k[[x, y]]
$$
Tout élément $f$ de $A$ s'écrit de manière unique sous la forme d'une série
$$
f = f_0 + f_1 xy + f_2 (xy)^2 + f_3 (xy)^3 + \ldots
$$
où
$$
f_n = a_{n, n} + a_{n, n + 1} y + a_{n + 1, n} x +
a_{n, n + 2} y^2 + a_{n + 2, n} x^2 + \ldots
$$
et la condition de la formule qui définit $A$ signifie que les
coefficients de $f_n$ engendrent une extension finie de $k^p$.
Cette présentation montre clairement que $A$ est une
sous-$k^p[[x, y]]$-algèbre de $k[[x, y]]$, complète pour la
topologie définie par l'idéal $xy$. De plus, nous avons manifestement
$$
A/xy A = C \times_k D
$$
où $k^p[[x]] \subset C \subset k[[x]]$ et
$k^p[[y]] \subset D \subset k[[y]]$ sont les sous-anneaux
de séries formelles de
Algèbre, exemple \ref{algebra-example-bad-dvr-char-p}.
Ainsi, $C$ et $D$ sont des anneaux de valuation discrète, et nous voyons que $A/ xy A$ est
noethérien. D'après Algèbre, lemme \ref{algebra-lemma-completion-Noetherian}, nous
concluons que $A$ est noethérien. Puisque $\dim(k[[x, y]]) = 2$, le
lemme \ref{algebra-lemma-integral-sub-dim-equal} d'Algèbre
donne $\dim(A) = 2$.

\medskip\noindent
Soit $f = \sum a_i x^i$ une série formelle telle que
$k^p(a_0, a_1, a_2, \ldots)$ soit de degré infini sur $k^p$.
Alors $f \not \in A$, mais $f^p \in A$. Posons
$$
B = A[f] \subset k[[x, y]]
$$
Comme $B$ est une $A$-algèbre finie, $B$ est noethérien.
De plus, $B$ est complet pour la topologie définie par l'idéal engendré par $xy$; voir
Algèbre, lemme \ref{algebra-lemma-completion-tensor}.
En fait, $B$ est libre sur $A$, de base $1, f, f^2, \ldots, f^{p - 1}$;
nous omettons la démonstration.

\medskip\noindent
Nous affirmons que l'anneau
$$
(B_y)^\wedge = (B[1/y])^\wedge = \lim B[1/y]/(xy)^n B[1/y] =
\lim B[1/y] / x^n B[1/y]
$$
n'est pas réduit. En effet, cet anneau est libre sur
$$
(A_y)^\wedge = (A[1/y])^\wedge = \lim A[1/y]/(xy)^n A[1/y] =
\lim A[1/y] / x^n A[1/y]
$$
de base $1, f, \ldots, f^{p - 1}$. Cependant, il existe un
élément $g \in (A_y)^\wedge$ tel que $f^p = g^p$. En effet,
il suffit de prendre $g = \sum a_i x^i$ (la même expression que
celle utilisée pour $f$), qui a un sens dans $(A_y)^\wedge$.
Nous voyons donc que
$$
(B_y)^\wedge = (A_y)^\wedge[f]/(f^p - g^p) \cong
(A_y)^\wedge[\tau]/(\tau^p)
$$
n'est pas réduit. En fait, cet exemple montre un peu plus.
En effet, remarquons que $(A_y)^\wedge$ est un anneau de valuation discrète
d'uniformisante $x$ et dont le corps résiduel est le corps des fractions
de l'anneau de valuation discrète $D$ défini ci-dessus. Nous voyons donc que même
$$
(B_y)^\wedge[1/(xy)] = ((B_y)^\wedge)_{xy}
$$
n'est pas réduit. Cela fournit un exemple du type suivant.

\begin{lemma}
\label{lemma-nonreduced-recompletion}
Il existe un anneau local noethérien intègre de dimension $2$, noté $(B, \mathfrak m)$,
complet pour la topologie définie par un idéal principal $I = (b)$, et un
élément $f \in \mathfrak m$, $f \not \in I$ tels que
le complété $I$-adique $C = (B_f)^\wedge$ du localisé
principal $B_f$ ne soit pas réduit et que, de plus,
$C_b = C[1/b] = (B_f)^\wedge[1/b]$ ne soit pas réduit.
\end{lemma}

\begin{proof}
Voir la discussion qui précède.
\end{proof}








\section{Un anneau local noethérien non caténaire}
\label{section-non-catenary-Noetherian-local}

\noindent
Bien qu'il existe une théorie satisfaisante de la dimension des anneaux locaux noethériens,
il existe des anneaux locaux noethériens non caténaires. On en trouve un exemple dans
\cite[Appendix, exemple 2]{Nagata}. Nous allons en fait présenter cet exemple
dans le cas le plus simple. Plus précisément, nous allons construire un anneau local noethérien intègre $A$
de dimension $2$ qui n'est pas universellement caténaire. (Remarquons que $A$ est
automatiquement caténaire; voir
Exercices, exercice
\ref{exercises-exercise-Noetherian-local-domain-dim-2-catenary}.)
L'existence d'un anneau local noethérien qui n'est pas universellement
caténaire entraîne celle d'un anneau local noethérien qui n'est pas
caténaire; nous l'expliciterons à la fin de cette section
dans l'exemple particulier considéré.

\medskip\noindent
Soit $k$ un corps, et considérons l'anneau de séries formelles
$k[[x]]$ en une indéterminée sur $k$. Soit
$$
z = \sum\nolimits_{i = 1}^\infty a_i x^i
$$
une série formelle. Nous supposons que $z$, considéré comme élément du corps
de séries de Laurent $k((x)) = k[[x]][1/x]$, est transcendant sur $k(x)$.
Posons
$$
z_j
=
x^{-j}(z - \sum\nolimits_{i = 1, \ldots, j - 1} a_i x^i)
=
\sum\nolimits_{i = j}^\infty a_i x^{i - j}
\in k[[x]].
$$
Remarquons que $z = xz_1$.
Soit $R$ le sous-anneau de $k[[x]]$ engendré par $x$, $z$ et tous les
$z_j$; autrement dit,
$$
R = k[x, z_1, z_2, z_3, \ldots ] \subset k[[x]].
$$
Considérons les idéaux $\mathfrak m = (x)$ et
$\mathfrak n = (x - 1, z_1, z_2 + a_1, z_3 + a_1 + a_2, \ldots)$ de $R$.

\medskip\noindent
Nous avons $xz_{j + 1} + a_j = z_j$. Ainsi, $R/\mathfrak m = k$
et $\mathfrak m$ est un idéal maximal. De plus, tout élément de $R$
qui n'appartient pas à $\mathfrak m$ a pour image un élément inversible de $k[[x]]$, et donc
$R_{\mathfrak m} \subset k[[x]]$. En fait, on en déduit facilement
que $R_{\mathfrak m}$ est un anneau de valuation discrète de corps
résiduel $k$.

\medskip\noindent
Nous affirmons que
$$
R/(x - 1) =
k[x, z_1, z_2, z_3, \ldots ]/(x - 1)
\cong
k[z].
$$
En effet, la relation ci-dessus implique que
$z_{j + 1} = z_j - a_j - (x - 1)z_{j + 1}$, et nous pouvons donc
exprimer la classe de $z_{j + 1}$ en fonction de $z_j$ dans
le quotient $R/(x - 1)$. Puisque le corps des fractions de $R$
est de degré de transcendance $2$ sur $k$ par construction, nous voyons que $z$ est
transcendant sur $k$ dans $R/(x - 1)$, d'où l'isomorphisme voulu.
Ainsi, $\mathfrak n = (x - 1, z)$ est un idéal maximal. En fait,
l'homomorphisme
$$
k[x, x^{-1}, z]_{(x - 1, z)} \longrightarrow R_{\mathfrak n}
$$
est un isomorphisme (puisque $x^{-1}$ est inversible dans $R_{\mathfrak n}$
et que $z_{j + 1} = x^{-1}z_j - a_j = \ldots = f_j(x, x^{-1}, z)$).
Cela montre que $R_{\mathfrak n}$ est un anneau local régulier
de dimension $2$ et de corps résiduel $k$.

\medskip\noindent
Soit $S$ la partie multiplicativement stable
$$
S =
(R \setminus \mathfrak m) \cap (R \setminus \mathfrak n) =
R \setminus (\mathfrak m \cup \mathfrak n)
$$
et posons $B = S^{-1}R$. Nous affirmons que
\begin{enumerate}
\item L'anneau $B$ est une $k$-algèbre.
\item Les idéaux maximaux de l'anneau $B$ sont les deux idéaux
$\mathfrak mB$ et $\mathfrak nB$.
\item Le corps résiduel en chacun de ces idéaux maximaux est $k$.
\item Nous avons $B_{\mathfrak mB} = R_{\mathfrak m}$
et $B_{\mathfrak nB} = R_{\mathfrak n}$,
qui sont des anneaux locaux noethériens réguliers de dimensions $1$ et $2$.
\item L'anneau $B$ est noethérien.
\end{enumerate}
Nous omettons les détails des vérifications.

\medskip\noindent
Étant donnée une $k$-algèbre $B$ possédant les propriétés ci-dessus, nous
obtenons un exemple comme suit. Posons $A = k + \text{rad}(B) \subset B$,
où $\text{rad}(B) = \mathfrak mB \cap \mathfrak nB$ est le radical de Jacobson.
On voit facilement que $B$ est une $A$-algèbre finie, et donc que $A$ est
noethérien par le théorème d'Eakin (voir \cite{Eakin}, ou
\cite[Appendix A1]{Nagata}, ou insérer ici une référence ultérieure).
De plus, $A$ est un anneau local intègre ayant le même corps des fractions que $B$ et
pour corps résiduel $k$. Puisque la dimension de $B$ est $2$, nous voyons que $A$
est également de dimension $2$, d'après
Algèbre, lemme \ref{algebra-lemma-integral-sub-dim-equal}.

\medskip\noindent
Si $A$ était universellement caténaire, la formule de dimension,
Algèbre, lemme \ref{algebra-lemma-dimension-formula},
donnerait $\dim(B_{\mathfrak mB}) = 2$, ce qui est contradictoire.

\medskip\noindent
Remarquons que $B$ est une $A$-algèbre engendrée par un élément.
Ainsi, $B = A[x]/\mathfrak p$ pour un idéal premier
$\mathfrak p$ de $A[x]$. Soit $\mathfrak m' \subset A[x]$
l'idéal maximal correspondant à $\mathfrak mB$. Alors, d'une
part, $\dim(A[x]_{\mathfrak m'}) = 3$ et, d'autre
part,
$$
(0)
\subset \mathfrak pA[x]_{\mathfrak m'}
\subset \mathfrak m'A[x]_{\mathfrak m'}
$$
est une chaîne maximale d'idéaux premiers. Ainsi, $A[x]_{\mathfrak m'}$ est
un exemple d'anneau local noethérien non caténaire.





\section{Existence d'anneaux locaux noethériens pathologiques}
\label{section-bad}

\noindent
Soit $(A, \mathfrak m, \kappa)$ un anneau local noethérien complet.
Dans \cite{Lech}, il est démontré que $A$ est le complété d'un anneau
local noethérien intègre si $\text{profondeur}(A) \geq 1$ et si $A$ contient soit
$\mathbf{Q}$, soit $\mathbf{F}_p$ comme sous-anneau, ou bien contient $\mathbf{Z}$
comme sous-anneau et $A$ est sans torsion comme $\mathbf{Z}$-module.
Cela fournit de nombreux exemples d'anneaux locaux noethériens intègres aux
propriétés « bizarres ».

\medskip\noindent
En appliquant ceci, par exemple, à $A = \mathbf{C}[[x, y]]/(y^2)$, nous obtenons
un anneau local noethérien intègre dont le complété n'est pas réduit.
Comparer avec la
section \ref{section-local-completion-nonreduced}.

\medskip\noindent
Dans \cite{LLPY}, on a trouvé des conditions qui caractérisent les cas où $A$ est
le complété d'un anneau local noethérien réduit.

\medskip\noindent
Dans \cite{Heitmann-completion-UFD}, il est démontré que $A$ est le complété
d'un anneau local noethérien factoriel $R$ si $\text{profondeur}(A) \geq 2$ et si $A$ contient
soit $\mathbf{Q}$, soit $\mathbf{F}_p$ comme sous-anneau, ou bien contient $\mathbf{Z}$
comme sous-anneau et $A$ est sans torsion comme $\mathbf{Z}$-module.
En particulier, $R$ est normal (Algèbre, lemme \ref{algebra-lemma-UFD-normal});
le hensélisé de $R$ est donc lui aussi un anneau normal intègre
(Compléments d'algèbre, lemme \ref{more-algebra-lemma-henselization-normal}).
Ainsi, $A$ comme ci-dessus est le complété d'un anneau local noethérien
hensélien, normal et intègre (puisque les complétés de $R$ et de son hensélisé coïncident;
voir Compléments d'algèbre, lemme \ref{more-algebra-lemma-henselization-noetherian}).

\medskip\noindent
Appliquons ce résultat pour obtenir un anneau local noethérien factoriel $R$ tel que
$R^\wedge \cong \mathbf{C}[[x, y, z, w]]/(wx, wy)$.
Remarquons que $\Spec(R^\wedge)$ est la
réunion d'une composante régulière de dimension $2$ et d'une composante régulière de dimension $3$.
L'anneau $R$ ne peut pas être universellement caténaire. Soit
$$
X \longrightarrow \Spec(R)
$$
l'éclatement de l'idéal maximal. Alors $X$ est un schéma intègre.
Il existe un point fermé $x \in X$ tel que $\dim(\mathcal{O}_{X, x}) = 2$;
plus précisément, au niveau de l'anneau local complet, nous choisissons $x$ sur la
transformée stricte de la composante de dimension $2$, mais non sur la transformée stricte
de la composante de dimension $3$. D'après
Morphismes, lemme \ref{morphisms-lemma-dimension-formula},
nous voyons que $R$ n'est pas universellement caténaire. Comparer avec la
section \ref{section-non-catenary-Noetherian-local}.

\medskip\noindent
L'anneau ci-dessus est caténaire (car c'est un anneau local noethérien factoriel de dimension $3$).
Cependant, dans \cite{Ogoma-example}, l'auteur construit un anneau local
noethérien normal intègre $R$ tel que $R^\wedge \cong \mathbf{C}[[x, y, z, w]]/(wx, wy)$
et que $R$ ne soit pas caténaire. Voir aussi \cite{Heitmann-Ogoma} et
\cite{Lech-YAPO}.

\medskip\noindent
Dans \cite{Heitmann-isolated}, il est démontré que $A$ est le complété
d'un anneau local noethérien $R$ à singularité isolée,
pourvu que $A$ contienne soit $\mathbf{Q}$, soit $\mathbf{F}_p$ comme sous-anneau,
ou que $A$ soit de caractéristique résiduelle $p > 0$ et que $p$ ne puisse devenir un
diviseur de zéro non nul dans aucun localisé propre de $A$.
Nous disons ici qu'un anneau local noethérien $R$
est à singularité isolée si $R_\mathfrak p$ est
un anneau local régulier pour tout idéal premier non maximal $\mathfrak p \subset R$.

\medskip\noindent
Les articles \cite{Nishimura-few} et \cite{Nishimura-few-II}
contiennent de longues listes d'anneaux locaux
noethériens « pathologiques » dont le complété est prescrit. En particulier, ils construisent
un exemple d'anneau local de Nagata, normal et intègre, de dimension $2$, dont le
complété est $\mathbf{C}[[x, y, z]]/(yz)$, et un autre dont le complété
est $\mathbf{C}[[x, y, z]]/(y^2 - z^3)$.

\medskip\noindent
À titre de digression, il est démontré dans \cite{Loepp} que $A$ est le complété d'un
anneau local noethérien excellent et intègre si $A$ est réduit, équidimensionnel,
et si aucun entier de $A$ n'est un diviseur de zéro. Cela ne conduit toutefois pas
à des anneaux locaux noethériens « pathologiques », puisque ceux obtenus sont excellents !




\section{Dimension dans les anneaux noethériens de Jacobson}
\label{section-noetherian-jacobson}

\noindent
Soit $k$ la clôture algébrique d'un corps fini.
Soient $A = k[x, y]$ et $X = \Spec(A)$.
Soit $C = V(x)$ l'axe des $y$ (on pourrait choisir tout autre
sous-schéma fermé intègre de dimension $1$ de $X$).
Soit $C_1, C_2, C_3, \ldots$ une énumération
des autres sous-schémas fermés intègres de $X$ de dimension $1$.
Soit $p_1, p_2, p_3, \ldots$ une
énumération des points fermés de $C$.

\medskip\noindent
Affirmation : pour tout $n$, il existe un fermé irréductible
$Z_n \subset X$ de dimension $1$ tel que
$$
\{p_n\} = Z_n \cap (C \cup C_1 \cup \ldots \cup C_n)
$$
ensemblistement. Pour cela, posons
$Y = C \cup C_1 \cup C_2 \cup \ldots \cup C_n$.
C'est un schéma algébrique affine réduit de dimension $1$ sur $k$.
Il suffit de trouver $f \in k[x, y]$ tel que $V(f) \cap  Y = \{p_n\}$
ensemblistement, car nous pouvons alors prendre pour $Z_n$ une composante
irréductible convenable de $V(f)$. Puisque l'application de restriction
$$
k[x, y] \longrightarrow \Gamma(Y, \mathcal{O}_Y)
$$
est surjective, il suffit de trouver une fonction régulière $g$ sur $Y$
dont l'ensemble des zéros est $\{p_n\}$ ensemblistement.
Pour voir que cela est possible, choisissons un diviseur de Cartier effectif
$D \subset Y$ dont le support est $p_n$ (cela est possible d'après
Variétés, lemme \ref{varieties-lemma-complement-codim-1-closed-points}).
Il suffit donc de montrer que $\mathcal{O}_Y(ND) \cong \mathcal{O}_Y$
pour un certain $N > 0$. Or le groupe de Picard d'un schéma algébrique
affine de dimension $1$ sur la clôture algébrique d'un corps fini est de torsion
(insérer ici une référence ultérieure), ce qui démontre l'affirmation.

\medskip\noindent
Choisissons $Z_n$ comme ci-dessus pour tout $n$. Puisque $k[x, y]$ est un anneau factoriel,
nous pouvons écrire $Z_n = V(f_n)$ pour un élément irréductible $f_n \in A$.
Soit $S \subset k[x, y]$ la partie multiplicative engendrée
par $f_1, f_2, f_3, \ldots$. Considérons l'anneau noethérien $B = S^{-1}A$.

\medskip\noindent
Il est clair que l'homomorphisme d'anneaux $A \to B$ identifie les anneaux locaux
correspondants et induit une injection $\Spec(B) \to \Spec(A)$.
De plus, en considérant la courbe $C_1$, nous voyons que
seuls les points de $C \cap C_1$ sont supprimés lors du passage
de $\Spec(A)$ à $\Spec(B)$.
En particulier, $\Spec(B)$ possède une infinité
d'idéaux maximaux correspondant à des idéaux
maximaux de $A$. En revanche, $xB$ est un idéal maximal, car
le spectre de $B/xB$ est constitué d'un unique idéal premier,
puisque nous avons supprimé tous les points fermés de $C = V(x)$ (mais non
le point générique).
Enfin, pour $i \geq 1$, considérons la courbe $C_i$.
Écrivons $C_i = V(g_i)$, où $g_i \in A$ est irréductible.
Si $C_i = Z_n$ pour un certain $n$, alors $g_iB$ est l'idéal unité.
Sinon, tous les points fermés de $C_i$ sauf un nombre fini
subsistent lors du passage de $A$ à $B$ :
plus précisément, seuls les points de
$(Z_1 \cup \ldots \cup Z_{i - 1} \cup C) \cap C_i$
sont supprimés de $C_i$.

\medskip\noindent
La structure du spectre de $B$ décrite ci-dessus
montre que $B$ est un anneau de Jacobson d'après
Algèbre, lemme \ref{algebra-lemma-noetherian-dim-1-Jacobson}.
Les idéaux maximaux sont ceux de $A$
qui appartiennent à $\Spec(B)$ (et ils sont en nombre infini),
ainsi que l'idéal maximal $xB$. Nous constatons donc que nous avons
des anneaux locaux de dimensions $1$ et $2$.

\begin{lemma}
\label{lemma-Noetherian-Jacobson}
Il existe un anneau de Jacobson noethérien intègre et universellement caténaire $B$
possédant des idéaux maximaux $\mathfrak m_1, \mathfrak m_2$ tels que
$\dim(B_{\mathfrak m_1}) = 1$ et $\dim(B_{\mathfrak m_2}) = 2$.
\end{lemma}

\begin{proof}
La construction de $B$ est donnée ci-dessus. Signalons simplement que
$B$ est universellement caténaire d'après
Algèbre, lemme \ref{algebra-lemma-localization-catenary} et
Morphismes, lemme \ref{morphisms-lemma-ubiquity-uc}.
\end{proof}




\section{Espace sous-jacent noethérien, schéma non noethérien}
\label{section-noetherian-not-noetherian}

\noindent
Nous donnons deux exemples montrant qu'un schéma dont l'espace topologique sous-jacent
est noethérien peut ne pas être un schéma noethérien.

\begin{example}
\label{example-many-variables}
Soit $k$ un corps, et soit $A = k[x_1, x_2, x_3, \dots] /
(x_1^2, x_2^2, x_3^2, \dots)$. Tout idéal premier de $A$ contient les
éléments nilpotents $x_1, x_2, x_3, \dots$; ainsi $\mathfrak p = (x_1, x_2, x_3,
\dots)$ est l'unique idéal premier de $A$. Par conséquent, l'espace
topologique sous-jacent à $\operatorname{Spec} A$ est réduit à un seul point et est
en particulier noethérien. Toutefois, l'idéal $\mathfrak p$ n'est manifestement pas
de type fini.
\end{example}

\begin{example}
\label{example-many-mononials}
Soit $k$ un corps, et soit $A \subseteq k[x, y]$ le
sous-anneau engendré par $k$ et les monômes $\{xy^i\}_{i \ge 0}$. Les idéaux
premiers de $A$ qui ne contiennent pas $x$ sont en correspondance biunivoque avec
les idéaux premiers de $A_x \cong k[x, x^{-1}, y]$. Si $\mathfrak p$ est un
idéal premier qui contient $x$, alors il contient tout $xy^i$, $i \ge 0$,
car $(xy^i)^2 = x(xy^{2i}) \in \mathfrak p$ et $\mathfrak p$ est
radical. Par conséquent, $\mathfrak p = (\{xy^i\}_{i \ge 0})$. Ainsi, l'espace
topologique sous-jacent à $\operatorname{Spec} A$ est noethérien, puisqu'il
est constitué des points du schéma noethérien
$\Spec(A[x, x^{-1}, y])$ et de l'idéal premier $\mathfrak p$.
Mais l'anneau $A$ n'est pas noethérien, car $\mathfrak p$
n'est pas de type fini. Notons que,
dans cet exemple, $A$ est également intègre.
\end{example}






\section{Variété non quasi-affine dont la normalisation est quasi-affine}
\label{section-nonquasi-affine}

\noindent
L'existence d'un exemple de ce type est mentionnée dans
\cite[II remarque 6.6.13]{EGA}, qui renvoie au cinquième volume des
EGA pour un tel exemple; or ce cinquième volume n'a pas paru.

\medskip\noindent
Soit $k$ un corps.
Soit $Y = \mathbf{A}^2_k \setminus \{(0, 0)\}$.
Nous allons construire un morphisme fini, surjectif et birationnel
$\pi : Y \longrightarrow X$
où $X$ est une variété sur $k$ telle que $X$ n'est pas quasi-affine.
Plus précisément, considérons les courbes suivantes dans $Y$ :
$$
\begin{matrix}
C_1 & : & x = 0 \\
C_2 & : & y = 0
\end{matrix}
$$
Remarquons que $C_1 \cap C_2 = \emptyset$. Choisissons l'isomorphisme
$\varphi : C_1 \to C_2$, $(0, y) \mapsto (y^{-1}, 0)$.
Nous affirmons qu'il existe un unique morphisme $\pi : Y \to X$ comme ci-dessus
tel que
$$
\xymatrix{
C_1
\ar@<1ex>[rr]^{\text{id}} \ar@<-1ex>[rr]_{\varphi}
& &
Y \ar[r]^\pi & X
}
$$
soit un diagramme coégalisateur dans la catégorie des variétés (et même dans
la catégorie des schémas). Admettons-le provisoirement et
montrons qu'un tel $X$ ne peut être quasi-affine. En effet, il est clair
que nous aurions
$$
\Gamma(X, \mathcal{O}_X) =
\{ f \in k[x, y] \mid f(0, y) = f(y^{-1}, 0)\} =
k \oplus (xy) \subset k[x, y].
$$
En particulier, ces fonctions ne séparent pas les points $(1, 0)$
et $(-1, 0)$, dont les images dans $X$ (comme nous le verrons plus bas) sont distinctes
(si la caractéristique de $k$ n'est pas $2$).

\medskip\noindent
Pour montrer que $X$ existe, considérons l'ouvert de Zariski
$D(x + y) \subset Y$ de $Y$. C'est le spectre
de l'anneau
$k[x, y, 1/(x + y)]$
et les courbes $C_1$, $C_2$ sont entièrement contenues dans
$D(x + y)$. De plus, le morphisme
$$
C_1 \amalg C_2
\longrightarrow
D(x + y) \cap Y = \Spec(k[x, y, 1/(x + y)])
$$
est une immersion fermée. Il résulte de
Compléments d'algèbre, lemme \ref{more-algebra-lemma-fibre-product-finite-type},
que l'anneau
$$
A =
\{f \in k[x, y, 1/(x + y)] \mid f(0, y) = f(y^{-1}, 0)\}
$$
est une $k$-algèbre de type fini. D'autre part, nous avons l'ouvert
$D(xy) \subset Y$ de $Y$, qui est disjoint des courbes $C_1$
et $C_2$. C'est le spectre de l'anneau
$$
B = k[x, y, 1/xy].
$$
Remarquons que $A_{xy} \cong B_{x + y}$ (car $A$ contient manifestement
les éléments $xyP(x, y)$ pour tout polynôme $P$, ainsi que l'élément $xy/(x + y)$).
Le schéma $X$ est obtenu par recollement des schémas affines
$\Spec(A)$ et $\Spec(B)$ au moyen de l'isomorphisme
$A_{xy} \cong B_{x + y}$; il est donc manifestement de type fini sur
$k$. Pour voir qu'il est séparé, il faut montrer que
l'homomorphisme d'anneaux $A \otimes_k B \to B_{x + y}$ est surjectif. Pour cela,
utilisons le fait que $A \otimes_k B$ contient l'élément
$xy/(x + y) \otimes 1/xy$, dont l'image est $1/(x + y)$.
Le morphisme $Y \to X$ est donné par les morphismes naturels
$D(x + y) \to \Spec(A)$ et $D(xy) \to \Spec(B)$.
Comme ils sont tous deux finis, il en résulte que $Y \to X$ est fini,
comme souhaité. Nous omettons la vérification que $X$ est bien le
coégalisateur du diagramme affiché ci-dessus; voir toutefois
(insérer ici une référence ultérieure sur les sommes amalgamées dans la catégorie des schémas
). Remarquons que le morphisme $\pi : Y \to X$ envoie bien
les points $(1, 0)$
et $(-1, 0)$ sur des points distincts de $X$, puisque la
fonction $(x + y^3)/(x + y)^2 \in A$ prend les valeurs
$1/1$ et, respectivement,\ $-1/(-1)^2 = -1$, qui sont toujours distinctes
(sauf en caractéristique $2$ -- au lecteur de choisir d'autres points
en caractéristique $2$). Résumons cette discussion sous
la forme d'un lemme.

\begin{lemma}
\label{lemma-quasi-affine-normalization-not-quasi-affine}
Soit $k$ un corps.
Il existe une variété $X$ dont la normalisation est quasi-affine, mais
qui n'est pas elle-même quasi-affine.
\end{lemma}

\begin{proof}
Voir la discussion ci-dessus et (insérer ici une référence ultérieure sur la normalisation).
\end{proof}



\section{Prise d'images schématiques}
\label{section-scheme-theoretic-image}

\noindent
Soit $k$ un corps. Soit $t$ une indéterminée. Soient $Y = \Spec(k[t])$
et $X = \coprod_{n \geq 1} \Spec(k[t]/(t^n))$. Notons
$f : X \to Y$ le morphisme
défini au moyen de l'immersion fermée $\Spec(k[t]/(t^n)) \to \Spec(k[t])$
pour tout $n \geq 1$. Dans ce cas, nous avons
\begin{enumerate}
\item L'image schématique
(Morphismes, définition \ref{morphisms-definition-scheme-theoretic-image})
de $f$ est $Y$. En revanche, l'image de $f$ est le point fermé
$t = 0$ de $Y$. Ainsi, la partie fermée sous-jacente à
l'image schématique de $f$
n'est pas égale à l'adhérence de l'image de $f$.
\item La formation de l'image schématique ne commute pas
à la restriction au sous-schéma ouvert $V = \Spec(k[t, 1/t]) \subset Y$.
En effet, l'image réciproque de $V$ dans $X$ est vide; l'image schématique
de $f|_{f^{-1}(V)} : f^{-1}(V) \to V$ est donc le schéma vide.
Celui-ci n'est pas égal à $Y \cap V$.
\end{enumerate}




\section{Images de parties localement fermées}
\label{section-non-chevalley}

\noindent
Le théorème de Chevalley affirme que l'image d'une partie constructible
par un morphisme de présentation finie entre schémas affines est constructible; voir
Algèbre, théorème \ref{algebra-theorem-chevalley}, et
Morphismes, section \ref{morphisms-section-constructible}.
Nous allons voir que l'énoncé analogue n'est pas vrai
pour les images de parties localement fermées.

\medskip\noindent
Soit $k$ un corps de caractéristique $0$. Considérons le morphisme de projection
$$
f :
X = \Spec(k[t, x_1, x_2, \ldots, y_1, y_2, \ldots])
\longrightarrow
\Spec(k[x_1, x_2, \ldots, y_1, y_2, \ldots]) = Y
$$
C'est un morphisme de présentation finie.
Soit $Z$ la partie fermée de $X$ définie par
$$
x_1(t - 1) = 0,\quad
x_2(t - 1)(t - 2) = 0,\quad
x_3(t - 1)(t - 2)(t - 3) = 0,\quad \ldots
$$
Soit $U = \bigcup_{j \geq 1} U_j$ l'ouvert de $X$ défini par
$$
U_j = \text{points où }y_j(t - 1)(t - 2) ... (t - j)\text{ est non nul}
$$
Alors nous avons
$$
f(Z \cap U_j) =
\text{points où }x_1, \ldots, x_j\text{ sont nuls et }y_j\text{ est non nul}
$$
Nous affirmons que $B = f(Z \cap U) = \bigcup_{j \geq 1} f(Z \cap U_j)$
n'est pas une réunion finie de parties localement fermées de $Y$.

\medskip\noindent
Démonstration de l'affirmation. Supposons que $B = A_1 \cup \cdots \cup A_m$
soit un recouvrement fini de $B$ par des parties localement fermées de $Y$.
Nous allons montrer par récurrence sur $n$ que $m \geq n$. Le cas initial
$n = 1$ est immédiat, puisque $B$ est non vide. Supposons $n > 1$ et
l'hypothèse de récurrence vraie pour $n - 1$. Comme l'adhérence de $B$ est
$(x_1 = 0)$, l'un des $A_i$ doit contenir une partie ouverte non vide de
$(x_1 = 0)$. Alors $A_i$ doit être ouvert dans $(x_1 = 0)$.
Mais une telle partie ouverte ne peut contenir aucun point où
$y_1 = 0$; en effet, pour les points de $B$, $y_1 = 0$ entraîne $x_2 = 0$,
ce qui montre que $B$ ne contient aucun voisinage de $(x, y)$ dans $(x_1 = 0)$.
Par conséquent, les $m - 1$ parties restantes induisent
un recouvrement constructible de $B \cap (y_1 = 0)$.
Or l'application de décalage à droite
$x_i \mapsto x_{i + 1}$, $y_i \mapsto y_{i + 1}$ identifie $B$
à $B \cap (y_1 = 0)$! L'hypothèse de récurrence donne donc
$m - 1 \geq n - 1$, d'où $m \geq n$. Ceci achève
la démonstration de l'étape de récurrence et établit l'affirmation.

\begin{lemma}
\label{lemma-no-chevalley}
Il existe un morphisme $f : X \to Y$ de présentation finie
entre schémas affines et une partie localement fermée $T$ de $X$
tels que $f(T)$ ne soit pas une réunion finie de parties localement fermées de $Y$.
\end{lemma}

\begin{proof}
Voir la discussion ci-dessus.
\end{proof}










\section[Sous-schéma localement fermé non ouvert dans un fermé]{Un sous-schéma localement fermé qui n'est pas ouvert dans un fermé}
\label{section-strange-immersion}

\noindent
Ceci est une copie de
Morphismes, exemple \ref{morphisms-example-thibaut}.
Voici un exemple d'immersion qui ne se compose pas d'une
immersion ouverte suivie d'une immersion fermée.
Soit $k$ un corps.
Soit $X = \Spec(k[x_1, x_2, x_3, \ldots])$.
Soit $U = \bigcup_{n = 1}^{\infty} D(x_n)$.
Alors $U \to X$ est une immersion ouverte.
Considérons les idéaux
$$
I_n =
(x_1^n, x_2^n, \ldots, x_{n - 1}^n, x_n - 1, x_{n + 1}, x_{n + 2}, \ldots)
\subset
k[x_1, x_2, x_3, \ldots][1/x_n].
$$
Remarquons que $I_n k[x_1, x_2, x_3, \ldots][1/x_nx_m] = (1)$
pour tout $m \not = n$. Par conséquent, les idéaux quasi-cohérents
$\widetilde I_n$ sur $D(x_n)$ coïncident sur $D(x_nx_m)$; plus précisément,
$\widetilde I_n|_{D(x_nx_m)} = \mathcal{O}_{D(x_n x_m)}$ si
$n \not = m$. Ces idéaux se recollent donc en un faisceau quasi-cohérent d'idéaux
$\mathcal{I} \subset \mathcal{O}_U$.
Soit $Z \subset U$ le sous-schéma fermé correspondant à
$\mathcal{I}$. Ainsi, $Z \to X$ est une immersion.

\medskip\noindent
Nous affirmons qu'il est impossible de factoriser $Z \to X$ sous la forme
$Z \to \overline{Z} \to X$, où $\overline{Z} \to X$ est une immersion fermée
et $Z \to \overline{Z}$ est une immersion ouverte. En effet, $\overline{Z}$ devrait
être défini par un idéal $I \subset k[x_1, x_2, x_3, \ldots]$
tel que $I_n = I k[x_1, x_2, x_3, \ldots][1/x_n]$.
Mais le seul élément $f \in k[x_1, x_2, x_3, \ldots]$
qui appartienne à tous les $I_n$ est $0$! Un tel $I$ n'existe donc pas.

\medskip\noindent
Le morphisme $Z \to X$ fournit aussi un exemple de comportement pathologique
des images schématiques d'immersions. En effet, les arguments ci-dessus
montrent que l'image schématique de l'immersion $Z \to X$ est $X$.
En revanche, nous constatons que
\begin{enumerate}
\item $Z$ n'est pas topologiquement dense dans $X$, et
\item l'image schématique de $Z = Z \cap U \to U$ est simplement
$Z$. Elle n'est pas égale à $U \cap X = U$; la formation
de l'image schématique, dans ce cas,
ne commute donc pas aux restrictions aux ouverts.
\end{enumerate}






\section{Non-existence d'ouverts convenables}
\label{section-nonexistence-opens}

\noindent
Cette section complète les résultats de
Propriétés, section \ref{properties-section-finding-affine-opens}.

\medskip\noindent
Soit $k$ un corps et soit $A = k[z_1, z_2, z_3, \ldots]/I$, où $I$
est l'idéal engendré par tous les produits deux à deux
$z_iz_j$, $i \not = j$, $i, j \in \mathbf{N}$. Posons $S = \Spec(A)$.
Soit $s \in S$ le point fermé correspondant à l'idéal
maximal $(z_i)$. Nous affirmons qu'il n'existe aucun ouvert
quasi-compact $V \subset S \setminus \{s\}$ qui soit
dense dans $S \setminus \{s\}$. Remarquons que $S \setminus \{s\} = \bigcup D(z_i)$.
Chaque $D(z_i)$ est ouvert et irréductible, de point générique
$\eta_i$. Nous en déduisons que $\eta_i \in V$ pour tout $i$.
Cependant, tout ouvert affine principal de $S \setminus \{s\}$ est
de la forme $D(f)$, où $f \in (z_1, z_2, \ldots)$. Alors
$f \in (z_1, \ldots, z_n)$ pour un certain $n$, et $D(f)$ ne contient
qu'un nombre fini des points $\eta_i$. Ainsi, $V$ ne peut être quasi-compact.

\medskip\noindent
Soit $k$ un corps et soit $B = k[x, z_1, z_2, z_3, \ldots]/J$, où $J$
est l'idéal engendré par les produits $xz_i$, $i \in \mathbf{N}$, et par
tous les produits deux à deux $z_iz_j$, $i \not = j$, $i, j \in \mathbf{N}$.
Posons $T = \Spec(B)$. Considérons l'ouvert principal $U = D(x)$.
Nous affirmons qu'il n'existe aucun ouvert quasi-compact $V \subset T$ tel que
$V \cap U = \emptyset$ et que $V \cup U$ soit dense dans $T$.
Soit $t \in T$ le point fermé correspondant à l'idéal
maximal $(x, z_i)$. L'adhérence de $U$ dans $T$ est
$\overline{U} = U \cup \{t\}$. Ainsi, $V \subset \bigcup_i D(z_i)$
est un ouvert quasi-compact. Les arguments du paragraphe précédent
montrent que $V$ ne peut être dense dans $\bigcup D(z_i)$.

\begin{lemma}
\label{lemma-complement-of-affine-does-not-contain-qc-dense-open}
Non-existence d'ouverts quasi-compacts dans des schémas affines :
\begin{enumerate}
\item Il existe un schéma affine $S$ et un ouvert affine $U \subset S$
tels qu'il n'existe aucun ouvert quasi-compact $V \subset S$ pour lequel
$U \cap V = \emptyset$ et $U \cup V$ est dense dans $S$.
\item Il existe un schéma affine $S$ et un point fermé $s \in S$ tels que
$S \setminus \{s\}$ ne contienne aucun ouvert quasi-compact dense.
\end{enumerate}
\end{lemma}

\begin{proof}
Voir la discussion ci-dessus.
\end{proof}

\noindent
Soit $X$ le recollement de deux copies du schéma affine $T$ (voir ci-dessus)
le long de l'ouvert affine $U$. Il existe donc un morphisme $\pi : X \to T$ et
$X = U_1 \cup U_2$ tels que $\pi$ envoie isomorphiquement $U_i$ sur $T$ et
$U_1 \cap U_2$ sur $U$. Remarquons que $X$ est quasi-séparé
(d'après Schémas, lemme \ref{schemes-lemma-characterize-quasi-separated})
et quasi-compact. Nous affirmons qu'il n'existe aucun ouvert séparé, dense et
quasi-compact $W \subset X$. Considérons en effet les deux points fermés
$x_1 \in U_1$, $x_2 \in U_2$ qui s'envoient sur le point fermé $t \in T$
introduit ci-dessus. Soit $\tilde \eta \in U_1 \cap U_2$ le point générique
qui s'envoie sur l'unique point générique $\eta$ de $U$.
Remarquons que $\tilde\eta \leadsto x_1$ et $\tilde\eta \leadsto x_2$
se trouvent au-dessus de la spécialisation $\eta \leadsto t$.
Comme $\pi|_W : W \to T$ est séparé, nous en concluons qu'on ne peut avoir simultanément
$x_1$ et $x_2 \in W$ (par le critère valuatif de séparation,
Schémas, lemme \ref{schemes-lemma-valuative-criterion-separatedness}).
Supposons $x_1 \not \in W$. Alors $W \cap U_1$ est un ouvert quasi-compact (car $X$
est quasi-séparé), dense dans $U_1$, qui ne contient pas $x_1$.
Observons maintenant qu'il existe un isomorphisme $(T, t) \cong (S, s)$
de schémas (envoyant $x$ sur $z_1$ et $z_i$ sur $z_{i + 1}$).
Le premier paragraphe de cette section fournit alors une contradiction.

\begin{lemma}
\label{lemma-no-dense-separated-quasi-compact-open-in-qcqs}
Il existe un schéma quasi-compact et quasi-séparé $X$ qui ne contient
aucun ouvert séparé, quasi-compact et dense.
\end{lemma}

\begin{proof}
Voir la discussion ci-dessus.
\end{proof}







\section{Non-existence d'un sous-schéma ouvert quasi-compact dense}
\label{section-nonexistence-qc-dense-open-subscheme}

\noindent
Soit $X$ un espace algébrique quasi-compact et quasi-séparé sur un corps
$k$. Nous savons que le lieu schématique $X' \subset X$ est un sous-espace
ouvert dense; voir
Propriétés des espaces, proposition
\ref{spaces-properties-proposition-locally-quasi-separated-open-dense-scheme}.
En fait, ce résultat vaut lorsque $X$ est raisonnable; voir
Espaces décents, proposition
\ref{decent-spaces-proposition-reasonable-open-dense-scheme}.
Il est naturel de demander si l'on peut trouver un sous-schéma ouvert
quasi-compact dense de $X$. Ce n'est en général pas possible.

\medskip\noindent
Supposons que la caractéristique de $k$ ne soit pas 2.
Soit $B = k[x, z_1, z_2, z_3, \ldots]/J$, où $J$ est l'idéal engendré par
les produits $xz_i$, $i \in \mathbf{N}$, et par tous les produits deux à deux
$z_iz_j$, $i \not = j$, $i, j \in \mathbf{N}$. Posons $U = \Spec(B)$.
Notons $0 \in U$ le point fermé dont toutes les coordonnées sont nulles.
Posons
$$
j : R = \Delta \amalg \Gamma \longrightarrow U \times_k U
$$
où $\Delta$ est l'image du morphisme diagonal de $U$ sur $k$ et
$$
\Gamma = \{((x, 0, 0, 0, \ldots), (-x, 0, 0, 0, \ldots))
\mid x \in \mathbf{A}^1_k, x \not = 0\}.
$$
Il est clair que $s, t : R \to U$ sont \'etales; ainsi,
$j$ est une relation d'équivalence \'etale. Le quotient $X = U/R$
est un espace algébrique (Espaces, théorème \ref{spaces-theorem-presentation}).
Remarquons que $j$ n'est pas une immersion, car
$(0, 0) \in \Delta$ appartient à l'adhérence de $\Gamma$.
Ainsi, $X$ n'est pas un schéma. En revanche, $X$ est quasi-séparé,
car $R$ est quasi-compact. Notons $0_X$ l'image du point $0 \in U$.
Nous affirmons que $X \setminus \{0_X\}$ est un schéma; plus précisément,
$$
X \setminus \{0_X\} =
\Spec\left(k[x^2, x^{-2}]\right) \amalg
\Spec\left(k[z_1, z_2, z_3, \ldots]/(z_iz_j)\right) \setminus \{0\}
$$
(détails omis). D'autre part, nous avons vu dans la
section \ref{section-nonexistence-opens} que le schéma
figurant au membre de droite ne contient
aucun ouvert quasi-compact dense.

\begin{lemma}
\label{lemma-nonexistence-qc-dense-open-subscheme}
Il existe un espace algébrique quasi-compact et quasi-séparé
qui ne contient aucun sous-schéma ouvert quasi-compact dense.
\end{lemma}

\begin{proof}
Voir la discussion ci-dessus.
\end{proof}

\noindent
En utilisant la construction de
Espaces, exemple \ref{spaces-example-non-representable-descent},
de la même manière que nous avons employé ci-dessus celle de
Espaces, exemple \ref{spaces-example-affine-line-involution},
on obtient un espace algébrique quasi-compact, quasi-séparé et
localement séparé qui ne contient aucun sous-schéma ouvert
quasi-compact dense.


\section{Schémas affines au-dessus d'espaces algébriques}
\label{section-embedding-affines}

\medskip\noindent
Supposons que $f : Y \to X$ soit un morphisme de schémas, que $f$
soit localement de type fini et que $Y$ soit affine. Il existe alors une immersion
$Y \to \mathbf{A}^n_X$ de $Y$ dans l'espace affine de dimension $n$ sur $X$.
Voir l'énoncé légèrement plus général de
Morphismes, lemme \ref{morphisms-lemma-quasi-affine-finite-type-over-S}.

\medskip\noindent
Supposons maintenant que $f : Y \to X$ soit un morphisme d'espaces algébriques,
$f$ étant localement de type fini et $Y$ un schéma affine. En général,
il n'existe pas nécessairement d'immersion de $Y$ dans l'espace
affine de dimension $n$ sur $X$.

\medskip\noindent
Un premier contre-exemple (peu élégant) est donné par $Y = \Spec(k)$ et
$X = [\mathbf{A}^1_k/\mathbf{Z}]$, où $k$ est un corps de caractéristique nulle
et $\mathbf{Z}$ agit sur $\mathbf{A}^1_k$ par la translation $(n, t) \mapsto t + n$.
En effet, pour tout morphisme $Y \to \mathbf{A}^n_X$ au-dessus de $X$, le changement de base
au revêtement $\mathbf{A}^1_k$ de $X$ donne une réunion disjointe infinie de copies de
$\mathbf{A}^1_k$ s'envoyant dans $\mathbf{A}^{n + 1}_k$, morphisme qui n'est pas une
immersion.

\medskip\noindent
Un second contre-exemple est $Y = \mathbf{A}^1_k \to X = \mathbf{A}^1_k/R$,
où $R = \{(t, t)\} \amalg \{(t, -t), t \not = 0\}$. En effet, dans
ce cas, le morphisme $Y \to \mathbf{A}^n_X$ serait donné par des
fonctions régulières $f_1, \ldots, f_n$ sur $Y$; le
produit fibré de $Y$ avec le revêtement
$\mathbf{A}^{n + 1}_k \to \mathbf{A}^n_X$
serait donc le schéma
$$
\{(f_1(t), \ldots, f_n(t), t)\} \amalg
\{(f_1(t), \ldots, f_n(t), -t), t \not = 0\}
$$
muni du morphisme évident vers $\mathbf{A}^{n + 1}_k$, qui n'est pas une immersion.
Remarquons que ce contre-exemple a $X$ quasi-séparé.

\begin{lemma}
\label{lemma-cannot-embed-into-affine}
Il existe un morphisme de type fini d'espaces algébriques $Y \to X$
tel que $Y$ soit affine et $X$ quasi-séparé, pour lequel il n'existe
aucune immersion $Y \to \mathbf{A}^n_X$ au-dessus de $X$.
\end{lemma}

\begin{proof}
Voir la discussion ci-dessus.
\end{proof}










\section{Image directe de modules quasi-cohérents}
\label{section-push-quasi-coherent}

\noindent
Dans Schémas, lemme \ref{schemes-lemma-push-forward-quasi-coherent},
nous avons montré que $f_*$ envoie les modules quasi-cohérents sur des modules
quasi-cohérents lorsque $f$ est quasi-compact et quasi-séparé. Voici des
exemples montrant que ces deux conditions sont nécessaires.

\medskip\noindent
Supposons que $Y = \Spec(A)$ soit un schéma affine et que
$X = \coprod_{n \in \mathbf{N}} Y$. Nous affirmons que $f_*\mathcal{O}_X$
n'est pas quasi-cohérent, où $f : X \to Y$ est le morphisme évident.
En effet, pour $a \in A$, nous avons
$$
f_*\mathcal{O}_X(D(a)) = \prod\nolimits_{n \in \mathbf{N}} A_a
$$
Pour que $f_*\mathcal{O}_X$ soit quasi-cohérent, il faudrait donc avoir
$$
\prod\nolimits_{n \in \mathbf{N}} A_a
=
\left(\prod\nolimits_{n \in \mathbf{N}} A\right)_a
$$
pour tout $a \in A$. Ce n'est pas vrai en général : si, par exemple,
$A = \mathbf{Z}$ et $a = 2$, alors $(1, 1/2, 1/4, 1/8, \ldots)$
est un élément du membre de gauche qui n'appartient pas à celui de droite.
Remarquons que $f$ est un morphisme séparé non quasi-compact.

\medskip\noindent
Soit $k$ un corps. Posons
$$
A = k[t, z, x_1, x_2, x_3, \ldots]/(tx_1z, t^2x_2^2z, t^3x_3^3z, \ldots)
$$
Soit $Y = \Spec(A)$. Soit $V \subset Y$ le sous-schéma ouvert
$V = D(x_1) \cup D(x_2) \cup \ldots$. Soit $X$ le recollement de deux copies de $Y$
le long de $V$. Soit $f : X \to Y$ le morphisme évident. Nous avons alors
une suite exacte
$$
0 \to f_*\mathcal{O}_X \to
\mathcal{O}_Y \oplus \mathcal{O}_Y \xrightarrow{(1, -1)} j_*\mathcal{O}_V
$$
où $j : V \to Y$ est le morphisme d'inclusion. Comme
$$
A \longrightarrow \prod A_{x_n}
$$
est injectif (détails omis), nous obtenons $\Gamma(Y, f_*\mathcal{O}_X) = A$.
En revanche, le noyau de l'application
$$
A_t \longrightarrow \prod A_{tx_n}
$$
est non nul, car il contient l'élément $z$. Par conséquent,
$\Gamma(D(t), f_*\mathcal{O}_X)$ est strictement plus grand que
$A_t$, puisqu'il contient $(z, 0)$. Ainsi, $f_*\mathcal{O}_X$
n'est pas quasi-cohérent. Remarquons que $f$ est quasi-compact mais
non quasi-séparé.

\begin{lemma}
\label{lemma-pushforward-quasi-coherent}
Schémas, lemme \ref{schemes-lemma-push-forward-quasi-coherent},
est optimal au sens où l'on ne peut supprimer ni l'hypothèse
de quasi-compacité ni celle de quasi-séparation.
\end{lemma}

\begin{proof}
Voir la discussion ci-dessus.
\end{proof}








\section{Un module non de type fini dont les germes sont libres de rang 1}
\label{section-nonfree}

\noindent
Soit $R = \mathbf{Q}[x]$. Posons
$M = \sum_{n \in \mathbf{N}} \frac{1}{x - n}R$, considéré comme sous-module du
corps des fractions de $R$. Alors $M$ n'est pas de type fini, mais,
pour tout idéal premier $\mathfrak p$ de $R$, nous avons
$M_{\mathfrak p} \cong R_{\mathfrak p}$ comme
$R_{\mathfrak p}$-modules.

\medskip\noindent
Un exemple du même genre est donné par $R = \mathbf{Z}$ et
$M = \sum_{p \text{ premier}} \frac{1}{p} \mathbf{Z} \subset \mathbf{Q}$,
qui est l'ensemble des fractions $\frac{a}{b}$ où $b$ est non nul et
sans facteur carré.








\section{Un idéal non inversible mais inversible dans les germes}
\label{section-locally-invertible-not-invertible}

\noindent
Soit $A$ un anneau intègre et soit $I \subset A$ un idéal non nul.
Rappelons que, lorsque nous disons que $I$ est inversible, nous entendons que
$I$ est inversible comme $A$-module.
Nous allons construire un exemple de cette situation où
$I$ n'est pas inversible, bien que $I_\mathfrak q = (f) \subset A_\mathfrak q$
soit un idéal principal non nul pour tout idéal premier $\mathfrak q \subset A$.
Dans la littérature, la propriété que $I_\mathfrak q$ soit principal
pour tout idéal premier $\mathfrak q$ s'exprime parfois en
disant que « $I$ est un idéal localement principal ». Nous ne pouvons employer
cette terminologie, car « local » signifie toujours ici
« local pour la topologie de Zariski » (ou pour toute autre topologie dans
laquelle nous travaillons).

\medskip\noindent
Soient $R = \mathbf{Q}[x]$ et $M = \sum \frac{1}{x - n}R$
le module construit dans la section \ref{section-nonfree}.
Considérons l'anneau\footnote{L'anneau $A$ est un exemple
d'anneau intègre non noethérien dont les anneaux locaux sont noethériens.}
$$
A = \text{Sym}^*_R(M)
$$
et l'idéal $I = M A = \bigoplus_{d \geq 1} \text{Sym}^d_R(M)$.
Comme $M$ n'est pas de type fini comme $R$-module,
$I$ ne peut être engendré par un nombre fini d'éléments comme idéal de $A$.
Puisqu'un module inversible est de type fini, il s'ensuit que
$I$ n'est pas inversible.
D'autre part, soit $\mathfrak p \subset R$ un idéal premier.
Par construction, $M_\mathfrak p \cong R_\mathfrak p$. Ainsi,
$$
A_\mathfrak p = \text{Sym}^*_{R_\mathfrak p}(M_\mathfrak p) \cong
\text{Sym}^*_{R_\mathfrak p}(R_\mathfrak p) =
R_\mathfrak p[T]
$$
comme $R_\mathfrak p$-algèbre graduée. Il s'ensuit que
$I_\mathfrak p \subset A_\mathfrak p$ est engendré par
le non-diviseur de zéro $T$.
Ainsi, pour tout idéal premier $\mathfrak q \subset A$,
$I_\mathfrak q$ est certainement engendré par un seul élément.

\begin{lemma}
\label{lemma-locally-principal-not-invertible}
Il existe un anneau intègre $A$ et un idéal non nul $I \subset A$
tels que $I_\mathfrak q \subset A_\mathfrak q$ soit un idéal
principal pour tout idéal premier $\mathfrak q \subset A$, mais que $I$ ne soit pas un
$A$-module inversible.
\end{lemma}

\begin{proof}
Voir la discussion ci-dessus.
\end{proof}









\section{Un module plat de type fini qui n'est pas projectif}
\label{section-finite-flat-not-projective}

\noindent
Il s'agit d'une copie de
Algèbre, remarque \ref{algebra-remark-warning}.
Il est faux qu'un $R$-module de type fini
plat sur $R$ soit automatiquement projectif. Un contre-exemple
s'obtient en prenant pour $R = \mathcal{C}^\infty(\mathbf{R})$
l'anneau des fonctions indéfiniment différentiables sur
$\mathbf{R}$, et $M = R_{\mathfrak m} = R/I$, où
$\mathfrak m = \{f \in R \mid f(0) = 0\}$ et
$I = \{f \in R \mid \exists \epsilon, \epsilon > 0 :
f(x) = 0\ \forall x, |x| < \epsilon\}$.

\medskip\noindent
Le morphisme $\Spec(R/I) \to \Spec(R)$ est également
un exemple d'immersion fermée plate qui n'est pas ouverte.

\begin{lemma}
\label{lemma-finite-flat-non-projective}
Modules plats pathologiques.
\begin{enumerate}
\item Il existe un anneau $R$ et un $R$-module plat de type fini $M$ qui n'est pas
projectif.
\item Il existe une immersion fermée plate qui n'est pas ouverte.
\end{enumerate}
\end{lemma}

\begin{proof}
Voir la discussion ci-dessus.
\end{proof}



\section{Un module projectif qui n'est pas localement libre}
\label{section-projective-not-locally-free}

\noindent
Nous donnons deux exemples : l'un où le rang est compris entre $0$ et $1$,
l'autre où le rang est $\aleph_0$.

\begin{lemma}
\label{lemma-ideal-generated-by-idempotents-projective}
Soit $R$ un anneau. Soit $I \subset R$ un idéal engendré par
une famille dénombrable d'idempotents. Alors $I$ est projectif
comme $R$-module.
\end{lemma}

\begin{proof}
Écrivons $I = (e_1, e_2, e_3, \ldots)$, où $e_n$ est un idempotent de $R$.
En remplaçant par récurrence $e_{n + 1}$ par $e_n + (1 - e_n)e_{n + 1}$,
nous pouvons supposer que $(e_1) \subset (e_2) \subset (e_3) \subset \ldots$,
et donc que $I = \bigcup_{n \geq 1} (e_n) = \colim_n e_nR$.
Dans ce cas,
$$
\Hom_R(I, M) = \Hom_R(\colim_n e_nR, M)
= \lim_n \Hom_R(e_nR, M) = \lim_n e_nM
$$
Remarquons que les applications de transition $e_{n + 1}M \to e_nM$ sont données
par la multiplication par $e_n$ et sont surjectives. Par conséquent, d'après
Algèbre, lemme \ref{algebra-lemma-ML-exact-sequence},
le foncteur $\Hom_R(I, M)$ est exact, c'est-à-dire que $I$ est un
$R$-module projectif.
\end{proof}

\noindent
Supposons que $P \subset Q$ soit une inclusion de $R$-modules, avec $Q$
est un $R$-module de type fini et $P$ est localement libre; voir
Algèbre, définition \ref{algebra-definition-locally-free}.
Supposons que $Q$ puisse être engendré par $N$ éléments comme $R$-module.
Il résulte alors de
Algèbre, lemme \ref{algebra-lemma-map-cannot-be-injective}
que $P$ est localement libre de type fini (ses parties libres étant de rang
au plus $N$). Dans ce cas, $P$ est un $R$-module de type fini; voir
Algèbre, lemme \ref{algebra-lemma-finite-projective}.

\medskip\noindent
En combinant ceci avec ce qui précède, nous voyons qu'un idéal qui n'est pas
de type fini mais qui est engendré par une famille dénombrable d'idempotents
est projectif sans être localement libre. Un exemple explicite est
$R = \prod_{n \in \mathbf{N}} \mathbf{F}_2$, avec
$I$ l'idéal engendré par les idempotents
$$
e_n = (1, 1, \ldots, 1, 0, \ldots )
$$
où la suite de $1$ est de longueur $n$.

\begin{lemma}
\label{lemma-ideal-projective-not-locally-free}
Il existe un anneau $R$ et un idéal $I$ tels que $I$ soit projectif comme
$R$-module, mais ne soit pas localement libre comme $R$-module.
\end{lemma}

\begin{proof}
Voir ci-dessus.
\end{proof}

\begin{remark}
\label{remark-projective-with-dual-finite-not-finite}
L'exemple ci-dessus où $R = \prod_{n \in \mathbf{N}} \mathbf{F}_2$ et
$I$ est l'idéal engendré par les idempotents
$e_n = (1, 1, \ldots, 1, 0, \ldots )$ fournit aussi un exemple $M = I$
de module projectif qui n'est pas de type fini, mais dont le dual est de type
fini et même libre de type fini. En reprenant la méthode de la démonstration du
lemme \ref{lemma-ideal-generated-by-idempotents-projective},
le lecteur montre en effet que $\Hom_R(I, R) \cong R$.
\end{remark}

\begin{lemma}
\label{lemma-chow-group-product}
Soit $K$ un corps.
Soient $C_i$, $i = 1, \ldots, n$, des courbes lisses, projectives et géométriquement
irréductibles sur $K$. Soit $P_i \in C_i(K)$ un point rationnel et
soit $Q_i \in C_i$ un point tel que $[\kappa(Q_i) : K] = 2$.
Alors $[P_1 \times \ldots \times P_n]$ est non nul dans
$\CH_0(U_1 \times_K \ldots \times_K U_n)$, où $U_i = C_i \setminus \{Q_i\}$.
\end{lemma}

\begin{proof}
Il existe une application degré
$\deg : \CH_0(C_1 \times_K \ldots \times_K C_n) \to \mathbf{Z}$
Comme chaque $Q_i$ est de degré $2$ sur $K$, nous voyons que
tout cycle de dimension zéro dont le support est contenu dans le « bord »
$$
C_1 \times_K \ldots \times_K C_n
\setminus
U_1 \times_K \ldots \times_K U_n
$$
a un degré divisible par $2$.
\end{proof}

\noindent
Nous pouvons construire, à l'aide du lemme précédent, un autre exemple de
module projectif non localement libre. Soient
$C_n$, $n = 1, 2, 3, \ldots$, des courbes lisses, projectives et géométriquement
irréductibles sur $\mathbf{Q}$, chacune munie d'une paire de points
$P_n, Q_n \in C_n$ tels que $\kappa(P_n) = \mathbf{Q}$ et que
$\kappa(Q_n)$ soit une extension quadratique de $\mathbf{Q}$.
Posons $U_n = C_n \setminus \{Q_n\}$; c'est une courbe affine.
Soit $\mathcal{L}_n$ l'inverse du faisceau d'idéaux de $P_n$
sur $U_n$. Remarquons que $c_1(\mathcal{L}_n) = [P_n]$ dans le groupe des
cycles de dimension zéro $\CH_0(U_n)$. Posons $A_n = \Gamma(U_n, \mathcal{O}_{U_n})$.
Soit $L_n = \Gamma(U_n, \mathcal{L}_n)$, qui est un module localement
libre de rang $1$ sur $A_n$. Posons
$$
B_n = A_1 \otimes_{\mathbf{Q}} A_2 \otimes_{\mathbf{Q}} \ldots
\otimes_{\mathbf{Q}} A_n
$$
de sorte que $\Spec(B_n) = U_1 \times \ldots \times U_n$, tous les produits
étant pris sur $\Spec(\mathbf{Q})$. Pour $i \leq n$, posons
$$
L_{n, i} =
A_1 \otimes_{\mathbf{Q}} \ldots \otimes_{\mathbf{Q}} L_i
\otimes_{\mathbf{Q}} \ldots \otimes_{\mathbf{Q}} A_n
$$
qui est un $B_n$-module localement libre de rang $1$. Remarquons qu'il s'agit
aussi des sections globales de $\text{pr}_i^*\mathcal{L}_i$. Posons
$$
B_\infty = \colim_n B_n
\quad\text{et}\quad
L_{\infty, i} = \colim_n L_{n, i}
$$
Enfin, posons
$$
M = \bigoplus\nolimits_{i \geq 1} L_{\infty, i}.
$$
C'est une somme directe de modules localement libres de type fini, donc un module projectif.
Nous affirmons que $M$ n'est pas localement libre. Supposons en effet que
$f \in B_\infty$ soit une fonction non nulle telle que $M_f$ soit libre
sur $(B_\infty)_f$. Soit $e_1, e_2, \ldots$ une base. Choisissons
$n \geq 1$ tel que $f \in B_n$.
Choisissons $m \geq n + 1$ tel que $e_1, \ldots, e_{n + 1}$ appartiennent à
$$
\bigoplus\nolimits_{1 \leq i \leq m} L_{m, i}.
$$
Comme les éléments $e_1, \ldots, e_{n + 1}$ font partie d'une base
après un changement de base fidèlement plat, nous en déduisons que
les classes de Chern
$$
c_i(\text{pr}_1^*\mathcal{L}_1 \oplus \ldots \oplus
\text{pr}_m^*\mathcal{L}_m), \quad i = m, m - 1, \ldots, m - n
$$
sont nulles dans le groupe de Chow de
$$
D(f) \subset U_1 \times \ldots \times U_m
$$
Comme $f$ est l'image réciproque d'une fonction sur $U_1 \times \ldots \times U_n$,
il en résulte en particulier que
$$
c_{m - n}(\mathcal{O}_W^{\oplus n} \oplus
\text{pr}_1^*\mathcal{L}_{n + 1} \oplus \ldots
\oplus \text{pr}_{m - n}^*\mathcal{L}_m) = 0.
$$
sur la variété
$$
W = (C_{n + 1} \times \ldots \times C_m)_K
$$
sur le corps $K = \mathbf{Q}(C_1 \times \ldots \times C_n)$.
Autrement dit, le cycle
$$
[(P_{n + 1} \times \ldots \times P_m)_K]
$$
est nul dans le groupe de Chow des cycles de dimension zéro sur $W$. Ceci contredit le
lemme \ref{lemma-chow-group-product}
ci-dessus, car les points $Q_i$, $n + 1 \leq i \leq m$,
induisent des points correspondants $Q_i'$ sur $(C_i)_K$ et, comme $K/\mathbf{Q}$
est géométriquement irréductible, nous avons $[\kappa(Q_i') : K] = 2$.

\begin{lemma}
\label{lemma-projective-not-locally-free}
Il existe un anneau dénombrable $R$ et un module projectif $M$
qui est somme directe d'une famille dénombrable de modules localement libres de rang $1$,
mais tel que $M$ ne soit pas localement libre.
\end{lemma}

\begin{proof}
Voir ci-dessus.
\end{proof}







\section{Anneau local de dimension zéro possédant un idéal plat non nul}
\label{section-zero-dimensional-flat-ideal}

\noindent
On trouve dans \cite{Lazard} et \cite{Autour}
un exemple d'anneau local de dimension zéro possédant un
idéal plat non nul. En voici la construction. Soit $k$ un corps.
Soient $X_i, Y_i$, $i \geq 1$, des indéterminées. Prenons
$R = k[X_i, Y_i]/(X_i - Y_i X_{i + 1}, Y_i^2)$. Notons $x_i$, respectivement,\ $y_i$
l'image de $X_i$, respectivement,\ $Y_i$ dans cet anneau. Remarquons que
$$
x_i = y_i x_{i + 1} = y_i y_{i +1} x_{i + 2} =
y_i y_{i + 1} y_{i + 2} x_{i + 3} = \ldots
$$
dans cet anneau. L'anneau $R$ ne possède qu'un
seul idéal premier, à savoir $\mathfrak m = (x_i, y_i)$. Nous affirmons que
l'idéal $I = (x_i)$ est plat comme $R$-module.

\medskip\noindent
Remarquons que l'annulateur de $x_i$ dans $R$ est l'idéal
$(x_1, x_2, x_3, \ldots, y_i, y_{i + 1}, y_{i + 2}, \ldots)$.
Considérons le $R$-module $M$ engendré par des éléments $e_i$, $i \geq 1$, et
les relations $e_i = y_i e_{i + 1}$. Alors $M$ est plat, car c'est la
limite inductive $\colim_i R$ de copies de $R$ dont les applications de transition sont
$$
R \xrightarrow{y_1} R \xrightarrow{y_2} R \xrightarrow{y_3} \ldots
$$
Remarquons que l'annulateur de $e_i$ dans $M$ est l'idéal
$(x_1, x_2, x_3, \ldots, y_i, y_{i + 1}, y_{i + 2}, \ldots)$.
Puisque tout élément de $M$, respectivement,\ de $I$, peut s'écrire
$f e_i$, respectivement,\ $h x_i$, pour certains $f, h \in R$, l'application
$M \to I$, $e_i \to x_i$, est un isomorphisme; ainsi, $I$ est plat.

\begin{lemma}
\label{lemma-zero-dimensional-flat-ideal}
\begin{slogan}
Anneau de dimension zéro possédant un idéal plat.
\end{slogan}
Il existe un anneau local $R$ possédant un unique idéal premier
et un idéal non nul $I \subset R$ qui est un $R$-module plat.
\end{lemma}

\begin{proof}
Voir la discussion ci-dessus.
\end{proof}








\section[Épimorphisme d'anneaux de dimension zéro non surjectif]{Un épimorphisme d'anneaux de dimension zéro qui n'est pas surjectif}
\label{section-epimorphism-not-surjective}

\noindent
On trouve dans \cite{Lazard-deux} et \cite{Autour} l'exemple suivant.
Soit $k$ un corps. Considérons l'homomorphisme d'anneaux
$$
k[x_1, x_2, \ldots, z_1, z_2, \ldots]/
(x_i^{4^i}, z_i^{4^i})
\longrightarrow
k[x_1, x_2, \ldots, y_1, y_2, \ldots]/
(x_i^{4^i}, y_i - x_{i + 1}y_{i + 1}^2)
$$
qui envoie $x_i$ sur $x_i$ et $z_i$ sur $x_iy_i$. Remarquons que $y_i^{4^{i + 1}}$
est nul dans le membre de droite, tandis que $y_1$ ne l'est pas (détails omis).
Cette application n'est pas surjective : on peut considérer ce qui précède comme une application
d'algèbres $\mathbf{Z}$-graduées en posant $\deg(x_i) = -1$, $\deg(z_i) = 0$
et $\deg(y_i) = 1$; il est alors clair que $y_1$ n'appartient pas à l'image.
Enfin, l'application est un épimorphisme, car
$$
y_{i - 1} \otimes 1 = x_i y_i^2 \otimes 1 = y_i \otimes x_i y_i =
x_i y_i \otimes y_i = 1 \otimes x_i y_i^2.
$$
par conséquent, le produit tensoriel du but sur la source est isomorphe
au but.

\begin{lemma}
\label{lemma-epi-not-surjective}
Il existe un épimorphisme d'anneaux locaux de dimension $0$
qui n'est pas une surjection.
\end{lemma}

\begin{proof}
Voir la discussion ci-dessus.
\end{proof}



\section{De type fini, non de présentation finie, plat en un idéal premier}
\label{section-ft-not-fp-flat-at-prime}

\noindent
Soit $k$ un corps. Considérons l'anneau local $A_0 = k[x, y]_{(x, y)}$.
Notons $\mathfrak p_{0, n} = (y + x^n + x^{2n + 1})$. C'est un idéal premier.
Posons
$$
A = A_0[z_1, z_2, z_3, \ldots]/(z_n z_m, z_n(y + x^n + x^{2n + 1}))
$$
Remarquons que $A \to A_0$ est une surjection dont le noyau est un idéal de
carré nul. Ainsi, $A$ est également un anneau local, et les idéaux premiers de $A$
sont en correspondance biunivoque avec ceux de $A_0$.
Notons $\mathfrak p_n$ l'idéal premier de $A$ correspondant à
$\mathfrak p_{0, n}$. Observons que $\mathfrak p_n$ est l'annulateur
de $z_n$ dans $A$. Soit
$$
C = A[z]/(xz^2 + z + y)[\frac{1}{2zx + 1}]
$$
Remarquons que $A \to C$ est un homomorphisme d'anneaux \'etale; voir
Algèbre, exemple \ref{algebra-example-make-standard-smooth}.
Soit $\mathfrak q \subset C$ l'idéal maximal engendré par
$x$, $y$, $z$ et tous les $z_n$. Comme $A \to C$ est plat,
l'annulateur de $z_n$ dans $C$ est $\mathfrak p_nC$. Calculons
\begin{align*}
C/\mathfrak p_n C
& =
A_0[z]/(xz^2 + z + y, y + x^n + x^{2n + 1})[1/(2zx + 1)] \\
& =
k[x]_{(x)}[z]/(xz^2 + z - x^n - x^{2n + 1})[1/(2zx + 1)] \\
& =
k[x]_{(x)}[z]/(z - x^n) \times
k[x]_{(x)}[z]/(xz + x^{n + 1} + 1)[1/(2zx + 1)] \\
& =
k[x]_{(x)} \times k(x)
\end{align*}
car $(z - x^n)(xz + x^{n + 1} + 1) = xz^2 + z - x^n - x^{2n + 1}$.
Nous obtenons donc $\mathfrak p_nC = \mathfrak r_n \cap \mathfrak q_n$,
où $\mathfrak r_n = \mathfrak p_nC + (z - x^n)C$ et
$\mathfrak q_n = \mathfrak p_nC + (xz + x^{n + 1} + 1)C$.
Comme $\mathfrak q_n + \mathfrak r_n = C$, nous avons aussi
$\mathfrak p_nC = \mathfrak r_n \mathfrak q_n$.
Il s'ensuit que $\mathfrak q_n$ est l'annulateur de $\xi_n = (z - x^n)z_n$.
Observons que, d'une part, $\mathfrak r_n \subset \mathfrak q$ et que,
d'autre part, $\mathfrak q_n + \mathfrak q = C$. Cela résulte par exemple du fait
que $\mathfrak q_n$ est un idéal maximal de $C$ distinct de $\mathfrak q$.
De même, $\mathfrak q_n + \mathfrak q_m = C$ pour $n \not = m$.
Posons maintenant
$$
B = \Im(C \longrightarrow C_{\mathfrak q})
$$
Les éléments $\xi_n$ s'envoient sur zéro dans $B$, car $xz + x^{n + 1} + 1$
n'appartient pas à $\mathfrak q$. Notons $\mathfrak q' \subset B$ l'image de
$\mathfrak q$. Par construction, $B$ est une $A$-algèbre de type fini et
$B_{\mathfrak q'} \cong C_{\mathfrak q}$. En particulier,
$B_{\mathfrak q'}$ est plat sur $A$.

\medskip\noindent
Nous affirmons qu'il n'existe aucun élément $g' \in B$, $g' \not \in \mathfrak q'$,
tel que $B_{g'}$ soit de présentation finie sur $A$. Esquissons la démonstration de
cette affirmation. Choisissons un élément $g \in C$ dont l'image est $g' \in B$.
Considérons l'application $C_g \to B_{g'}$. D'après
Algèbre, lemme \ref{algebra-lemma-finite-presentation-independent},
$B_{g'}$ est de présentation finie sur $A$ si et seulement si le noyau
de $C_g \to B_{g'}$ est de type fini. Mais l'élément $g \in C$
n'appartient pas à $\mathfrak q$ et s'envoie donc sur un élément non nul de
$A_0[z]/(xz^2 + z + y)$. Par conséquent, $g$ ne peut appartenir qu'à un nombre fini
des idéaux premiers $\mathfrak q_n$, car les idéaux premiers
$(y + x^n + x^{2n + 1}, xz + x^{n + 1} + 1)$ forment une famille infinie
de points de codimension 1 de l'espace noethérien irréductible de dimension 2
$\Spec(k[x, y, z]/(xz^2 + z + y))$. L'application
$$
\bigoplus\nolimits_{g \not \in \mathfrak q_n} C/\mathfrak q_n
\longrightarrow
C_g, \quad
(c_n) \longrightarrow \sum c_n \xi_n
$$
est injective et son image est le noyau de $C_g \to B_{g'}$. Nous omettons la
démonstration de cet énoncé. (Indication : écrire $A = A_0 \oplus I$ comme $A_0$-module,
où $I$ est le noyau de $A \to A_0$. De même, écrire $C = C_0 \oplus IC$.
Écrire
$IC = \bigoplus Cz_n \cong \bigoplus (C/\mathfrak r_n \oplus C/\mathfrak q_n)$
et étudier l'effet de la multiplication par $g$ sur les facteurs.)
Ceci achève l'esquisse de la démonstration de l'affirmation.
Cela montre également que $B_{g'}$ n'est plat sur $A$ pour aucun $g'$ comme ci-dessus.
En effet, s'il était plat, l'annulateur de l'image de $z_n$ dans
$B_{g'}$ serait $\mathfrak p_nB_{g'}$ et ne contiendrait pas $z - x^n$.

\medskip\noindent
Nous pouvons par conséquent répondre par la négative à une question posée dans
\cite[Part I, Remarques (3.4.7) (\romannumeral 5)]{GruRay}.
En voici un énoncé précis.

\begin{lemma}
\label{lemma-example-raynaud-gruson}
Il existe un anneau local $A$, un homomorphisme d'anneaux de type fini $A \to B$ et un idéal premier
$\mathfrak q$ au-dessus de $\mathfrak m_A$ tels que $B_{\mathfrak q}$ soit plat
sur $A$ et que, pour tout élément $g \in B$, $g \not \in \mathfrak q$,
l'anneau $B_g$ ne soit ni de présentation finie sur $A$, ni plat sur $A$.
\end{lemma}

\begin{proof}
Voir la discussion ci-dessus.
\end{proof}


\section{De type fini et plat, mais non de présentation finie}
\label{section-finite-type-flat-not-finite-presentation}

\noindent
Dans cette section, nous donnons des exemples d'homomorphismes d'anneaux et de morphismes
qui sont de type fini et plats, mais non de présentation finie.

\medskip\noindent
Soit $R$ un anneau possédant un idéal $I$ tel que $R/I$ soit un module
plat de type fini mais non projectif; voir la
section \ref{section-finite-flat-not-projective}
pour un exemple explicite. Cela signifie que $I$ n'est pas
de type fini; voir
Algèbre, lemme \ref{algebra-lemma-finitely-generated-pure-ideal}.
Remarquons que $I = I^2$; voir
Algèbre, lemme \ref{algebra-lemma-pure}.
Dans nos exemples, l'anneau de base sera
$R$ et, par conséquent, le schéma de base $S = \Spec(R)$.

\medskip\noindent
Considérons l'homomorphisme d'anneaux $R \to R \oplus R/I\epsilon$, où
$\epsilon^2 = 0$ par convention. C'est un homomorphisme fini et plat
qui n'est pas de présentation finie. Tous les anneaux des fibres
sont des intersections complètes et sont géométriquement irréductibles.

\medskip\noindent
Soit $A = R[x, y]/(xy, ay; a \in I)$. Remarquons que, comme $R$-module,
nous avons $A = \bigoplus_{i \geq 0} Ry^i \oplus \bigoplus_{j > 0} R/Ix^j$.
Ainsi, $R \to A$ est un homomorphisme d'anneaux plat de type fini qui n'est pas de
présentation finie. Chaque anneau de fibre
est isomorphe soit à $\kappa(\mathfrak p)[x, y]/(xy)$,
soit à $\kappa(\mathfrak p)[x]$.

\medskip\noindent
Nous pouvons transformer l'exemple précédent en un morphisme projectif en
prenant $B = R[X_0, X_1, X_2]/(X_1X_2, aX_2; a \in I)$.
Dans ce cas, $X = \text{Proj}(B) \to S$ est un morphisme propre et plat
qui n'est pas de présentation finie et tel que, pour tout $s \in S$,
la fibre $X_s$ soit isomorphe soit à $\mathbf{P}^1_s$, soit au
sous-schéma fermé de $\mathbf{P}^2_s$ défini par l'annulation de
$X_1X_2$ (c'est une courbe nodale projective de genre arithmétique $0$).

\medskip\noindent
Soit $M = R \oplus R \oplus R/I$. Posons $B = \text{Sym}_R(M)$, l'algèbre
symétrique de $M$. Posons $X = \text{Proj}(B)$.
Alors $X \to S$ est un morphisme propre et plat, non de présentation finie,
tel que, pour $s \in S$, la fibre géométrique soit isomorphe soit à
$\mathbf{P}^1_s$, soit à $\mathbf{P}^2_s$. En particulier, ces fibres sont
lisses et géométriquement irréductibles.

\begin{lemma}
\label{lemma-finite-type-flat-not-finitely-presented}
Il existe des exemples de
\begin{enumerate}
\item homomorphisme d'anneaux plat de type fini, non de présentation finie,
dont les anneaux des fibres sont des intersections complètes géométriquement irréductibles;
\item homomorphisme d'anneaux plat de type fini, non de présentation finie,
dont les anneaux des fibres sont des intersections complètes de dimension 1,
géométriquement connexes et géométriquement réduites;
\item morphisme propre et plat de schémas $X \to S$, non de présentation finie,
dont chaque fibre est isomorphe soit à $\mathbf{P}^1_s$, soit au lieu d'annulation de
$X_1X_2$ dans $\mathbf{P}^2_s$; et
\item morphisme propre et plat de schémas $X \to S$ dont chaque
fibre est isomorphe soit à $\mathbf{P}^1_s$, soit à $\mathbf{P}^2_s$,
qui n'est pas de présentation finie.
\end{enumerate}
\end{lemma}

\begin{proof}
Voir la discussion ci-dessus.
\end{proof}



\section{Topologie d'un homomorphisme d'anneaux de type fini}
\label{section-topology-finite-type}

\noindent
Soit $A \to B$ un homomorphisme local d'anneaux locaux intègres.
Si $A$ est noethérien, si $A \to B$ est essentiellement de type fini
et si $A/\mathfrak m_A \subset B/\mathfrak m_B$ est fini, alors il
existe un idéal premier $\mathfrak q \subset B$, $\mathfrak q \not = \mathfrak m_B$,
tel que $A \to B/\mathfrak q$ soit la localisation d'un homomorphisme d'anneaux
quasi-fini. Voir
Compléments sur les morphismes, lemme
\ref{more-morphisms-lemma-quasi-finite-quasi-section-meeting-nearby-open}.

\medskip\noindent
Dans cette section, nous donnons un exemple montrant que ce résultat devient faux si $A$
n'est plus noethérien. Soit en effet $k$ un corps et posons
$$
A = \{a_0 + a_1 x + a_2 x^2 + \ldots \mid a_0 \in k, a_i \in k((y))
\text{ pour }i\geq 1\}
$$
et
$$
C = \{a_0 + a_1 x + a_2 x^2 + \ldots \mid a_0 \in k[y], a_i \in k((y))
\text{ pour }i\geq 1\}.
$$
L'inclusion $A \to C$ est de type fini, car $C$ est engendré par $y$
sur $A$. Nous affirmons que $A$ est un anneau local d'idéal maximal
$\mathfrak m = \{a_1 x + a_2 x^2 + \ldots \in A\}$, sans autre idéal
premier que $(0)$ et $\mathfrak m$. En effet, un élément
$f = a_0 + a_1 x + a_2 x^2 + \ldots$ de $A$ est inversible dès que
$a_0 \not = 0$. Si $\mathfrak q \subset A$ est un idéal premier non nul
et si $f = a_i x^i + \ldots \in \mathfrak q$, les propriétés des
séries formelles montrent que, pour tout $g \in k((y))$, l'élément
$g^{i + 1} x^{i + 1} \in \mathfrak q$, c'est-à-dire $gx \in \mathfrak q$.
Cela prouve que $\mathfrak q = \mathfrak m$.

\medskip\noindent
Quant au spectre de l'anneau $C$, le même raisonnement montre
que tout idéal premier non nul contient l'idéal premier
$\mathfrak p = \{a_1 x + a_2 x^2 + \ldots \in C\}$, qui se trouve
au-dessus de $\mathfrak m$. Ainsi, le seul idéal premier de $C$ au-dessus de
$(0)$ est $(0)$. Posons $\mathfrak m_C = yC + \mathfrak p$ et
$B = C_{\mathfrak m_C}$. Alors $A \to B$ est l'exemple recherché.

\begin{lemma}
\label{lemma-topology-finite-type}
Il existe un homomorphisme local $A \to B$ d'anneaux locaux intègres qui est
essentiellement de type fini et tel que $A/\mathfrak m_A \to B/\mathfrak m_B$
soit fini, avec la propriété que, pour tout idéal premier
$\mathfrak q \not = \mathfrak m_B$ de $B$, l'homomorphisme d'anneaux
$A \to B/\mathfrak q$ n'est pas la localisation d'un homomorphisme quasi-fini.
\end{lemma}

\begin{proof}
Voir la discussion ci-dessus.
\end{proof}



\section{Pur mais non universellement pur}
\label{section-pure-not-universally}

\noindent
Soit $k$ un corps. Soit
$$
R = k[[x, xy, xy^2, \ldots]] \subset k[[x, y]].
$$
Autrement dit, une série formelle $f \in k[[x, y]]$ appartient à $R$ si et seulement si
$f(0, y)$ est constante. En particulier, $R[1/x] = k[[x, y]][1/x]$
et $R/xR$ est un anneau local dont l'idéal maximal est de carré
nul. Notons $R[y] \subset k[[x, y]]$ l'ensemble des séries formelles
$f \in k[[x, y]]$ telles que $f(0, y)$ soit un polynôme en $y$. Alors
$R \to R[y]$ est un homomorphisme d'anneaux de type fini mais non de présentation finie
qui induit un isomorphisme après inversion de $x$. Il existe aussi une
surjection $R[y]/xR[y] \to k[y]$ dont le noyau est de carré nul. Considérons
l'homomorphisme d'anneaux de présentation finie $R \to S = R[t]/(xt - xy)$.
De nouveau, $R[1/x] \to S[1/x]$ est un isomorphisme et, dans ce cas,
$S/xS \cong (R/xR)[t]/(xy)$ se surjecte sur $k[t]$ avec un noyau nilpotent.
Il existe une surjection $S \to R[y]$, $t \longmapsto y$, qui induit
un isomorphisme après inversion de $x$ et une surjection de noyau nilpotent
modulo $x$. Le noyau de $S \to R[y]$ est donc localement nilpotent.
En particulier, $S \to R[y]$ induit un homéomorphisme universel sur les spectres.

\medskip\noindent
Nous affirmons d'abord que $S$, considéré comme $S$-module, est relativement pur
sur $R$. Comme l'inversion de $x$ donne un isomorphisme, il suffit
de le vérifier en l'idéal maximal $\mathfrak m \subset R$. Puisque
$R$ est complet pour son idéal maximal, il est hensélien;
il suffit donc de vérifier que, pour tout idéal premier $\mathfrak p \subset R$,
$\mathfrak p \not = \mathfrak m$, l'unique idéal premier $\mathfrak q$
de $S$ au-dessus de $\mathfrak p$ satisfait
$\mathfrak mS + \mathfrak q \not = S$. Comme $\mathfrak p \not = \mathfrak m$,
il correspond à un unique idéal premier de $k[[x, y]][1/x]$. Ainsi,
soit $\mathfrak p = (0)$, soit $\mathfrak p = (f)$ pour un certain
élément irréductible $f \in k[[x, y]]$ non associé à $x$
(nous utilisons ici que $k[[x, y]]$ est factoriel -- insérer ici une référence ultérieure).
Dans le premier cas, $\mathfrak q = (0)$ et le résultat est clair. Dans le
second, nous pouvons multiplier $f$ par une unité afin que $f \in R[y]$
(préparation de Weierstrass; détails omis). On voit alors facilement que
$R[y]/fR[y] \cong k[[x, y]]/(f)$; ainsi, $f$ définit un idéal premier
de $R[y]$ et $\mathfrak mR[y] + fR[y] \not = R[y]$.
Comme $S \to R[y]$ induit un homéomorphisme universel sur les spectres, nous en déduisons
également le résultat voulu pour $S$.

\medskip\noindent
Nous affirmons ensuite que $S$ n'est pas universellement relativement pur sur $R$.
Pour le voir, il suffit de trouver un anneau de valuation
$\mathcal{O}$ et un homomorphisme local $R \to \mathcal{O}$ tels que
$\Spec(R[y] \otimes_R \mathcal{O}) \to \Spec(\mathcal{O})$
n'atteigne pas le point fermé de $\Spec(\mathcal{O})$.
De manière équivalente, il faut trouver $\varphi : R \to \mathcal{O}$ tel que
$\varphi(x) \not = 0$ et $v(\varphi(x)) > v(\varphi(xy))$, où $v$
est la valuation de $\mathcal{O}$.
(En effet, cela signifie que la valuation de $y$ est négative.)
Considérons pour cela l'homomorphisme d'anneaux
$$
R
\longrightarrow
\{a_0 + a_1 x + a_2 x^2 + \ldots \mid a_0 \in k[y^{-1}], a_i \in k((y))\}
$$
défini de manière évidente. On peut trouver un anneau de valuation $\mathcal{O}$
dominant le localisé du membre de droite en l'idéal
maximal $(y^{-1}, x)$, ce qui conclut.

\begin{lemma}
\label{lemma-pure-not-universally-pure}
Il existe un morphisme de schémas affines de présentation finie
$X \to S$ et un $\mathcal{O}_X$-module $\mathcal{F}$ de présentation finie
tels que $\mathcal{F}$ soit pur relativement à $S$, mais non universellement
pur relativement à $S$.
\end{lemma}

\begin{proof}
Voir la discussion ci-dessus.
\end{proof}



\section[Homomorphisme formellement lisse non plat]{Un homomorphisme d'anneaux formellement lisse qui n'est pas plat}
\label{section-formally-smooth-nonflat}

\noindent
Soit $k$ un corps. Considérons la $k$-algèbre $k[\mathbf{Q}]$.
C'est la $k$-algèbre ayant pour base les $x_\alpha, \alpha \in \mathbf{Q}$,
dont la multiplication est déterminée par $x_\alpha x_\beta = x_{\alpha + \beta}$.
(En particulier, $x_0 = 1$.)
Considérons l'homomorphisme de $k$-algèbres
$$
k[\mathbf{Q}] \longrightarrow k, \quad
x_\alpha \longmapsto 1.
$$
Il est surjectif, de noyau $J$ engendré par les éléments $x_\alpha - 1$.
Calculons $J/J^2$. Remarquons que la multiplication par $x_\alpha$ sur $J/J^2$
est l'application identité. Notons $z_\alpha$ la classe de $x_\alpha - 1$ modulo
$J^2$. Ces classes engendrent $J/J^2$. Comme
$$
(x_\alpha - 1)(x_\beta - 1) = x_{\alpha + \beta} - x_\alpha - x_\beta + 1 =
(x_{\alpha + \beta} - 1) - (x_\alpha - 1) - (x_\beta - 1)
$$
nous avons $z_{\alpha + \beta} = z_\alpha + z_\beta$ dans $J/J^2$.
Un élément général de $J/J^2$ est de la forme $\sum \lambda_\alpha z_\alpha$,
où $\lambda_\alpha \in k$ (un nombre fini seulement étant non nuls). Si la
caractéristique de $k$ est $p > 0$, alors
$$
0 = pz_{\alpha/p} = z_{\alpha/p} + \ldots + z_{\alpha/p} = z_\alpha
$$
et nous obtenons $J/J^2 = 0$. Si la caractéristique de $k$ est nulle, alors
$$
J/J^2 = \mathbf{Q} \otimes_{\mathbf{Z}} k \cong k
$$
(détails omis), qui n'est pas nul.

\medskip\noindent
Nous affirmons que $k[\mathbf{Q}] \to k$ est un homomorphisme d'anneaux formellement lisse
si la caractéristique de $k$ est positive. Supposons en effet donné un
diagramme commutatif formé des flèches pleines
$$
\xymatrix{
k \ar[r] \ar@{..>}[rd] & A \\
k[\mathbf{Q}] \ar[u] \ar[r]^\varphi & A' \ar[u]
}
$$
où $A' \to A$ est une surjection dont le noyau $I$ est de carré nul.
Pour montrer que $k[\mathbf{Q}] \to k$ est formellement lisse, il faut prouver
que $\varphi$ se factorise par $k$. Comme $\varphi(x_\alpha - 1)$ s'envoie
sur zéro dans $A$, l'application $\varphi$ induit une application
$\overline{\varphi} : J/J^2 \to I$ dont l'annulation est l'obstruction
à la factorisation voulue. Puisque $J/J^2 = 0$ lorsque la caractéristique
est $p > 0$, nous obtenons le résultat voulu : $k[\mathbf{Q}] \to k$ est
formellement lisse dans ce cas. Enfin, cet homomorphisme n'est pas plat :
par exemple, le non-diviseur de zéro $x_2 - 1$ s'envoie sur zéro.

\begin{lemma}
\label{lemma-formally-smooth-nonflat}
Il existe un homomorphisme d'anneaux formellement lisse qui n'est pas plat.
\end{lemma}

\begin{proof}
Voir la discussion ci-dessus.
\end{proof}


\section{Un homomorphisme d'anneaux formellement \'etale qui n'est pas plat}
\label{section-formally-etale-nonflat}

\noindent
Dans cette section, nous donnons un contre-exemple à la dernière phrase de
\cite[0, exemple 19.10.3(i)]{EGA} (ce point ne figurait pas parmi ceux relevés
dans leurs listes ultérieures d'errata). Considérons $A \to A/J$, où $A$ est un anneau local
et $J$ un idéal propre non nul tel que $J^2 = J$ (ainsi, $J$ n'est pas de type
fini); l'anneau de valuation d'un corps non archimédien algébriquement clos,
avec pour $J$ son idéal maximal, fournit un tel $(A, J)$. Ces
homomorphismes quotients non plats sont formellement \'etales. Supposons en effet donné un
diagramme commutatif
$$
\xymatrix{
A/J \ar[r] & R/I \\
A \ar[u] \ar[r]^\varphi & R \ar[u]
}
$$
où $I$ est un idéal de l'anneau $R$ tel que $I^2 = 0$. Alors $A \to R$
se factorise de manière unique par $A/J$, car
$$
\varphi(J) = \varphi(J^2) \subset (\varphi(J)R)^2 \subset I^2 = 0.
$$
Cela fournit également un contre-exemple au cas formellement \'etale
du « théorème de structure » pour les morphismes localement de type fini et formellement \'etales
de \cite[IV, théorème 18.4.6(i)]{EGA}
(mais non à la partie (ii), qui est celle que l'on emploie effectivement
en pratique). L'erreur dans la démonstration de ce dernier énoncé
consiste, à sa toute dernière étape, à invoquer l'énoncé incorrect
\cite[0, exemple 19.3.10(i)]{EGA};
c'est ainsi que le contre-exemple précédent s'y introduit.

\begin{lemma}
\label{lemma-formally-etale-not-presented}
Il existe des homomorphismes d'anneaux formellement \'etales non plats.
\end{lemma}

\begin{proof}
Voir la discussion ci-dessus.
\end{proof}



\section{Un homomorphisme d'anneaux formellement \'etale à complexe cotangent non trivial}
\label{section-formally-etale-nontrivial-cotangent-complex}

\noindent
Soit $k$ un corps. Considérons l'anneau
$$
R = k[\{x_n\}_{n \geq 1}, \{y_n\}_{n \geq 1}]/(
x_1y_1, x_{nm}^m - x_n, y_{nm}^m - y_n)
$$
Soit $A$ le localisé en l'idéal maximal engendré par
tous les $x_n, y_n$, et notons $J \subset A$ l'idéal maximal. Posons $B = A/J$.
Par construction, $J^2 = J$; ainsi, $A \to B$ est formellement \'etale (voir la
section \ref{section-formally-etale-nonflat}).
Nous affirmons que l'élément $x_1 \otimes y_1$ est un élément
non nul du noyau de
$$
J \otimes_A J \longrightarrow J.
$$
En effet, $(A, J)$ est la limite inductive des localisés
$(A_n, J_n)$ des anneaux
$$
R_n = k[x_n, y_n]/(x_n^n y_n^n)
$$
en leurs idéaux maximaux correspondants. Alors
$x_1 \otimes y_1$ correspond à l'élément
$x_n^n \otimes y_n^n \in J_n \otimes_{A_n} J_n$, qui est non nul
(par un calcul explicite que nous omettons). Puisque $\otimes$ commute
aux limites inductives, nous concluons. D'après
\cite[III section 3.3]{cotangent}
nous voyons que $J$ n'est pas faiblement régulier. Par conséquent, d'après
\cite[III Proposition 3.3.3]{cotangent}
nous voyons que le complexe cotangent $L_{B/A}$ n'est pas nul. En fait, nous pouvons
être plus précis. Nous avons $H_0(L_{B/A}) = \Omega_{B/A}$ et
$H_1(L_{B/A}) = 0$, car $J/J^2 = 0$. Mais la suite exacte à cinq termes
de la suite spectrale fondamentale de Quillen
(voir Cotangent, remarque \ref{cotangent-remark-elucidate-ss}, ou
\cite[Corollary 8.2.6]{Reinhard})
et la non-annulation de
$\text{Tor}_2^A(B, B) = \Ker(J \otimes_A J \to J)$ montrent que
$H_2(L_{B/A})$ est non nul.

\begin{lemma}
\label{lemma-formally-etale-nontrivial-cotangent-complex}
Il existe un homomorphisme d'anneaux surjectif et formellement \'etale $A \to B$
tel que $L_{B/A}$ ne soit pas nul.
\end{lemma}

\begin{proof}
Voir la discussion ci-dessus.
\end{proof}






\section[Plat et formellement non ramifié, non formellement \'etale]{Plat et formellement non ramifié n'implique pas formellement \'etale}
\label{section-flat-formally-unramified-not-formally-etale}

\noindent
Dans Compléments sur les morphismes, lemme
\ref{more-morphisms-lemma-unramified-flat-formally-etale}
il est montré qu'un morphisme de schémas non ramifié et plat $X \to S$
est formellement \'etale. Le but de cette section est de donner deux exemples
montrant qu'on ne peut remplacer « non ramifié » par « formellement non ramifié ».
Le premier exploite des propriétés particulières des anneaux parfaits, tandis que le
second montre que le résultat échoue même pour des homomorphismes d'anneaux réguliers noethériens.

\begin{lemma}
\label{lemma-perfect-closure-polynomial-ring}
Soit $A = \mathbb{F}_p[T]$ l'anneau de polynômes en une indéterminée sur
$\mathbb{F}_p$. Notons $A_{perf}$ la clôture parfaite de $A$.
Alors $A \rightarrow A_{perf}$ est plat et formellement non ramifié,
mais non formellement \'etale.
\end{lemma}

\begin{proof}
Remarquons que, pour le Frobenius $F_A : A \to A$, la copie d'arrivée de $A$
est un module libre sur la copie de départ, de base $\{1, T, \dots, T^{p - 1}\}$.
Ainsi, $F_A$ est fidèlement plat et, par conséquent,
$A \to A_{perf}$ l'est aussi, puisqu'il est une limite inductive d'homomorphismes fidèlement plats.
Comme $A_{perf}$ est un anneau parfait, le Frobenius relatif
$F_{A_{perf}/A}$ est surjectif. Autrement dit,
$A_{perf} = A[A_{perf}^p]$, ce qui implique immédiatement
$\Omega_{A_{perf}/A} = 0$. Alors
$A \rightarrow A_{perf}$ est formellement non ramifié d'après
Compléments sur les morphismes, lemme
\ref{more-morphisms-lemma-formally-unramified-differentials}

\medskip\noindent
Il suffit de montrer que $A \rightarrow A_{perf}$
n'est pas formellement lisse. Comme $A$ est une
$\mathbb{F}_p$-algèbre lisse, le complexe cotangent
$L_{A/\mathbb{F}_p} \simeq \Omega_{A/\mathbb{F}_p}[0]$
est concentré en degré $0$; voir
Cotangent, lemme \ref{cotangent-lemma-when-projective}. De plus,
$L_{A_{perf}/\mathbb{F}_p} = 0$ dans $D(A_{perf})$
d'après Cotangent, lemme \ref{cotangent-lemma-perfect-zero}.
Considérons le triangle distingué de complexes cotangents
$$
L_{A/\mathbb{F}_p} \otimes_A A_{perf} \to
L_{A_{perf}/\mathbb{F}_p} \to
L_{A_{perf}/A} \to
(L_{A/\mathbb{F}_p} \otimes_A A_{perf})[1]
$$
dans $D(A_{perf})$; voir Cotangent, section \ref{cotangent-section-triangle}.
Nous obtenons $L_{A_{perf}/A} = \Omega_{A/\mathbb{F}_p} \otimes_A A_{perf}[1]$,
c'est-à-dire que $L_{A_{perf}/A}$ est égal à un module libre de rang $1$ sur $A_{perf}$
placé en degré $-1$. Ainsi, $A \rightarrow A_{perf}$
n'est pas formellement lisse d'après
Compléments sur les morphismes, lemme \ref{more-morphisms-lemma-NL-formally-smooth},
et
Cotangent, lemme \ref{cotangent-lemma-relation-with-naive-cotangent-complex}.
\end{proof}

\noindent
L'exemple suivant porte également sur des anneaux de caractéristique positive, mais il est peut-être
un peu plus surprenant. Son inconvénient est qu'il requiert une meilleure connaissance des
phénomènes de caractéristique $p$ que l'exemple précédent. Rappelons qu'un
anneau $A$ de caractéristique positive est dit $F$-fini si le Frobenius de $A$
est fini.

\begin{lemma}
\label{lemma-completion-etale}
Soit $(A, \mathfrak m, \kappa)$ un anneau local noethérien de caractéristique
$p > 0$ tel que $[\kappa : \kappa^p] < \infty$.
Alors l'application canonique $A \to A^\wedge$ vers le complété de $A$
est plate et formellement non ramifiée. Cependant, si $A$ est régulier mais non
excellent, cette application n'est pas formellement \'etale.
\end{lemma}

\begin{proof}
La platitude du complété est donnée par
Algèbre, lemme \ref{algebra-lemma-completion-flat}.
Pour montrer que l'application est formellement non ramifiée, il
suffit de montrer que $\Omega_{A^\wedge/A} = 0$; voir
Algèbre, lemme \ref{algebra-lemma-characterize-formally-unramified}.

\medskip\noindent
Esquissons une démonstration. Choisissons $x_1, \ldots, x_r \in A$ dont les images forment une $p$-base
$\overline{x}_1, \ldots, \overline{x}_r$ de $\kappa$, c'est-à-dire
telle que $\kappa$ soit engendré minimalement par les $\overline{x}_i$ sur $\kappa^p$.
Choisissons un système minimal de générateurs $y_1, \ldots, y_s$ de $\mathfrak m$.
Pour tout $n$, les éléments $x_1, \ldots, x_r, y_1, \ldots, y_s$ engendrent
$A/\mathfrak m^n$ sur $(A/\mathfrak m^n)^p$ par Frobenius.
Certains détails sont omis. Nous en déduisons que $F : A^\wedge \to A^\wedge$
est fini. Par conséquent, $\Omega_{A^\wedge/A}$ est un $A^\wedge$-module de type fini.
D'autre part, pour tout $a \in A^\wedge$ et tout $n$, il existe
$a_0 \in A$ tel que $a - a_0 \in \mathfrak m^nA^\wedge$.
Nous en déduisons que $\text{d}(a) \in \bigcap \mathfrak m^n \Omega_{A^\wedge/A}$,
ce qui implique que $\text{d}(a)$ est nul d'après
Algèbre, lemme \ref{algebra-lemma-intersect-powers-ideal-module-zero}.
Ainsi, $\Omega_{A^\wedge/A} = 0$.

\medskip\noindent
Supposons $A$ régulier. Alors, par le théorème de structure de Cohen,
$x_1, \ldots, x_r, y_1, \ldots, y_s$ est une $p$-base de l'anneau
$A^\wedge$, c'est-à-dire que nous avons
$$
A^\wedge = \bigoplus\nolimits_{I, J} (A^\wedge)^p
x_1^{i_1} \ldots x_r^{i_r} y_1^{j_1} \ldots y_s^{j_s}
$$
où $I = (i_1, \ldots, i_r)$, $J = (j_1, \ldots, j_s)$ et
$0 \leq i_a, j_b \leq p - 1$. Détails omis. En particulier,
$\Omega_{A^\wedge}$ est un $A^\wedge$-module libre ayant pour base
$\text{d}(x_1), \ldots, \text{d}(x_r), \text{d}(y_1), \ldots, \text{d}(y_s)$.

\medskip\noindent
Si maintenant $A \to A^\wedge$ est formellement \'etale, ou même seulement formellement lisse,
alors $\NL_{A^\wedge/A}$ a une cohomologie nulle en degrés $-1, 0$
d'après Algèbre, proposition \ref{algebra-proposition-characterize-formally-smooth}.
La suite de Jacobi-Zariski
(Algèbre, lemme \ref{algebra-lemma-exact-sequence-NL}) appliquée aux homomorphismes d'anneaux
$\mathbf{F}_p \to A \to A^\wedge$ donne alors un isomorphisme
$\Omega_A \otimes_A A^\wedge \cong \Omega_{A^\wedge}$.
Nous obtenons donc que $\Omega_A$ est libre ayant pour base
$\text{d}(x_1), \ldots, \text{d}(x_r), \text{d}(y_1), \ldots, \text{d}(y_s)$.
En considérant les corps des fractions et en utilisant la normalité de $A$,
nous concluons que $a \in A$ est une puissance $p$-ième si et seulement si son image dans
$A^\wedge$ est une puissance $p$-ième (détails omis; utiliser
Algèbre, lemme \ref{algebra-lemma-derivative-zero-pth-power}).
Une seconde conséquence est que les opérateurs $\partial/\partial x_a$ et
$\partial/\partial y_b$ sont définis sur $A$.

\medskip\noindent
Nous allons montrer que ce qui précède implique que $A$ est fini
sur $A^p$, avec pour $p$-base $x_1, \ldots, x_r, y_1, \ldots, y_s$.
Cela contredira la non-excellence de $A$ par un résultat de
Kunz; voir \cite[Corollary 2.6]{Kun76}.
Soit en effet $a \in A$ et écrivons
$$
a =  \sum\nolimits_{I, J} (a_{I, J})^p
x_1^{i_1} \ldots x_r^{i_r} y_1^{j_1} \ldots y_s^{j_s}
$$
avec $a_{I, J} \in A^\wedge$. Pour achever la démonstration, il suffit de
montrer que $a_{I, J} \in A$. En appliquant l'opérateur
$$
(\partial/\partial x_1)^{p - 1} \ldots
(\partial/\partial x_r)^{p - 1}
(\partial/\partial y_1)^{p - 1} \ldots
(\partial/\partial y_s)^{p - 1}
$$
aux deux membres, nous concluons que $a_{I, J}^p \in A$, où
$I = (p - 1, \ldots, p - 1)$ et $J = (p - 1, \ldots, p - 1)$.
D'après la remarque précédente, cela implique aussi $a_{I, J} \in A$.
Après avoir remplacé $a$ par $a' = a - a_{I, J}^p x^I y^J$,
nous pouvons employer un opérateur différentiel d'ordre abaissé de $1$ pour obtenir
qu'un autre coefficient $a_{I, J}$ appartient à $A$, et ainsi de suite.
\end{proof}

\begin{remark}
\label{remark-reference-existence-regular-nonexcellent-rings}
On peut construire des anneaux réguliers non excellents dont le corps résiduel a une $p$-base finie,
même dans le corps de fonctions de $\mathbb{P}^2_k$, sur un corps de
caractéristique $p$, $k = \overline{k}$. Voir
\cite[$\mathsection 4.1$]{DS18}.
\end{remark}

\noindent
La démonstration du lemme \ref{lemma-completion-etale} montre en fait un peu plus.

\begin{lemma}
\label{lemma-excellent-regular-local-rings}
Soit $(A, \mathfrak m, \kappa)$ un anneau local régulier de caractéristique
$p > 0$. Supposons $[\kappa : \kappa^p] < \infty$. Alors $A$ est excellent
si et seulement si $A \to A^\wedge$ est formellement \'etale.
\end{lemma}

\begin{proof}
L'implication réciproque résulte du lemme \ref{lemma-completion-etale}.
Pour l'implication directe, le
lemme \ref{lemma-completion-etale}
montre déjà que $A \to A^\wedge$ est formellement non ramifié, ou, de manière équivalente,
que $\Omega_{A^\wedge/A}$ est nul.
Il suffit donc de montrer que l'application de complétion est formellement lisse lorsque
$A$ est excellent. Par la désingularisation de N\'eron-Popescu,
$A \to A^\wedge$ s'écrit comme une limite inductive filtrante
de $A$-algèbres lisses
(Lissage d'homomorphismes d'anneaux, théorème \ref{smoothing-theorem-popescu}).
Par conséquent, $\NL_{A^\wedge/A}$ a une cohomologie nulle en degré $-1$.
Ainsi, $A \to A^\wedge$ est formellement lisse d'après Algèbre, proposition
\ref{algebra-proposition-characterize-formally-smooth}.
\end{proof}










\section[Idéaux d'idempotents et localisation]{Idéaux engendrés par des familles d'idempotents et localisation}
\label{section-ideal-locally-idempotents}

\noindent
Soit $R$ un anneau. Considérons l'anneau
$$
B(R) = R[x_n; n \in \mathbf{Z}]/(x_n(x_n - 1), x_nx_m; n \not = m)
$$
On montre facilement que tout idéal premier $\mathfrak q \subset B(R)$
est soit de la forme
$$
\mathfrak q = \mathfrak pB(R) + (x_n; n \in \mathbf{Z})
$$
soit de la forme
$$
\mathfrak q =
\mathfrak pB(R) + (x_n - 1) + (x_m; n \not = m, m \in \mathbf{Z}).
$$
Nous voyons donc que
$$
\Spec(B(R)) =
\Spec(R) \amalg \coprod\nolimits_{n \in \mathbf{Z}} \Spec(R)
$$
où la topologie n'est pas simplement celle de la réunion disjointe. Elle possède les
propriétés suivantes : chaque copie indexée par $n \in \mathbf{Z}$
est un sous-schéma ouvert, à savoir l'ouvert standard $D(x_n)$.
La copie « centrale » de $\Spec(R)$ appartient à l'adhérence de la réunion
de toute famille infinie des autres copies de $\Spec(R)$.
Remarquons que cette dernière copie de $\Spec(R)$ est définie par l'idéal
$(x_n, n \in \mathbf{Z})$, qui est engendré par les idempotents $x_n$.
Nous voyons ainsi que, si $\Spec(R)$ est connexe,
la décomposition précédente est exactement la décomposition de
$\Spec(B(R))$ en composantes connexes.

\medskip\noindent
Prenons ensuite $A = \mathbf{C}[x, y]/((y - x^2 + 1)(y + x^2 - 1))$.
Le spectre de $A$ possède deux composantes irréductibles
$C_1 = \Spec(A_1)$, $C_2 = \Spec(A_2)$,
où $A_1 = \mathbf{C}[x, y]/(y - x^2 + 1)$ et
$A_2 = \mathbf{C}[x, y]/(y + x^2 - 1)$. Remarquons qu'elles sont
paramétrées par $(x, y) = (t, t^2 - 1)$ et $(x, y) = (t, -t^2 + 1)$,
et se rencontrent en $P = (-1, 0)$ et $Q = (1, 0)$. Construisons une version
tordue de $B(A)$ en recollant $B(A_1)$ à $B(A_2)$ de la manière suivante :
au-dessus de $P$, faisons correspondre $x_n \in B(A_1) \otimes \kappa(P)$
à $x_n \in B(A_2) \otimes \kappa(P)$, tandis qu'au-dessus de $Q$,
nous faisons correspondre $x_n \in B(A_1) \otimes \kappa(Q)$
à $x_{n + 1} \in B(A_2) \otimes \kappa(Q)$.
Notons $B^{twist}(A)$ la $A$-algèbre ainsi obtenue.
Détails omis. Par construction,
$B^{twist}(A)$ est localement pour Zariski sur $A$ isomorphe à la version
non tordue. C'est le cas au-dessus des deux ouverts principaux
$\Spec(A) \setminus \{P\}$
et $\Spec(A) \setminus \{Q\}$.
Cependant, le recollement choisi produit assez de « monodromie » pour que
$\Spec(B^{twist}(A))$ soit connexe (détails omis).
Enfin, il existe une copie centrale
$\Spec(A) \to \Spec(B^{twist}(A))$
qui donne un sous-schéma fermé qui est, localement pour Zariski
sur $B^{twist}(A)$, défini par des idéaux engendrés par des idempotents, mais
ne l'est pas globalement (car $B^{twist}(A)$ ne possède aucun idempotent non trivial).

\begin{lemma}
\label{lemma-not-generated-idempotents}
Il existe un schéma affine $X = \Spec(A)$ et un
sous-schéma fermé $T \subset X$ tel que $T$ soit, localement pour Zariski
sur $X$, défini par des idéaux engendrés par des idempotents, mais que
$T$ ne soit pas défini par un idéal engendré par des idempotents.
\end{lemma}

\begin{proof}
Voir ci-dessus.
\end{proof}



\section{Un homomorphisme d'anneaux qui induit des isomorphismes sur les anneaux locaux sans être ind-\'etale}
\label{section-not-ind-etale}

\noindent
Remarquons que l'homomorphisme d'anneaux $R \to B(R)$ construit dans la
section \ref{section-ideal-locally-idempotents} est
une limite inductive de produits finis de copies de $R$. Ainsi, $R \to B(R)$
est ind-Zariski; voir
Cohomologie pro-\'etale, définition \ref{proetale-definition-ind-zariski}.
Considérons ensuite l'homomorphisme d'anneaux $A \to B^{twist}(A)$
construit dans la section \ref{section-ideal-locally-idempotents}.
Comme cet homomorphisme est, localement pour Zariski sur $\Spec(A)$, isomorphe à un
homomorphisme ind-Zariski $R \to B(R)$, nous concluons qu'il induit des isomorphismes sur les anneaux locaux
(voir Cohomologie pro-\'etale, lemme \ref{proetale-lemma-ind-zariski-implies}).
La discussion de la section \ref{section-ideal-locally-idempotents}
montre qu'il existe une section
$B^{twist}(A) \to A$ dont le noyau n'est pas engendré par des idempotents.
Or, si $A \to B^{twist}(A)$ était ind-\'etale, c'est-à-dire si
$B^{twist}(A) = \colim A_i$ avec $A \to A_i$ \'etale,
alors le noyau de $A_i \to A$ serait engendré par un idempotent
(Algèbre, lemmes \ref{algebra-lemma-map-between-etale} et
\ref{algebra-lemma-surjective-flat-finitely-presented}).
Cela contredirait le résultat précédent.

\begin{lemma}
\label{lemma-not-ind-etale}
Il existe un homomorphisme d'anneaux $A \to B$ qui induit des isomorphismes sur les anneaux locaux, mais
qui n'est pas ind-\'etale. A fortiori, il n'est pas ind-Zariski.
\end{lemma}

\begin{proof}
Voir la discussion ci-dessus.
\end{proof}




\section[Faisceau quasi-cohérent non flasque d'un module injectif]{Faisceau quasi-cohérent non flasque associé à un module injectif}
\label{section-nonflasque}

\noindent
Pour d'autres exemples de ce type, voir \cite[Expos\'e II, Appendix I]{SGA6},
où Illusie expose certains exemples dus à Verdier.

\medskip\noindent
Considérons le schéma affine $X = \Spec(A)$, où
$$
A = k[x, y, z_1, z_2, \ldots]/(x^nz_n)
$$
est l'anneau de
Propriétés, exemple \ref{properties-example-does-not-work-in-general}.
Posons $I = (x) \subset A$. Considérons l'ouvert quasi-compact $U = D(x)$ de $X$.
Nous avons vu loc.\ cit.\ qu'il existe une section
$s \in \mathcal{O}_X(U)$ qui ne provient d'aucun homomorphisme de $A$-modules
$I^n \to A$, quel que soit $n \geq 0$.

\medskip\noindent
Soit $\alpha : A \to J$ un plongement de $A$ dans un $A$-module injectif.
Posons $Q = J/\alpha(A)$ et notons $\beta : J \to Q$ l'application quotient.
Nous affirmons que l'application
$$
\Gamma(X, \widetilde{J})
\longrightarrow
\Gamma(U, \widetilde{J})
$$
n'est pas surjective. Plus précisément, nous affirmons que $\alpha(s)$
n'appartient pas à son image. Pour le voir, raisonnons par l'absurde.
Supposons donc qu'il existe un élément $x \in J$ qui se restreigne à
$\alpha(s)$ sur $U$. Alors $\beta(x) \in Q$ se restreint à $0$ sur $U$. Nous savons donc que
$I^n\beta(x) = 0$ pour un certain $n$; voir
Propriétés,
lemme \ref{properties-lemma-sections-over-quasi-compact-open-in-affine}.
Nous obtenons par conséquent un morphisme
$\varphi : I^n \to A$, $h \mapsto \alpha^{-1}(hx)$. Il est immédiat que
ce morphisme $\varphi$ donne la section $s$ par l'application de
Propriétés,
lemme \ref{properties-lemma-sections-over-quasi-compact-open-in-affine},
ce qui est absurde.

\begin{lemma}
\label{lemma-nonflasque}
Il existe un schéma affine $X = \Spec(A)$ et un $A$-module injectif $J$
tels que $\widetilde{J}$ ne soit pas un faisceau flasque sur $X$.
Même la restriction $\Gamma(X, \widetilde{J}) \to \Gamma(U, \widetilde{J})$,
où $U$ est un ouvert standard, n'est pas nécessairement surjective.
\end{lemma}

\begin{proof}
Voir ci-dessus.
\end{proof}

\noindent
En fait, une construction analogue fournit un exemple de module injectif
dont la cohomologie du faisceau quasi-cohérent associé sur un ouvert
quasi-compact n'est pas nulle. Plus précisément, partons de l'anneau
$$
A = k[x, y, w_1, u_1, w_2, u_2, \ldots]/(x^nw_n, y^nu_n, u_n^2, w_n^2)
$$
où $k$ est un corps. Choisissons un homomorphisme injectif $A \to I$, où
$I$ est un $A$-module injectif. Nous affirmons que l'élément $1/xy$ de
$A_{xy} \subset I_{xy}$ n'appartient pas à l'image de
$I_x \oplus I_y \to I_{xy}$. En raisonnant par l'absurde, supposons que
$$
\frac{1}{xy} = \frac{i}{x^n} + \frac{j}{y^n}
$$
pour un certain $n \geq 1$ et des $i, j \in I$. En chassant les dénominateurs, nous obtenons
$$
(xy)^{n + m - 1} = x^my^{n + m}i + x^{n + m}y^mj
$$
pour un certain $m \geq 0$. En multipliant par $u_{n + m}w_{n + m}$,
nous voyons que $u_{n + m}w_{n + m}(xy)^{n + m - 1} = 0$ dans $A$, ce
qui donne la contradiction cherchée.
Posons $U = D(x) \cup D(y) \subset X = \Spec(A)$. Pour tout $A$-module
$M$, nous avons une suite exacte
$$
0 \to H^0(U, \widetilde{M}) \to M_x \oplus M_y \to M_{xy} \to
H^1(U, \widetilde{M}) \to 0
$$
par Mayer--Vietoris. Nous en concluons que $H^1(U, \widetilde{I})$ n'est pas nul.

\begin{lemma}
\label{lemma-nonvanishing}
Il existe un schéma affine $X = \Spec(A)$ dont l'espace
topologique sous-jacent est noethérien et un
$A$-module injectif $I$ tels que $\widetilde{I}$ ait un $H^1$ non nul
sur un ouvert quasi-compact $U$ de $X$.
\end{lemma}

\begin{proof}
Voir ci-dessus. Notons que $\Spec(A) = \Spec(k[x, y])$ comme espaces topologiques.
\end{proof}









\section{Un schéma en groupes plat non séparé}
\label{section-non-separated-group-scheme}

\noindent
Tout schéma en groupes sur un corps est séparé; voir
Groupoïdes, lemme \ref{groupoids-lemma-group-scheme-over-field-separated}.
Ce n'est pas vrai pour les schémas en groupes sur une base.

\medskip\noindent
Soit $k$ un corps. Posons $S = \Spec(k[x]) = \mathbf{A}^1_k$.
Soit $G$ la droite affine dont l'origine $0$ est doublée (voir
Schémas, exemple \ref{schemes-example-affine-space-zero-doubled}),
considérée comme un schéma sur $S$. Ainsi, une fibre de $G \to S$ est soit
réduite à un point, soit un ensemble à deux éléments (l'un dans $U$, l'autre dans $V$).
Nous pouvons donc munir ces fibres d'une structure de groupe (en
prenant l'élément de $U$ pour zéro de la structure de groupe).
Plus précisément, $G$ possède deux ouverts $U, V$ qui se projettent isomorphiquement
sur $S$ et tels que $U \cap V$ se projette isomorphiquement sur
$S \setminus \{0\}$. Alors
$$
G \times_S G = U \times_S U \cup V \times_S U \cup
U \times_S V \cup V \times_S V
$$
où chaque morceau est isomorphe à $S$. Nous pouvons donc définir une multiplication
$m : G \times_S G \to G$ comme l'unique $S$-morphisme qui envoie le premier
et le dernier morceau dans $U$, et les deux morceaux intermédiaires dans $V$. Cela concorde avec
la description fibre par fibre donnée ci-dessus. Nous omettons la
vérification que cela définit une structure de schéma en groupes.

\begin{lemma}
\label{lemma-non-separated-group-scheme}
Il existe sur la droite affine un schéma en groupes plat et de type fini
qui n'est pas séparé.
\end{lemma}

\begin{proof}
Voir la discussion ci-dessus.
\end{proof}

\begin{lemma}
\label{lemma-non-quasi-separated-group-scheme}
Il existe sur l'espace affine de dimension infinie un schéma en groupes
plat et de type fini qui n'est pas quasi-séparé.
\end{lemma}

\begin{proof}
La même construction peut être effectuée avec l'espace affine de dimension infinie
$S = \mathbf{A}^\infty_k = \Spec k[x_1, x_2, \ldots]$ comme base
et avec l'origine $0 \in S$, qui correspond à l'idéal maximal
$(x_1, x_2, \ldots)$, comme point fermé doublé dans $G$.
Le schéma en groupes $G \rightarrow S$ ainsi obtenu
n'est pas quasi-séparé, comme l'explique
Schémas, exemple \ref{schemes-example-not-quasi-separated}.
\end{proof}



\section{Un schéma en groupes non plat à composante neutre plate}
\label{section-non-flat-group-scheme}

\noindent
Soit $X \to S$ un monomorphisme de schémas. Posons $G = S \amalg X$.
Soit $m : G \times_S G \to G$ le $S$-morphisme
$$
G \times_S G = X \times_S X \amalg X \amalg X \amalg S
\longrightarrow G = X \amalg S
$$
qui envoie $X \times_S X$ et $S$ dans $S$, ainsi que chacun des
deux exemplaires de $X$ dans $X$ par le morphisme identité.
Cela définit une loi de groupe. Pour le voir, il faut montrer que
$m \circ (m \times \text{id}_G) = m \circ (\text{id}_G \times m)$
comme applications $G \times_S G \times_S G \to G$. En décomposant
$G \times_S G \times_S G$ en morceaux comme ci-dessus, nous constatons que
l'égalité doit être vérifiée sur chacun des $8$ morceaux.
Chaque morceau est isomorphe à $S$, à $X$, à $X \times_S X$ ou à
$X \times_S X \times_S X$. De plus, les deux applications envoient
respectivement ces morceaux dans $S$, $X$, $S$, $X$. Cela étant, le fait que
$X \to S$ soit un monomorphisme implique que $X \times_S X \cong X$
et $X \times_S X \times_S X \cong X$, et qu'il existe dans chaque cas
un unique $S$-morphisme $S \to S$ ou $X \to X$. Nous voyons donc
que $m \circ (m \times \text{id}_G) = m \circ (\text{id}_G \times m)$.
Ainsi, en prenant pour $X \to S$ n'importe quel
monomorphisme non plat de schémas (par exemple une immersion fermée),
nous obtenons un exemple de schéma en groupes
sur une base $S$ dont la composante neutre est $S$ (donc plate),
mais qui n'est pas plat.

\begin{lemma}
\label{lemma-non-flat-group-scheme}
Il existe un schéma en groupes $G$ sur une base $S$ dont la composante
neutre est plate sur $S$, mais qui n'est pas plat sur $S$.
\end{lemma}

\begin{proof}
Voir la discussion ci-dessus.
\end{proof}




\section{Un espace algébrique en groupes non séparé sur un corps}
\label{section-non-separated-group-space}

\noindent
Tout schéma en groupes sur un corps est séparé; voir
Groupoïdes, lemme \ref{groupoids-lemma-group-scheme-over-field-separated}.
Ce n'est pas vrai pour les espaces algébriques en groupes sur un corps
(mais voir la fin de cette section pour des résultats positifs).

\medskip\noindent
Soit $k$ un corps de caractéristique zéro.
Considérons l'espace algébrique $G = \mathbf{A}^1_k/\mathbf{Z}$ de
Espaces, exemple \ref{spaces-example-affine-line-translation}.
Par construction, $G$ est le faisceau fppf associé au préfaisceau
$$
T \longmapsto \Gamma(T, \mathcal{O}_T) / \mathbf{Z}
$$
sur la catégorie des schémas sur $k$. La loi d'addition évidente du préfaisceau
induit une addition $m : G \times G \to G$ qui fait de $G$ un espace
algébrique en groupes sur $\Spec(k)$. Notons que $G$ n'est pas séparé
(ni même quasi-séparé ou localement séparé). En revanche,
$G \to \Spec(k)$ est de type fini !

\begin{lemma}
\label{lemma-non-separated-group-space}
Il existe sur un corps un espace algébrique en groupes de type fini
qui n'est pas séparé (ni même quasi-séparé ou localement séparé).
\end{lemma}

\begin{proof}
Voir la discussion ci-dessus.
\end{proof}

\noindent
Résultats positifs : si l'espace algébrique en groupes $G$ est
quasi-séparé ou localement séparé, ou plus généralement s'il est
un espace algébrique décent, alors $G$ est en fait séparé; voir
Compléments sur les groupoïdes dans les espaces, lemme
\ref{spaces-more-groupoids-lemma-group-scheme-over-field-separated}.
De plus, un espace algébrique en groupes de type fini et séparé sur un
corps est en fait un schéma d'après Compléments sur les groupoïdes dans les espaces, lemme
\ref{spaces-more-groupoids-lemma-group-space-scheme-locally-finite-type-over-k}.
L'idée de la preuve est que le lieu schématique est ouvert et dense; voir
Propriétés des espaces, proposition
\ref{spaces-properties-proposition-locally-quasi-separated-open-dense-scheme}
ou
Espaces décents, théorème \ref{decent-spaces-theorem-decent-open-dense-scheme}.
En prenant les translatés de cet ouvert, nous voyons que
tout point de $G$ possède un voisinage ouvert qui est un schéma.




\section{Spécialisations entre points d'une même fibre d'un morphisme étale}
\label{section-specializations-fibre-etale}

\noindent
Si $f : X \to Y$ est un morphisme étale ou, plus généralement, localement
quasi-fini de schémas, il n'existe aucune spécialisation entre des points
d'une même fibre; voir
Morphismes, lemme \ref{morphisms-lemma-locally-quasi-finite-fibres}.
En revanche, cela est faux en général pour les morphismes d'espaces algébriques.

\medskip\noindent
Pour construire un exemple, soit $k$ un corps.
Posons
$$
P = k[u, u^{-1}, y, \{x_n\}_{n \in \mathbf{Z}}].
$$
Considérons l'action
de $\mathbf{Z}$ sur $P$ par des homomorphismes de $k$-algèbres, engendrée par l'automorphisme
$\tau$ défini par $\tau(u) = u$, $\tau(y) = uy$ et
$\tau(x_n) = x_{n + 1}$. Pour $d \geq 1$, posons
$I_d = ((1 - u^d)y, x_n - x_{n + d}, n \in \mathbf{Z})$.
Alors $V(I_d) \subset \Spec(P)$ est le lieu des points fixes de $\tau^d$.
Soit $S \subset P$ la partie multiplicativement stable engendrée par $y$ et
tous les $1 - u^d$, $d \in \mathbf{N}$. Nous
voyons alors que $\mathbf{Z}$ agit librement sur $U = \Spec(S^{-1}P)$.
Soit $X = U/\mathbf{Z}$ l'espace algébrique quotient; voir
Espaces, définition \ref{spaces-definition-quotient}.

\medskip\noindent
Considérons les idéaux premiers $\mathfrak p_n = (x_n, x_{n + 1}, \ldots)$ de
$S^{-1}P$. Notons que $\tau(\mathfrak p_n) = \mathfrak p_{n + 1}$.
Chacun d'eux définit donc un point $\xi_n \in U$ dont l'image dans $X$ est
le même point $x$ de $X$. De plus, nous avons les spécialisations
$$
\ldots \leadsto \xi_n \leadsto \xi_{n - 1} \leadsto \ldots
$$
Nous en concluons que $U \to X$ est un exemple du type annoncé.

\begin{lemma}
\label{lemma-specializations-fibre-etale}
Il existe un morphisme étale d'espaces algébriques $f : X \to Y$
et une spécialisation non triviale de points $x \leadsto x'$ dans $|X|$, telle que
$f(x) = f(x')$ dans $|Y|$.
\end{lemma}

\begin{proof}
Voir la discussion ci-dessus.
\end{proof}







\section{Un torseur qui n'est pas un torseur fppf}
\label{section-torsor-not-fppf}

\noindent
Dans
Groupoïdes, remarque \ref{groupoids-remark-fun-with-torsors},
nous demandons si tout torseur sous $G$ est un torseur fppf sous $G$.
Nous montrons dans cette section que ce n'est pas toujours le cas.

\medskip\noindent
Soit $k$ un corps. Dans la suite, tous les schémas et champs sont sur $k$.
Soit $G \to \Spec(k)$ le schéma en groupes
$$
G = (\mu_{2, k})^\infty =
\mu_{2, k} \times_k \mu_{2, k} \times_k \mu_{2, k} \times_k \ldots =
\lim_n (\mu_{2, k})^n
$$
où $\mu_{2, k}$ est le schéma en groupes des racines carrées de l'unité sur
$\Spec(k)$; voir
Groupoïdes, exemple \ref{groupoids-example-roots-of-unity}.
Comme $G$ est une limite projective de schémas affines, c'est un schéma en groupes
affine. C'est en fait le spectre de l'anneau
$k[t_1, t_2, t_3, \ldots]/(t_i^2 - 1)$. L'application de multiplication
$m : G \times_k G \to G$ est donnée au niveau des algèbres par
$t_i \mapsto t_i \otimes t_i$.

\medskip\noindent
Nous affirmons que tout torseur sous $G$ sur $k$ est de la forme
$$
P = \Spec(k[x_1, x_2, x_3, \ldots]/(x_i^2 - a_i))
$$
pour certains $a_i \in k^*$; l'action de $G$ sur cet espace, $G \times_k P \to P$,
est donnée par $x_i \to t_i \otimes x_i$ au niveau des algèbres.
Nous omettons la démonstration.
En fait, pour l'exemple, il nous suffit que $P$ soit un torseur sous $G$,
ce qui est clair puisque, sur $k' = k(\sqrt{a_1}, \sqrt{a_2}, \ldots)$,
le schéma $P$ devient isomorphe à $G$ de manière $G$-équivariante.
Notons que $P$ est trivial si et seulement si $k' = k$ : en effet, si
$P$ possède un point rationnel sur $k$, tous les $a_i$ sont des carrés.

\medskip\noindent
Nous affirmons que $P$ est un torseur fppf si et seulement si l'extension de corps
$k' = k(\sqrt{a_1}, \sqrt{a_2}, \ldots)/k$ est finie.
Si $k'$ est une extension finie de $k$, alors
$\{\Spec(k') \to \Spec(k)\}$
est un recouvrement fppf qui trivialise $P$, et nous voyons que $P$ est bien
un torseur fppf. Réciproquement, supposons que $P$ soit un torseur fppf sous $G$.
Cela signifie qu'il existe un recouvrement fppf
$\{S_i \to \Spec(k)\}$ tel que chaque $P_{S_i}$ soit trivial.
Choisissons un $i$ tel que $S_i$ ne soit pas vide. Soit $s \in S_i$ un point
fermé. D'après
Variétés, lemme \ref{varieties-lemma-locally-finite-type-Jacobson},
l'extension de corps $\kappa(s)/k$ est finie et, par construction,
$P_{\kappa(s)}$ possède un point rationnel sur $\kappa(s)$. Nous voyons donc que
$k \subset k' \subset \kappa(s)$ et $k'$ est une extension finie de $k$.

\medskip\noindent
Pour obtenir un exemple explicite, prenons $k = \mathbf{Q}$ et $a_i = i$,
par exemple (ou bien $a_i$ égal au $i$-ième nombre premier, si l'on préfère).

\begin{lemma}
\label{lemma-torsors-principal-spaces-not-equal}
Soient $S$ un schéma et $G$ un schéma en groupes sur $S$.
Le champ $G\textit{-Principal}$ qui classifie les espaces principaux homogènes sous $G$
(voir Exemples de champs, sous-section
\ref{examples-stacks-subsection-principal-homogeneous-spaces})
et le champ $G\textit{-Torsors}$ qui classifie les torseurs fppf sous $G$
(voir Exemples de champs, sous-section
\ref{examples-stacks-subsection-fppf-torsors})
ne sont pas équivalents en général.
\end{lemma}

\begin{proof}
La discussion ci-dessus montre que le foncteur
$G\textit{-Torsors} \to G\textit{-Principal}$ n'est pas essentiellement
surjectif en général.
\end{proof}









\section[Champ non algébrique à recouvrement plat quasi compact]{Champ non algébrique admettant un recouvrement plat quasi compact}
\label{section-not-algebraic-stack}

\noindent
Dans cette section, nous décrivons brièvement un exemple dû à Brian Conrad.
On peut trouver cet exemple en ligne à
\href{https://mathoverflow.net/questions/15082/fpqc-covers-of-stacks/15269#15269}{cette adresse}.
Notre exemple est légèrement différent.

\medskip\noindent
Soit $k$ un corps algébriquement clos.
Dans la suite, tous les schémas et champs sont sur $k$.
Soit $G \to \Spec(k)$ un schéma en groupes affine.
Dans Exemples de champs,
lemme \ref{examples-stacks-lemma-classifying-stacks},
nous avons donné plusieurs manières équivalentes de voir
$\mathcal{X} = [\Spec(k)/G]$ comme un champ en groupoïdes sur
$(\Sch/\Spec(k))_{fppf}$. En particulier, $\mathcal{X}$ classifie les
torseurs fppf sous $G$. Plus précisément, un $1$-morphisme $T \to \mathcal{X}$
correspond à un torseur fppf sous $G_T$, noté $P$, sur $T$, et les $2$-flèches
correspondent aux isomorphismes de torseurs. Il s'ensuit que
le $1$-morphisme diagonal
$$
\Delta :
\mathcal{X}
\longrightarrow
\mathcal{X} \times_{\Spec(k)} \mathcal{X}
$$
est représentable et affine. En effet, pour toute paire
de torseurs fppf sous $G_T$, notés $P_1, P_2$, sur un schéma $T/k$, le
schéma $\mathit{Isom}(P_1, P_2)$ est affine sur $T$.
Le torseur trivial sous $G$ sur $\Spec(k)$ définit un $1$-morphisme
$$
f : \Spec(k) \longrightarrow \mathcal{X}.
$$
Nous affirmons qu'il s'agit d'un $1$-morphisme surjectif. En effet,
par définition, pour tout $1$-morphisme $T \to \mathcal{X}$, il existe
un recouvrement fppf $\{T_i \to T\}$ tel que $P_{T_i}$ soit isomorphe
au torseur trivial sous $G_{T_i}$. Les composés
$T_i \to T \to \mathcal{X}$ se factorisent donc par $f$. Il est alors clair que
la projection $T \times_\mathcal{X} \Spec(k) \to T$
est surjective (c'est ainsi que nous définissons la surjectivité de $f$; voir
Champs algébriques,
définition \ref{algebraic-definition-relative-representable-property}).
On montre de même que $f$ est quasi compact et plat (nous omettons les détails).
Notons également que
$$
\Spec(k) \times_\mathcal{X} \Spec(k) \cong G
$$
comme schémas sur $k$.

\medskip\noindent
Supposons qu'il existe un morphisme lisse surjectif
$p : U \to \mathcal{X}$, où $U$ est un schéma.
Considérons le produit fibré
$$
\xymatrix{
W \ar[d] \ar[r] & U \ar[d] \\
\Spec(k) \ar[r] & \mathcal{X}
}
$$
Alors $W$ est un schéma lisse non vide sur $k$ et possède donc
un point rationnel sur $k$. Ainsi, $f$ se factorise par $U$.
Nous obtenons donc
$$
G \cong
\Spec(k) \times_\mathcal{X} \Spec(k) \cong
(\Spec(k) \times_k \Spec(k))
\times_{(U \times_k U)}
(U \times_\mathcal{X} U)
$$
et, puisque les projections $U \times_\mathcal{X} U \to U$ sont
supposées lisses, nous en concluons que
$U \times_\mathcal{X} U \to U \times_k U$ est
localement de type fini; voir
Morphismes, lemme \ref{morphisms-lemma-permanence-finite-type}.
Il s'ensuit que, dans ce cas, $G$ est localement de type fini sur $k$.
Nous avons ainsi démontré le lemme suivant (qui admet une
généralisation considérable).

\begin{lemma}
\label{lemma-BG-algebraic}
Soit $k$ un corps. Soit $G$ un schéma en groupes affine sur $k$.
Si le champ $[\Spec(k)/G]$ possède un recouvrement lisse par un
schéma, alors $G$ est de type fini sur $k$.
\end{lemma}

\begin{proof}
Voir la discussion ci-dessus.
\end{proof}

\noindent
Pour obtenir un exemple explicite comme dans le titre de cette section, prenons
$G = (\mu_{2, k})^\infty$, le schéma en groupes de la
section \ref{section-torsor-not-fppf}, qui n'est pas localement de type fini
sur $k$. La discussion ci-dessus montre que
$\mathcal{X} = [\Spec(k)/G]$ possède les propriétés (1) et (2) de
Champs algébriques, définition \ref{algebraic-definition-algebraic-stack},
mais non la propriété (3). Ainsi, $\mathcal{X}$ n'est pas un champ algébrique.
En revanche, il existe bien un schéma $U$ et un morphisme surjectif,
plat et quasi compact $U \to \mathcal{X}$, à savoir le morphisme
$f : \Spec(k) \to \mathcal{X}$ étudié ci-dessus.






\section{Commutation aux limites sur les objets, mais non aux limites}
\label{section-limit-preserving}

\noindent
Soit $S$ un schéma non vide. Soit $\mathcal{G}$ un faisceau abélien injectif
sur $(\Sch/S)_{fppf}$. Nous obtenons un champ en groupoïdes
$$
\mathcal{G}\textit{-Torsors} \longrightarrow (\Sch/S)_{fppf}
$$
sur $S$; voir Exemples de champs, lemme
\ref{examples-stacks-lemma-torsors-sheaf-stack-in-groupoids}.
Ce champ commute aux limites sur les objets au-dessus de $(\Sch/S)_{fppf}$ (voir
Critères de représentabilité, section \ref{criteria-section-limit-preserving}),
car tout torseur sous $\mathcal{G}$ est trivial. En revanche,
$\mathcal{G}\textit{-Torsors}$ ne commute en général pas aux limites (voir
Axiomes d'Artin, définition \ref{artin-definition-limit-preserving}),
car le faisceau $\mathcal{G}$ ne commute pas nécessairement aux limites. Par exemple,
prenons un faisceau injectif non nul quelconque $\mathcal{I}$ et posons
$\mathcal{G} = \prod_{n \in \mathbf{Z}} \mathcal{I}$ pour obtenir un
exemple.

\begin{lemma}
\label{lemma-limit-preserving-on-objects-not-limit-preserving}
Soit $S$ un schéma non vide. Il existe un champ en groupoïdes
$p : \mathcal{X} \to (\Sch/S)_{fppf}$
tel que $p$ commute aux limites sur les objets, mais que $\mathcal{X}$ ne
commute pas aux limites.
\end{lemma}

\begin{proof}
Voir la discussion ci-dessus.
\end{proof}



\section{Un champ classifiant non algébrique}
\label{section-non-algebraic}

\noindent
Soit $S = \Spec(\mathbf{F}_p)$ et notons $\mu_p$ le
schéma en groupes des racines $p$-ièmes de l'unité sur $S$. Dans
Groupoïdes dans les espaces, section \ref{spaces-groupoids-section-stacks},
nous avons introduit le champ quotient $[S/\mu_p]$ et, dans
Exemples de champs, section \ref{examples-stacks-section-group-quotient-stacks},
nous avons montré que $[S/\mu_p]$ est le champ classifiant les
torseurs fppf sous $\mu_p$ : pour tout schéma $T$ sur $S$, la catégorie
$\Mor_S(T, [S/\mu_p])$ est canoniquement équivalente à la
catégorie des torseurs fppf sous $\mu_p$ sur $T$. Enfin, dans
Critères de représentabilité, théorème
\ref{criteria-theorem-flat-groupoid-gives-algebraic-stack}
nous avons vu que $[S/\mu_p]$ est un champ algébrique.

\medskip\noindent
Nous pouvons maintenant poser la question : « Qu'en est-il de la catégorie fibrée en groupoïdes
$\mathcal{S}$ qui classifie les torseurs étales sous $\mu_p$ ? » (Autrement dit,
$\mathcal{S}$ est une catégorie au-dessus de $\Sch/S$ dont la catégorie fibre au-dessus d'un
schéma $T$ est la catégorie des torseurs étales sous $\mu_p$ sur $T$.)

\medskip\noindent
La première objection est qu'il ne s'agit pas d'un champ pour la topologie fppf,
car la descente des objets n'y est pas satisfaite. Par exemple, le
torseur sous $\mu_p$ donné par $\Spec(\mathbf{F}_p(t)[x]/(x^p - t))$ sur
$T = \Spec(\mathbf{F}_p(t))$ est localement trivial pour la topologie fppf, mais
non pour la topologie étale.

\medskip\noindent
Pour remédier à ce premier problème, travaillons avec la topologie étale :
la descente des objets y est satisfaite. En effet,
$\mathcal{S}$ est un champ en groupoïdes sur $(\Sch/S)_\etale$.
De plus, la diagonale
$\Delta : \mathcal{S} \to \mathcal{S} \times \mathcal{S}$
est représentable (par des schémas). En effet, étant donnés deux
torseurs sous $\mu_p$ (qu'ils soient ou non localement triviaux pour la topologie étale),
le faisceau de leurs isomorphismes est représentable par un schéma.

\medskip\noindent
Nous pouvons donc enfin demander s'il existe un schéma $U$ et un
$1$-morphisme lisse et surjectif $U \to \mathcal{S}$. Nous montrerons de deux
manières que cela est impossible : par un argument direct (que nous conseillons
au lecteur de passer) et par un argument utilisant un résultat général.

\medskip\noindent
Argument direct (esquisse) : notons que le $1$-morphisme
$\mathcal{S} \to \Spec(\mathbf{F}_p)$ satisfait au
critère de relèvement infinitésimal de lissité formelle.
En effet, pour tout épaississement infinitésimal du premier ordre
de schémas $T \to T'$, le noyau de $\mu_p(T') \to \mu_p(T)$
est isomorphe aux sections du faisceau d'idéaux de $T$ dans $T'$, et
ainsi $H^1_\etale(T, \mu_p) = H^1_\etale(T', \mu_p)$.
De plus, le champ $\mathcal{S}$ commute aux limites. Par conséquent,
si $U \to \mathcal{S}$ est lisse, alors $U \to \Spec(\mathbf{F}_p)$
commute aux limites et satisfait au
critère de relèvement infinitésimal de lissité formelle.
Il en résulte que $U$ est lisse sur $\mathbf{F}_p$.
En particulier, $U$ est réduit, donc $H^1_\etale(U, \mu_p) = 0$.
Ainsi, $U \to \mathcal{S}$ se factorise en
$U \to \Spec(\mathbf{F}_p) \to \mathcal{S}$
et la première flèche est lisse. Par descente de la lissité,
la lissité de $U \to \mathcal{S}$ impliquerait celle de
$\Spec(\mathbf{F}_p) \to \mathcal{S}$. Or ce n'est
pas le cas, puisque
$\Spec(\mathbf{F}_p) \times_\mathcal{S} \Spec(\mathbf{F}_p)$
est $\mu_p$, qui n'est pas lisse sur $\Spec(\mathbf{F}_p)$.


\medskip\noindent
Argument structurel : dans
Critères de représentabilité, section \ref{criteria-section-stacks-etale},
nous avons vu que l'on peut considérer les champs algébriques comme les
champs en groupoïdes pour la topologie étale dont la diagonale est
représentable par des espaces algébriques admettant un recouvrement lisse.
Ainsi, s'il existe un morphisme lisse surjectif $U \to \mathcal{S}$, alors
$\mathcal{S}$ est un champ algébrique et satisfait en particulier à la
descente pour la topologie fppf. Mais nous avons vu ci-dessus que $\mathcal{S}$
ne satisfait pas à la descente pour la topologie fppf.

\medskip\noindent
De manière informelle, les arguments ci-dessus montrent que le champ classifiant,
pour la topologie étale, les torseurs localement triviaux pour cette topologie
sous un schéma en groupes $G$ sur une base $B$ est algébrique si et seulement
si $G$ est lisse sur $B$. L'un des avantages de travailler avec
la topologie fppf est qu'il suffit de supposer que $G \to B$
est plat et localement de présentation finie. En fait, le champ quotient
(pour la topologie fppf) $[B/G]$ est algébrique si et seulement
si $G \to B$ est plat et localement de présentation finie; voir
Critères de représentabilité, lemme \ref{criteria-lemma-BG-algebraic}.







\section{Faisceau muni d'un recouvrement plat quasi-compact qui n'est pas algébrique}
\label{section-not-algebraic}

\noindent
Considérons le foncteur $F = (\mathbf{P}^1)^\infty$, c'est-à-dire que, pour un schéma $T$,
$F(T)$ est l'ensemble des $f = (f_1, f_2, f_3, \ldots)$ où chaque
$f_i : T \to \mathbf{P}^1$ est un morphisme de schémas. Remarquons que
$\mathbf{P}^1$ vérifie la propriété de faisceau pour les recouvrements fpqc; voir
Descente, lemme \ref{descent-lemma-fpqc-universal-effective-epimorphisms}.
Un produit de faisceaux étant un faisceau, $F$ vérifie aussi la propriété de faisceau pour
la topologie fpqc. La diagonale de $F$ est représentable : si $f : T \to F$
et $g : S \to F$ sont des morphismes, alors $T \times_F S$ est l'intersection schématique
des sous-schémas fermés $T \times_{f_i, \mathbf{P}^1, g_i} S$
dans le schéma $T \times S$. Considérons le schéma en groupes $\text{SL}_2$, muni
d'un morphisme affine, lisse et surjectif $\text{SL}_2 \to \mathbf{P}^1$.
Considérons ensuite $U = (\text{SL}_2)^\infty$ et son morphisme canonique (produit)
$U \to F$. Remarquons que $U$ est un schéma affine. Nous affirmons que le morphisme
$U \to F$ est plat, surjectif et universellement ouvert. En effet, supposons donné
un morphisme $f : T \to F$. Alors $Z = T \times_F U$ est le produit fibré infini
des schémas $Z_i = T \times_{f_i, \mathbf{P}^1} \text{SL}_2$
sur $T$. Chacun des morphismes $Z_i \to T$ est lisse, surjectif et
affine, ce qui entraîne que
$$
Z = Z_1 \times_T Z_2 \times_T Z_3 \times_T \ldots
$$
est un schéma plat et affine sur $T$. Un argument élémentaire de limite montre aussi que
$Z \to T$ est ouvert.

\medskip\noindent
En revanche, nous affirmons que $F$ n'est pas un espace algébrique.
En effet, si $F$ était un espace algébrique, ce serait un espace algébrique
quasi-compact et séparé (d'après notre description des produits fibrés sur $F$).
La cohomologie des faisceaux quasi-cohérents s'annulerait donc au-dessus d'un
certain degré (voir
Cohomologie des espaces, proposition \ref{spaces-cohomology-proposition-vanishing},
et les remarques qui la précèdent). Mais il est clair qu'en prenant l'image réciproque
de $\mathcal{O}(-2, -2, \ldots, -2)$ par la projection
$$
(\mathbf{P}^1)^\infty \longrightarrow (\mathbf{P}^1)^n
$$
(qui possède une section), on peut obtenir un faisceau quasi-cohérent dont la cohomologie
est non nulle en degré $n$. Nous obtenons ainsi une réponse à une question posée
par Anton Geraschenko sur MathOverflow.

\begin{lemma}
\label{lemma-not-algebraic}
Il existe un foncteur $F : \Sch^{opp} \to \textit{Ensembles}$
qui satisfait à la condition de faisceau pour la topologie fpqc, dont la diagonale
$\Delta : F \to F \times F$ est représentable, et tel qu'il existe un morphisme
surjectif, plat, universellement ouvert et quasi-compact
$U \to F$, où $U$ est un schéma, mais tel que $F$ ne soit pas un espace algébrique.
\end{lemma}

\begin{proof}
Voir la discussion qui précède.
\end{proof}







\section[Topologie \'etale, Zariski et recouvrements finis \'etales]{La topologie \'etale face à celle de Zariski et aux recouvrements finis et \'etales}
\label{section-etale-zariski-finite-etale}

\noindent
Nous donnons dans cette section un exemple trouvé par B.~Bhatt en réponse
à une question de C.~Deninger, qui montre que la topologie \'etale
des variétés algébriques sur $\mathbf{C}$ est strictement plus fine que la
topologie engendrée par les immersions ouvertes et les morphismes finis et \'etales.
De tels exemples sont certainement connus des spécialistes et remontent peut-être
à Artin et Grothendieck\footnote{Selon une anecdote, Grothendieck aurait d'abord
tenté de définir la topologie \'etale au moyen des recouvrements ouverts de Zariski et des
recouvrements finis et \'etales, avant qu'Artin ne le convainque du contraire
(probablement à partir d'exemples comme celui consigné ici).}.
Si vous connaissez un exemple publié, veuillez écrire à
\href{mailto:stacks.project@gmail.com}{stacks.project@gmail.com}.

\medskip\noindent
Soit $X$ une courbe affine lisse et connexe sur $\mathbf{C}$.
Choisissons un morphisme fini $g: Y \to X$ de degré $3$, où
$Y$ est une courbe lisse connexe, et un point $x \in X$
dont la fibre $g^{-1}(x) = \{u, y\}$ est formée d'un point non ramifié
$u$ (de degré 1) et d'un point ramifié $y$ (de degré 2). En remplaçant
$X$ par un voisinage ouvert de Zariski de $x$, nous pouvons supposer que $g$
est non ramifié partout sauf en $y$. Soit $U = Y \setminus \{y\}$.
Alors $f=g|_U: U \to X$ est un morphisme \'etale surjectif et
$f^{-1}(x) = \{u\}$. 

\begin{lemma}
\label{lemma-etale-not-dominated-zariski-finite-etale}
Le recouvrement \'etale $U \to X$ n'admet aucun raffinement par une composition finie
d'immersions ouvertes et de morphismes finis et \'etales.
\end{lemma}

\begin{proof}
Raisonnons par l'absurde et supposons donnée une tour de morphismes
$$
W_m \to W_{m-1} \to \dots \to W_1 \to W_0 = X
$$
où chaque morphisme $W_i \to W_{i - 1}$ est soit une immersion ouverte, soit un
morphisme fini et \'etale, ainsi qu'un relèvement $W_m \to U$ de $W_m \to X$ le long de
$f : U \to X$ et un point $w \in W_m$ d'image $u \in U$
(et donc $x \in X$). En remplaçant chaque $W_k$ par la composante connexe
qui contient l'image $w_k \in W_k$ de $w$, nous pouvons supposer
que chaque $W_k$ est lisse et connexe.

\medskip\noindent
Pour chaque $0 \le i \le m$, considérons le produit fibré $Z_i = Y \times_X W_i$.
Comme $W_i \to X$ est \'etale et $Y$ est lisse, la courbe $Z_i$ est une réunion
disjointe de courbes lisses connexes, et les morphismes $Z_i \to W_i$ sont finis
de degré $3$ par changement de base. La fibre de $Z_i \to W_i$
au-dessus de $w_i$ est manifestement formée de deux points $z_i, v_i$ qui s'envoient
respectivement sur $y, u$ dans $Y$. Le morphisme $Z_i \to W_i$ est non ramifié en $v_i$
et ramifié de degré $2$ en $z_i$. Posons $n_i = \#\pi_0(Z_i)$. Comme la fibre de
$Z_i \to W_i$ possède $2$ points, nous avons $n_i \in \{1, 2\}$. De plus, $n_i$ ne peut
que croître avec $i$, car les morphismes de transition $Z_i \to Z_{i - 1}$ sont dominants.
Nous avons clairement $n_0 = 1$. D'autre part, $n_m \geq 2$ : il existe un
$X$-morphisme $W_m \to U \subset Y$, donc le changement de base $Z_m \to W_m$
possède une section. Il existe alors un plus petit indice $i \ge 1$ tel que $n_{i - 1} = 1$
et $n_i = 2$. Comme les $n_j$ sont déterminés par l'irréductibilité de $Z_j$,
le morphisme $W_i \to W_{i - 1}$ ne peut être une immersion ouverte;
il doit donc être fini et \'etale. Puisque $n_i = 2$, la courbe lisse $Z_i$
possède deux composantes connexes; puisque le morphisme vers $W_i$ est fini, l'une de
ces composantes s'envoie nécessairement isomorphiquement sur $W_i$.
Nous pouvons donc écrire $Z_i = W_i \sqcup Z'_i$, où $Z'_i \to W_i$ est de degré $2$.
Le morphisme induit $W_i \subset Z_i \to Z_{i - 1}$ est dominant
entre deux courbes connexes, toutes deux finies sur $W_{i - 1}$.
Un tel morphisme est nécessairement surjectif. Nous pouvons donc choisir un point $t \in W_i$
d'image $z_{i - 1}$. Nous obtenons des extensions
$$
\mathcal{O}_{W_{i - 1}, w_{i - 1}} \to
\mathcal{O}_{Z_{i - 1}, z_{i - 1}} \to
\mathcal{O}_{W_i, t}
$$
d'anneaux de valuation discrète. La première est ramifiée, tandis que
l'extension composée est non ramifiée puisque $W_i \to W_{i - 1}$ est \'etale; c'est
impossible, d'où la contradiction.
\end{proof}










\section{Faisceaux et spécialisations}
\label{section-sheaves}

\noindent
Dans la suite, nous fixons un grand site \'etale $\Sch_\etale$
construit dans
Topologies, définition \ref{topologies-definition-big-etale-site}.
De plus, par « schéma », nous entendrons un objet de ce site.
Rappelons que, si $x, x'$ sont des points d'un schéma $X$, on dit que $x$ est une
{\it spécialisation} de $x'$, ou l'on écrit $x' \leadsto x$, si
$x \in \overline{\{x'\}}$. C'est notamment le cas lorsque $x = x'$.

\medskip\noindent
Considérons le foncteur $F : \Sch_\etale \to \textit{Ab}$
défini par les règles suivantes :
$$
F(X) = \prod\nolimits_{x \in X}
\prod\nolimits_{x' \in X, x' \leadsto x, x' \not = x}
\mathbf{Z}/2\mathbf{Z}
$$
Pour tout schéma $X$, nous notons $|X|$ son ensemble sous-jacent.
Un élément $a \in F(X)$ sera considéré comme une application ensembliste
$|X| \times |X| \to \mathbf{Z}/2\mathbf{Z}$, $(x, x') \mapsto a(x, x')$,
nulle si $x = x'$ ou si $x$ n'est pas une spécialisation de $x'$.
Pour un morphisme de schémas $f : X \to Y$, nous définissons
$$
F(f) : F(Y) \longrightarrow F(X)
$$
par la règle suivante : pour $b \in F(Y)$, nous posons
$$
F(f)(b)(x, x') =
\left\{
\begin{matrix}
0 & \text{si }x\text{ n'est pas une spécialisation de }x' \\
b(f(x), f(x')) & \text{sinon.}
\end{matrix}
\right.
$$
Remarquons que cela définit bien un élément de $F(X)$. Nous affirmons que, si
$f : X \to Y$ et $g : Y \to Z$ sont des morphismes composables, alors
$F(f) \circ F(g) = F(g \circ f)$. En effet, soient $c \in F(Z)$ et
$x' \leadsto x$ une spécialisation entre deux points de $X$; alors
$$
F(g \circ f)(c)(x, x') = c(g(f(x)), g(f(x'))) = F(f)(F(g)(c))(x, x')
$$
car $f(x') \leadsto f(x)$. (Cela vaut aussi si $f(x) = f(x')$.)

\medskip\noindent
Soit $G$ le faisceau associé à $F$ pour la topologie \'etale.

\medskip\noindent
Nous affirmons que, si $X$ est un schéma et $x' \leadsto x$ une spécialisation
avec $x' \not = x$, alors $G(X) \not = 0$. En effet, soit $a \in F(X)$
un élément tel que $a$, considéré comme la fonction
$|X| \times |X| \to \mathbf{Z}/2\mathbf{Z}$, est non nul en $(x, x')$.
Soit $\{f_i : U_i \to X\}$ un recouvrement \'etale de $X$. Nous pouvons choisir un
indice $i$ et un point $u_i \in U_i$ tels que $f_i(u_i) = x$. Comme les générisations
se relèvent le long des morphismes plats (voir
Morphismes, lemme \ref{morphisms-lemma-generalizations-lift-flat}),
nous pouvons trouver une spécialisation $u'_i \leadsto u_i$
telle que $f_i(u'_i) = x'$. Notre construction montre alors que
$F(f_i)(a) \not = 0$. Ainsi, $a$ détermine un élément non nul de $G(X)$.

\medskip\noindent
Remarquons que, si $X = \Spec(k)$, où $k$ est un corps (ou, plus généralement,
un anneau dont tous les idéaux premiers sont maximaux), alors $F(X) = 0$;
de plus, pour tout morphisme \'etale $U \to X$, nous avons $F(U) = 0$, car il
n'existe aucune spécialisation entre des points distincts dans les fibres d'un
morphisme \'etale. Par conséquent, $G(X) = 0$.

\medskip\noindent
Supposons que $X \subset X'$ soit un épaississement; voir
Compléments sur les morphismes, définition \ref{more-morphisms-definition-thickening}.
Alors la catégorie des schémas \'etales sur $X'$ est équivalente à la
catégorie des schémas \'etales sur $X$ par le foncteur de changement de base
$U' \mapsto U = U' \times_{X'} X$; voir
Cohomologie \'etale,
théorème \ref{etale-cohomology-theorem-topological-invariance}.
Puisque, dans cette situation, on a toujours $F(U) = F(U')$,
nous avons aussi $G(X) = G(X')$.

\medskip\noindent
Comme variante, nous pouvons considérer le préfaisceau $F_n$ qui associe
à un schéma $X$ l'ensemble des applications
$a : |X|^{n + 1} \to \mathbf{Z}/2\mathbf{Z}$ telles que $a(x_0, \ldots, x_n)$
ne soit non nul que si $x_n \leadsto \ldots \leadsto x_0$ est une suite de
spécialisations et $x_n \not = x_{n - 1} \not = \ldots \not = x_0$.
Soit $G_n$ le faisceau associé à $F_n$.
Exactement comme ci-dessus, on montre que $G_n$ est non nul
si $\dim(X) \geq n$ et nul si $\dim(X) < n$.

\begin{lemma}
\label{lemma-sheaf-zero-on-low-dimension}
Il existe un faisceau de groupes abéliens $G$ sur
$\Sch_\etale$ qui possède les propriétés suivantes :
\begin{enumerate}
\item $G(X) = 0$ dès que $\dim(X) < n$,
\item $G(X)$ n'est pas nul si $\dim(X) \geq n$, et
\item si $X \subset X'$ est un épaississement, alors $G(X) = G(X')$.
\end{enumerate}
\end{lemma}

\begin{proof}
Voir la discussion qui précède.
\end{proof}

\begin{remark}
\label{remark-specialization}
Voici quelques remarques :
\begin{enumerate}
\item Les préfaisceaux $F$ et $F_n$ sont séparés.
\item Il se trouve que $F$ et $F_n$ ne sont pas des faisceaux.
\item On peut montrer que $G$ et $G_n$ sont en fait des faisceaux pour la topologie fppf.
\end{enumerate}
Nous démontrerons ces résultats si nous en avons besoin.
\end{remark}


\section{Faisceaux et fonctions constructibles}
\label{section-constructible-functions}

\noindent
Dans la suite, nous fixons un grand site \'etale $\Sch_\etale$ construit dans
Topologies, définition \ref{topologies-definition-big-etale-site}.
De plus, par « schéma », nous entendrons un objet de ce site.
Dans cette section, nous appelons {\it partition constructible}
d'un schéma $X$ toute décomposition localement finie en réunion disjointe
$X = \coprod_{i \in I} X_i$ telle que chaque $X_i \subset X$ soit une
partie localement constructible de $X$. Localement finie signifie que,
pour tout ouvert quasi-compact $U \subset X$, il n'existe qu'un nombre fini
d'indices $i \in I$ tels que $X_i \cap U$ soit non vide.
Remarquons que, si $f : X \to Y$ est un morphisme de schémas et
$Y = \coprod Y_j$ une partition constructible, alors
$X = \coprod f^{-1}(Y_j)$ est une partition constructible de $X$.
Étant donnés un ensemble $S$ et un schéma $X$, une {\it fonction constructible}
$f : |X| \to S$ est une application telle que $X = \coprod_{s \in S} f^{-1}(s)$
soit une partition constructible de $X$.
Si $G$ est un groupe abstrait et $a, b : |X| \to G$ des fonctions
constructibles, alors $ab : |X| \to G, x \mapsto a(x)b(x)$ est encore une fonction
constructible. En effet, deux partitions constructibles quelconques
admettent une troisième partition qui les raffine toutes deux.

\medskip\noindent
Soit $A$ un groupe abélien quelconque. Pour tout schéma $X$, nous définissons
$$
F(X) =
\frac{\{a : |X| \to A \mid a \text{ est une fonction constructible}\}}{\{\text{fonctions localement constantes }|X| \to A\}}
$$
Nous considérons simplement un élément $a$ de $F(X)$ comme une fonction définie
à une fonction localement constante près. Étant donnés un morphisme de schémas
$f : X \to Y$ et un élément $b \in F(Y)$, nous posons
$F(f)(b) = b \circ f$. Ainsi, $F$ est un préfaisceau sur $\Sch_\etale$.

\medskip\noindent
Remarquons que, si $\{f_i : U_i \to X\}$ est un recouvrement fppf et si $a \in F(X)$
est tel que $F(f_i)(a) = 0$ dans $F(U_i)$, alors $a \circ f_i$ est une fonction
localement constante pour tout $i$. Il en résulte que $a$ est elle-même localement
constante, puisque les morphismes $f_i$ sont ouverts. Donc $a = 0$ dans $F(X)$.
Ainsi, $F$ est un préfaisceau séparé (pour la topologie fppf,
donc a fortiori pour la topologie \'etale).

\medskip\noindent
Soit $G$ le faisceau associé à $F$ pour la topologie \'etale.
Comme $F$ est séparé et que, par exemple, $F(X) \not = 0$ lorsque
$X$ est le spectre d'un anneau de valuation discrète, nous voyons que $G$
n'est pas nul.

\medskip\noindent
Soit $X = \Spec(k)$, où $k$ est un corps.
Alors tout recouvrement \'etale de $X$ admet un raffinement
de la forme $\{\Spec(k') \to \Spec(k)\}$, où $k'/k$
est une extension finie séparable. Puisque $F(\Spec(k')) = 0$,
nous avons $G(X) = 0$.

\medskip\noindent
Supposons que $X \subset X'$ soit un épaississement; voir
Compléments sur les morphismes, définition \ref{more-morphisms-definition-thickening}.
Alors la catégorie des schémas \'etales sur $X'$ est équivalente à la
catégorie des schémas \'etales sur $X$ par le foncteur de changement de base
$U' \mapsto U = U' \times_{X'} X$; voir
Cohomologie \'etale,
théorème \ref{etale-cohomology-theorem-topological-invariance}.
Comme $F(U) = F(U')$ dans cette situation, nous avons aussi $G(X) = G(X')$.

\medskip\noindent
Le faisceau $G$ commute aux limites; voir
Limites d'espaces,
définition \ref{spaces-limits-definition-locally-finite-presentation}.
En effet, soit $R$ un anneau écrit comme une limite inductive filtrante
$R = \colim_i R_i$ d'anneaux. Posons $X = \Spec(R)$ et
$X_i = \Spec(R_i)$, de sorte que $X = \lim_i X_i$. Alors
$G(X) = \colim_i G(X_i)$. Pour le démontrer, on établit d'abord que
toute partition constructible de $\Spec(R)$ provient d'une partition constructible
d'un certain $\Spec(R_i)$. On obtient ainsi le résultat pour $F$. Pour
obtenir le résultat pour le faisceau associé, on utilise que tout homomorphisme \'etale d'anneaux
$R \to R'$ provient d'un homomorphisme \'etale $R_i \to R_i'$ pour un certain $i$.
Nous omettons les détails.

\begin{lemma}
\label{lemma-weird-sheaf}
Il existe un faisceau de groupes abéliens $G$ sur
$\Sch_\etale$ qui possède les propriétés suivantes :
\begin{enumerate}
\item $G(\Spec(k)) = 0$ pour tout corps $k$,
\item $G$ commute aux limites,
\item si $X \subset X'$ est un épaississement, alors $G(X) = G(X')$, et
\item $G$ n'est pas nul.
\end{enumerate}
\end{lemma}

\begin{proof}
Voir la discussion qui précède.
\end{proof}




\section{Le site lisse-\'etale n'est pas fonctoriel}
\label{section-lisse-etale-not-functorial}

\noindent
Le site {\it lisse-\'etale}
$X_{lisse,\etale}$ de $X$ est la catégorie des schémas lisses sur $X$
munie des recouvrements \'etales (usuels); voir
Cohomologie des champs, section \ref{stacks-cohomology-section-lisse-etale}.
Soit $f : X  \to Y$ un morphisme de schémas.
Il existe un foncteur
$$
u : Y_{lisse,\etale} \longrightarrow X_{lisse,\etale},\quad
V/Y \longmapsto V \times_Y X
$$
qui est continu. Nous obtenons donc une paire de foncteurs adjoints
$$
u^s :
\Sh(X_{lisse,\etale})
\longrightarrow
\Sh(Y_{lisse,\etale}),
\quad
u_s :
\Sh(Y_{lisse,\etale})
\longrightarrow
\Sh(X_{lisse,\etale}),
$$
voir Sites, section \ref{sites-section-continuous-functors}.
Nous affirmons qu'en général $u$ ne définit {\bf pas} un morphisme de sites; voir
Sites, définition \ref{sites-definition-morphism-sites}.
Autrement dit, nous affirmons que $u_s$ n'est pas exact à gauche en général. Notons que
les préfaisceaux représentables sont des faisceaux sur les sites lisses-\'etales. Par
Sites, lemme \ref{sites-lemma-pullback-representable-sheaf},
nous avons donc $u_sh_V = h_{V \times_Y X}$. Considérons deux morphismes
$$
\xymatrix{
V_1 \ar[rd] \ar@<1ex>[rr]^a \ar@<-1ex>[rr]_b & & V_2 \ar[ld] \\
& Y
}
$$
de schémas $V_1, V_2$ lisses sur $Y$. Si $u_s$ était exact à gauche, nous
aurions alors
$$
u_s \text{Égaliseur}(h_a, h_b : h_{V_1} \to h_{V_2})
=
\text{Égaliseur}(h_{a \times 1}, h_{b \times 1} :
h_{V_1 \times_Y X} \to h_{V_2 \times_Y X})
$$
Nous choisirons les morphismes $a, b : V_1 \to V_2$ de sorte qu'il n'existe
aucun morphisme d'un schéma lisse sur $Y$ vers $(a = b) \subset V_1$, c'est-à-dire
que le membre de gauche soit le faisceau vide, mais qu'après
changement de base à $X$ l'égalisateur soit non vide et lisse sur $X$.
Un exemple rudimentaire consiste à prendre $X = \Spec(\mathbf{F}_p)$,
$Y = \Spec(\mathbf{Z})$ et $V_1 = V_2 = \mathbf{A}^1_\mathbf{Z}$,
avec $a(x) = x$ et $b(x) = x + p$. Remarquons que l'égalisateur
de $a$ et $b$ est la fibre de $\mathbf{A}^1_\mathbf{Z}$ au-dessus de $(p)$.

\begin{lemma}
\label{lemma-lisse-etale-not-functorial}
Le site lisse-\'etale n'est pas fonctoriel, même pour les morphismes de schémas.
\end{lemma}

\begin{proof}
Voir la discussion qui précède.
\end{proof}






\section{Faisceaux sur la catégorie des schémas noethériens}
\label{section-sheaves-locally-Noetherian}

\noindent
Soit $S$ un schéma localement noethérien. Comme dans
Axiomes d'Artin, section \ref{artin-section-noetherian},
considérons le foncteur d'inclusion
$$
u : (\textit{Noetherian}/S)_{fppf} \longrightarrow (\Sch/S)_{fppf}
$$
du site fppf des schémas localement noethériens sur $S$ dans un
grand site fppf de $S$. Comme expliqué dans la section citée,
ce foncteur est continu. Nous obtenons donc une paire de foncteurs adjoints
$$
u^s :
\Sh((\Sch/S)_{fppf})
\longrightarrow
\Sh((\textit{Noetherian}/S)_{fppf})
$$
et
$$
u_s :
\Sh((\textit{Noetherian}/S)_{fppf})
\longrightarrow
\Sh((\Sch/S)_{fppf})
$$
voir Sites, section \ref{sites-section-continuous-functors}.
Nous affirmons toutefois qu'en général $u$ ne définit pas un
morphisme de sites; voir Sites, définition \ref{sites-definition-morphism-sites}.
Autrement dit, le foncteur $u_s$ n'est pas exact à gauche en général.

\medskip\noindent
Soit $p$ un nombre premier et posons $S = \Spec(\mathbf{F}_p)$.
Considérons le morphisme injectif de faisceaux
$$
a : \mathcal{F} \longrightarrow \mathcal{G}
$$
sur $(\textit{Noetherian}/S)_{fppf}$ défini comme suit : pour $U$
un schéma localement noethérien sur $S$, nous définissons
$$
\mathcal{G}(U) = \Gamma(U, \mathcal{O}_U)^* =
\Mor_S(U, \mathbf{G}_{m, S})
$$
et nous prenons
$$
\begin{aligned}
\mathcal{F}(U) = \{f \in \mathcal{G}(U) \mid {}&
\text{localement pour la topologie fppf, }f
\text{ admet des racines de degré une puissance de }p\\
&\text{arbitrairement grande}\}
\end{aligned}
$$
Une $\mathbf{F}_p$-algèbre noethérienne $A$ possède un nilradical nilpotent
$I \subset A$; les racines de degré une puissance de $p$ de $1$ dans $A$ sont les
éléments de la forme $1 + a$, $a \in I$, et aucune racine non triviale
de degré une puissance de $p$ de $1$ n'admet donc des racines de degré une puissance de $p$ arbitrairement grande. Nous en déduisons
que $\mathcal{F}(U)$ est un groupe abélien sans $p$-torsion pour tout schéma
localement noethérien $U$; nous omettons certains détails. Il s'ensuit que
$p : \mathcal{F} \to \mathcal{F}$
est un morphisme injectif de faisceaux abéliens sur $(\textit{Noetherian}/S)_{fppf}$.

\medskip\noindent
Pour obtenir une contradiction, supposons $u_s$ exact. Alors
$p : u_s\mathcal{F} \to u_s\mathcal{F}$ est lui aussi injectif,
et $(u_s\mathcal{F})(V)$ est un groupe abélien sans $p$-torsion
pour tout $V$ sur $S$. Puisque les
préfaisceaux représentables sont des faisceaux sur les sites fppf, il résulte de
Sites, lemme \ref{sites-lemma-pullback-representable-sheaf},
que $u_s\mathcal{G}$ est représenté par
$\mathbf{G}_{m, S}$. Comme $u_s\mathcal{F} \to u_s\mathcal{G}$
est injectif, nous obtenons un sous-groupe sans $p$-torsion
$$
(u_s\mathcal{F})(V) \subset \Gamma(V, \mathcal{O}_V)^*
$$
pour tout schéma $V$ sur $S$, avec la propriété suivante :
pour tout morphisme $V \to U$ de schémas sur $S$, où $U$ est
localement noethérien, le sous-groupe
$$
\mathcal{F}(U) \subset \Gamma(U, \mathcal{O}_U)^*
$$
s'envoie dans le sous-groupe $(u_s\mathcal{F})(V)$ par l'application de restriction
$\Gamma(U, \mathcal{O}_U)^* \to \Gamma(V, \mathcal{O}_V)^*$.

\medskip\noindent
La contradiction proprement dite s'obtient comme suit :
posons $k = \bigcup_{n \geq 0} \mathbf{F}_p(t^{1/{p^n}})$
et
$$
B = k \otimes_{\mathbf{F}_p(t)} k
$$
et soit $V = \Spec(B)$. Nous avons les deux morphismes de projection
$V \to \Spec(k)$ correspondant aux deux coprojections $k \to B$;
comme $\Spec(k)$ est noethérien, nous concluons que le sous-groupe
$$
(u_s\mathcal{F})(V) \subset B^*
$$
contient $k^* \otimes 1$ et $1 \otimes k^*$. C'est une contradiction,
car
$$
(t^{1/p} \otimes 1) \cdot (1 \otimes t^{-1/p}) =
t^{1/p} \otimes t^{-1/p}
$$
est une unité non triviale de $p$-torsion dans $B$.

\begin{lemma}
\label{lemma-not-a-morphism-of-sites-noetherian-to-all}
Pour $S = \Spec(\mathbf{F}_p)$, le foncteur d'inclusion
$(\textit{Noetherian}/S)_{fppf} \to (\Sch/S)_{fppf}$
ne définit pas un morphisme de sites.
\end{lemma}

\begin{proof}
Voir la discussion qui précède.
\end{proof}





\section{Image directe dérivée des modules quasi-cohérents}
\label{section-derived-push-quasi-coherent}

\noindent
Soit $k$ un corps de caractéristique $p > 0$. Soit $S = \Spec(k[x])$.
Soit $G = \mathbf{Z}/p\mathbf{Z}$, considéré soit comme groupe abstrait,
soit comme schéma en groupes constant sur $S$. Considérons le champ algébrique
$\mathcal{X} = [S/G]$, où $G$ agit trivialement sur $S$; voir
Exemples de champs, remarque \ref{examples-stacks-remark-X-mod-G-group},
et
Critères de représentabilité, lemme \ref{criteria-lemma-BG-algebraic}.
Considérons le morphisme structural
$$
f : \mathcal{X} \longrightarrow S
$$
Ce morphisme est quasi-compact et quasi-séparé. Nous obtenons donc un foncteur
$$
Rf_{\QCoh, *} :
D^{+}_\QCoh(\mathcal{O}_\mathcal{X})
\longrightarrow
D^{+}_\QCoh(\mathcal{O}_S),
$$
voir
Catégories dérivées des champs, proposition
\ref{stacks-perfect-proposition-derived-direct-image-quasi-coherent}.
Calculons $Rf_{\QCoh, *}\mathcal{O}_\mathcal{X}$.
Comme $D_\QCoh(\mathcal{O}_S)$ est équivalente à
la catégorie dérivée des $k[x]$-modules (voir
Catégories dérivées des schémas, lemme
\ref{perfect-lemma-affine-compare-bounded})
cela revient à calculer
$R\Gamma(\mathcal{X}, \mathcal{O}_\mathcal{X})$.
Nous pouvons pour cela utiliser le recouvrement $S \to \mathcal{X}$ et la suite
spectrale
$$
H^q(S \times_\mathcal{X} \ldots \times_\mathcal{X} S, O)
\Rightarrow H^{p + q}(\mathcal{X}, \mathcal{O}_\mathcal{X})
$$
voir
Cohomologie des champs, proposition
\ref{stacks-cohomology-proposition-smooth-covering-compute-cohomology}.
Remarquons que
$$
S \times_\mathcal{X} \ldots \times_\mathcal{X} S = S \times G^p
$$
qui est affine. Le complexe
$$
k[x] \to \text{Map}(G, k[x]) \to \text{Map}(G^2, k[x]) \to \ldots
$$
calcule donc $R\Gamma(\mathcal{X}, \mathcal{O}_\mathcal{X})$.
Ici, pour $\varphi \in \text{Map}(G^{p - 1}, k[x])$, sa différentielle est
l'application qui envoie $(g_1, \ldots, g_p)$ sur
$$
\varphi(g_2, \ldots, g_p) +
\sum\nolimits_{i = 1}^{p - 1}
(-1)^i\varphi(g_1, \ldots, g_i + g_{i + 1}, \ldots, g_p)
+ (-1)^p\varphi(g_1, \ldots, g_{p - 1}).
$$
C'est précisément le complexe qui calcule la cohomologie du groupe $G$ agissant
trivialement sur $k[x]$ (insérer ici une référence ultérieure). La cohomologie
du groupe cyclique $G$ à coefficients dans $k[x]$ est exactement une copie de $k[x]$ en chaque
degré cohomologique $\geq 0$ (insérer ici une référence ultérieure). Nous en concluons
que
$$
Rf_*\mathcal{O}_\mathcal{X} = \bigoplus\nolimits_{n \geq 0} \mathcal{O}_S[-n]
$$
Considérons maintenant le complexe
$$
E = \bigoplus\nolimits_{m \geq 0} \mathcal{O}_\mathcal{X}[m]
$$
C'est un objet de $D_\QCoh(\mathcal{O}_\mathcal{X})$. Interrompons la
discussion pour établir un résultat général.

\begin{lemma}
\label{lemma-is-limit}
Soit $\mathcal{X}$ un champ algébrique. Soit $K$ un objet de
$D(\mathcal{O}_\mathcal{X})$ dont les faisceaux de cohomologie sont localement
quasi-cohérents (Faisceaux sur les champs, définition
\ref{stacks-sheaves-definition-locally-quasi-coherent})
et satisfont à la propriété de changement de base plat (Cohomologie des champs,
définition \ref{stacks-cohomology-definition-flat-base-change}).
Alors il existe un triangle distingué
$$
K \to
\prod\nolimits_{n \geq 0} \tau_{\geq -n} K \to
\prod\nolimits_{n \geq 0} \tau_{\geq -n} K \to K[1]
$$
dans $D(\mathcal{O}_\mathcal{X})$. Autrement dit, $K$ est la limite dérivée
de ses troncations canoniques.
\end{lemma}

\begin{proof}
Rappelons que nous travaillons sur le « grand site fppf » $\mathcal{X}_{fppf}$
de $\mathcal{X}$ (selon nos conventions
pour les faisceaux de $\mathcal{O}_\mathcal{X}$-modules dans les chapitres
Faisceaux sur les champs et Cohomologie sur les champs). Soit $\mathcal{B}$ l'ensemble
des objets $x$ de $\mathcal{X}_{fppf}$ qui sont au-dessus d'un schéma affine $U$.
En combinant
Faisceaux sur les champs, lemmes
\ref{stacks-sheaves-lemma-compare-fppf-etale},
\ref{stacks-sheaves-lemma-cohomology-restriction},
Descente, lemme \ref{descent-lemma-quasi-coherent-and-flat-base-change},
et
Cohomologie des schémas, lemme
\ref{coherent-lemma-quasi-coherent-affine-cohomology-zero}
nous voyons que $H^p(x, \mathcal{F}) = 0$ si $\mathcal{F}$ est
localement quasi-cohérent et $x \in \mathcal{B}$.
L'assertion découle alors de
Cohomologie sur les sites, lemme \ref{sites-cohomology-lemma-is-limit-dimension},
avec $d = 0$.
\end{proof}

\begin{lemma}
\label{lemma-sum-is-product}
Soit $\mathcal{X}$ un champ algébrique. Si $\mathcal{F}_n$ est une famille
de faisceaux localement quasi-cohérents satisfaisant à la propriété de changement de base plat sur
$\mathcal{X}$, alors $\oplus_n \mathcal{F}_n[n] \to \prod_n \mathcal{F}_n[n]$
est un isomorphisme dans $D(\mathcal{O}_\mathcal{X})$.
\end{lemma}

\begin{proof}
En effet, le lemme \ref{lemma-is-limit} montre que la somme directe
est isomorphe au produit.
\end{proof}

\noindent
Reprenons la discussion. Puisqu'un module quasi-cohérent est localement
quasi-cohérent et satisfait à la propriété de changement de base plat
(Faisceaux sur les champs, lemme \ref{stacks-sheaves-lemma-quasi-coherent}), nous obtenons
$$
E = \prod\nolimits_{m \geq 0} \mathcal{O}_\mathcal{X}[m]
$$
Puisque la cohomologie commute
aux limites, nous avons
$$
Rf_*E = \prod\nolimits_{m \geq 0}
\left(\bigoplus\nolimits_{n \geq 0} \mathcal{O}_S[m - n]\right)
$$
Remarquons que ce complexe n'est pas un objet de $D_\QCoh(\mathcal{O}_S)$,
car son faisceau de cohomologie en degré $0$ est un produit infini de copies
de $\mathcal{O}_S$, qui n'est même pas un
$\mathcal{O}_S$-module localement quasi-cohérent.

\begin{lemma}
\label{lemma-push-not-OK}
Un morphisme quasi-compact et quasi-séparé
$f : \mathcal{X} \to \mathcal{Y}$ de champs algébriques
n'induit pas nécessairement un foncteur
$Rf_* : D_\QCoh(\mathcal{O}_\mathcal{X}) \to
D_\QCoh(\mathcal{O}_\mathcal{Y})$.
\end{lemma}

\begin{proof}
Voir la discussion qui précède.
\end{proof}



\section{Une grande catégorie abélienne}
\label{section-big}

\noindent
Le but de cette section est de donner un exemple d'une « grande » catégorie abélienne
$\mathcal{A}$ et d'objets $M, N$ tels que la collection des classes d'isomorphisme
d'extensions $\Ext_\mathcal{A}(M, N)$ ne soit pas un ensemble.
Cet exemple est dû à Freyd; voir \cite[page 131, Exercise A]{Freyd}.

\medskip\noindent
Nous définissons $\mathcal{A}$ comme suit. Un objet de $\mathcal{A}$
est un triplet $(M, \alpha, f)$, où $M$ est un groupe abélien,
$\alpha$ un ordinal et $f : \alpha \to \text{End}(M)$ une application.
Un morphisme $(M, \alpha, f) \to (M', \alpha', f')$ est donné par un homomorphisme
de groupes abéliens $\varphi : M \to M'$ tel que, pour {\it tout} ordinal $\beta$,
on ait
$$
\varphi \circ f(\beta) = f'(\beta) \circ \varphi
$$
La convention est de poser $f(\beta) = 0$ si $\beta$ n'appartient pas à
$\alpha$ et, de même, $f'(\beta)$ est nul si $\beta$ n'est pas un
élément de $\alpha'$. Nous omettons la vérification que la catégorie ainsi définie
est abélienne.

\medskip\noindent
Considérons l'objet $Z = (\mathbf{Z}, \emptyset, f)$, c'est-à-dire que tous les
opérateurs sont nuls. L'observation est que,
calculé dans $\mathcal{A}$, le groupe $\Ext^1_\mathcal{A}(Z, Z)$ est une classe propre
et non un ensemble. En effet, pour tout ordinal $\alpha$, nous pouvons construire une extension
$(M, \alpha + 1, f)$ de $Z$ par $Z$ dont le groupe sous-jacent est
$M = \mathbf{Z} \oplus \mathbf{Z}$ et pour laquelle $f$ est
partout nulle sauf en
$$
f(\alpha) =
\left(
\begin{matrix}
0 & 1 \\
0 & 0
\end{matrix}
\right).
$$
Cela produit manifestement une classe propre de classes d'isomorphisme d'extensions.
En particulier, la catégorie dérivée de $\mathcal{A}$ possède des classes propres pour
ses collections de morphismes; voir
Catégories dérivées, lemme \ref{derived-lemma-ext-1}. Il faut donc faire preuve
de prudence dans la définition des quotients de Verdier de catégories
triangulées.

\begin{lemma}
\label{lemma-big-abelian-category}
Il existe une « grande » catégorie abélienne $\mathcal{A}$ dont les
groupes $\Ext$ sont des classes propres.
\end{lemma}

\begin{proof}
Voir la discussion qui précède.
\end{proof}





\section{Points faiblement associés et densité schématique}
\label{section-weak-ass-dense}

\noindent
Soit $k$ un corps. Soit $R = k[z, x_i, y_i]/(z^2, zx_iy_i)$, où
$i$ parcourt les éléments de $\mathbf{N}$. Remarquons que $R = R_0 \oplus M_0$, où
$R_0 = k[x_i, y_i]$ est un sous-anneau et $M_0$ un idéal de carré nul,
avec $M_0 \cong R_0/(x_iy_i)$ comme $R_0$-module. L'idéal premier
$\mathfrak p = (z, x_i)$ est faiblement associé à $R$ considéré comme $R$-module
(Algèbre, définition \ref{algebra-definition-weakly-associated}).
En effet, l'élément $z$ de $R_\mathfrak p$ est non nul mais annulé
par $\mathfrak pR_\mathfrak p$. Considérons d'autre part le sous-schéma
ouvert
$$
U = \bigcup D(x_i) \subset \Spec(R) = S
$$
Nous affirmons que $U \subset S$ est schématiquement dense
(Morphismes, définition \ref{morphisms-definition-scheme-theoretically-dense}).
Pour le démontrer, il suffit de voir que $\mathcal{O}_S \to j_*\mathcal{O}_U$
est injectif, où $j : U \to S$ est le morphisme d'inclusion; voir
Morphismes, lemme \ref{morphisms-lemma-characterize-scheme-theoretically-dense}.
Traduit en termes algébriques, il faut montrer que, pour tout $g \in R$,
l'application
$$
R_g \longrightarrow \prod R_{x_ig}
$$
est injective. Écrivons $g = g_0 + m_0$, avec $g_0 \in R_0$ et $m_0 \in M_0$.
Alors $R_g = R_{g_0}$ (détails omis). Nous pouvons donc supposer $g \in R_0$.
Nous pouvons aussi supposer $g$ non nul. Alors $R_g = (R_0)_g \oplus (M_0)_g$.
Puisque $R_0$ est intègre, l'application $(R_0)_g \to \prod (R_0)_{x_ig}$ est
injective. Si $g \in (x_iy_i)$, alors $(M_0)_g = 0$ et il n'y a rien
à démontrer. Si $g \not \in (x_iy_i)$, comme
$(x_iy_i)$ est un idéal radical de $R_0$, il reste à montrer que
$M_0 \to \prod (M_0)_{x_ig}$ est injective. Le noyau de
$R_0 \to M_0 \to (M_0)_{x_n}$ est $(x_iy_i, y_n)$. Comme
$(x_iy_i, y_n)$ est un idéal radical, si $g \not \in (x_iy_i, y_n)$,
alors le noyau de $R_0 \to M_0 \to (M_0)_{x_ng}$ est $(x_iy_i, y_n)$.
Puisque $g \not \in (x_iy_i, y_n)$ pour tout $n \gg 0$, nous concluons que
le noyau est contenu dans $\bigcap_{n \gg 0} (x_iy_i, y_n) = (x_iy_i)$,
comme voulu.

\medskip\noindent
Voici un second exemple, dû à Ofer Gabber. Soit $k$ un corps, et soient
$R$, resp.\ $R'$, l'anneau des fonctions $\mathbf{N} \to k$, resp.\ l'anneau
des fonctions $\mathbf{N} \to k$ constantes à partir d'un certain rang. Alors $\Spec(R)$,
resp.\ $\Spec(R')$, est le compactifié de
Stone-{\v C}ech\footnote{Tout élément $f \in R$ est de la
forme $ue$, où $u$ est une unité et $e$ un idempotent. Alors
Algèbre, lemme \ref{algebra-lemma-ring-with-only-minimal-primes},
montre que $\Spec(R)$ est séparé. D'autre part, $\mathbf{N}$ muni de la
topologie discrète peut être considéré comme un ouvert dense. Étant donnée une application d'ensembles
$\mathbf{N} \to X$ vers un espace topologique séparé et quasi-compact $X$,
nous obtenons un homomorphisme d'anneaux $\mathcal{C}^0(X; k) \to R$, où $\mathcal{C}^0(X; k)$
est la $k$-algèbre des applications localement constantes $X \to k$. Cela donne
$\Spec(R) \to \Spec(\mathcal{C}^0(X; k)) = X$ et démontre la propriété
universelle.} $\beta\mathbf{N}$, resp.\ le compactifié à
un point\footnote{On montre ici qu'il n'existe effectivement
qu'un seul idéal maximal supplémentaire dans $R'$.}
$\mathbf{N}^* = \mathbf{N} \cup \{\infty\}$.
Tous les points sont faiblement associés, puisque tous les idéaux premiers sont minimaux dans les
anneaux $R$ et $R'$.

\begin{lemma}
\label{lemma-example-schematically-dense-missing-weakly-associated-point}
Il existe un schéma réduit $X$ et un ouvert schématiquement dense
$U \subset X$ tel qu'un certain point faiblement associé $x \in X$ n'appartienne pas à $U$.
\end{lemma}

\begin{proof}
Dans le premier exemple, $\mathfrak p \not \in U$ par construction.
Dans les exemples de Gabber, les schémas $\Spec(R)$ et $\Spec(R')$ sont réduits.
\end{proof}



\section{Exemple de non-additivité des traces}
\label{section-non-additive}

\noindent
Soit $k$ un corps et soit $R = k[\epsilon]$ l'anneau des nombres duaux
sur $k$. Autrement dit, $R = k[x]/(x^2)$ et $\epsilon$ est la classe
de $x$ dans $R$. Considérons la suite exacte courte de complexes
$$
\xymatrix{
0 \ar[d] \ar[r] &
R \ar[d]^\epsilon \ar[r]_1 &
R \ar[d] \\
R \ar[r]^1 & R \ar[r] & 0
}
$$
Ici, les colonnes sont les complexes, la première ligne est placée en degré $0$ et
la seconde en degré $1$. Désignons le premier complexe (c'est-à-dire la
colonne de gauche) par $A^\bullet$, le deuxième par $B^\bullet$ et le troisième par
$C^\bullet$. Nous affirmons que le diagramme
\begin{equation}
\label{equation-commutes-up-to-homotopy}
\vcenter{
\xymatrix{
A^\bullet \ar[d]_{1 + \epsilon} \ar[r] &
B^\bullet \ar[r] \ar[d]_1 &
C^\bullet \ar[d]_1 \\
A^\bullet \ar[r] & B^\bullet \ar[r] & C^\bullet
}
}
\end{equation}
est commutatif dans $K(R)$, c'est-à-dire qu'il s'agit d'un diagramme de complexes commutatif à homotopie près.
En effet, le carré de droite est commutatif et celui de gauche l'est modulo
l'homotopie $1 : A^1 \to B^0$. En revanche,
$$
\text{Tr}_{A^\bullet}(1 + \epsilon) + \text{Tr}_{C^\bullet}(1)
\not = \text{Tr}_{B^\bullet}(1).
$$

\begin{lemma}
\label{lemma-nonadditivity-of-trace}
Il existe un anneau $R$, un triangle distingué
$(K, L, M, \alpha, \beta, \gamma)$ dans la catégorie homotopique $K(R)$,
et un endomorphisme $(a, b, c)$ de ce triangle distingué, tels que
$K$, $L$, $M$ soient des complexes parfaits et que
$\text{Tr}_K(a) + \text{Tr}_M(c) \not = \text{Tr}_L(b)$.
\end{lemma}

\begin{proof}
Considérons l'exemple ci-dessus. L'application $\gamma : C^\bullet \to A^\bullet[1]$
est donnée par multiplication par $\epsilon$ en degré $0$; voir
Catégories dérivées, définition \ref{derived-definition-distinguished-triangle}.
Il est donc également vrai que
$$
\xymatrix{
C^\bullet \ar[d] \ar[r]_\gamma & A^\bullet[1] \ar[d] \\
C^\bullet \ar[r]^\gamma & A^\bullet[1]
}
$$
est commutatif dans $K(R)$, puisque $\epsilon(1 + \epsilon) = \epsilon$.
Nous avons donc bien un morphisme de triangles distingués.
\end{proof}




\section{Cohomologie non finie du faisceau structural d'un schéma projectif}
\label{section-nonfinite-h0}

\noindent
Soient $R$ un anneau non noethérien quelconque et $I$ un idéal de $R$
qui n'est pas de type fini. Soit $S = R[x, y]/xyI$, considéré
comme anneau gradué où $x$ et $y$ sont de degré $1$.
Alors $X = \text{Proj}(S)$ est un schéma projectif sur $R$.
Nous affirmons que $H^0(X, \mathcal{O}_X)$ n'est pas
un $R$-module de type fini.

\medskip\noindent
En décomposant les anneaux $S_x$, $S_y$ et $S_{xy}$ suivant le bidegré,
nous pouvons écrire
$$
S_x = \bigoplus_{\substack{(a, b) \in \mathbf{Z}^2 \\ b \geq 0}} R'_{a, b},
\quad\quad
S_y = \bigoplus_{\substack{(a, b) \in \mathbf{Z}^2 \\ a \geq 0}} R''_{a, b},
\quad\quad
S_{xy} = \bigoplus_{\substack{(a, b) \in \mathbf{Z}^2}} R/I,
$$
où les facteurs directs d'indice $(a, b)$ sont en degré $a + b$ et où
$$
R'_{a, b} =
\left\{
\begin{matrix}
R & b = 0, \\
R/I & b \ne 0,
\end{matrix}
\right.
\quad\quad
R''_{a, b} =
\left\{
\begin{matrix}
R & a = 0, \\
R/I & a \ne 0.
\end{matrix}
\right.
$$
Le $R$-module $H^0(X, \mathcal{O}_X)$ se calcule comme le noyau de l'application
$$
(S_x)_0 \times (S_y)_0 \to (S_{xy})_0
$$
définie par $(F, G) \mapsto F - G$. Ici, la notation $A_0$ désigne la
composante de degré $0$ d'un anneau $\mathbf{Z}$-gradué $A$. D'après la description explicite
ci-dessus de $S_x, S_y, S_{xy}$, ce noyau est isomorphe
au sous-module $M \subset R^2$ formé des couples $(F, G)$
tels que $F - G \in I$. Puisque l'application $M \to I$ de $R$-modules
donnée par $(F, G) \mapsto F - G$ est surjective,
$M$ n'est pas un $R$-module de type fini.

\medskip\noindent
La surjection $R[x, y] \to S$ correspond à une
immersion fermée $i : X \to \mathbf{P}^1_R$. Alors $\mathcal{F} = i_*\mathcal{O}_X$
est un $\mathcal{O}_{\mathbf{P}^1_R}$-module quasi-cohérent de type fini
dont le $H^0$ est le module $M$ non de type fini
calculé ci-dessus.

\begin{lemma}
\label{lemma-nonfinite-h0}
$H^0$ non fini.
\begin{enumerate}
\item Il existe un anneau $R$ et un schéma projectif $X$
sur $R$ tels que $H^0(X, \mathcal{O}_X)$ ne soit pas
un $R$-module de type fini.
\item Il existe un anneau $R$ et un module quasi-cohérent de type fini
$\mathcal{F}$ sur $\mathbf{P}^1_R$ tels que
$H^0(\mathbf{P}^1_R, \mathcal{F})$ ne soit pas un $R$-module de type fini.
\end{enumerate}
\end{lemma}

\begin{proof}
Voir la discussion qui précède.
\end{proof}









\section{Être projectif n'est pas local sur la base}
\label{section-non-descending-property-projective}

\noindent
Dans le chapitre sur la descente, nous avons vu que de nombreuses propriétés des morphismes
sont locales sur la base, même pour la topologie fpqc. Voir
Descente, sections \ref{descent-section-descending-properties-morphisms},
\ref{descent-section-descending-properties-morphisms-fpqc} et
\ref{descent-section-descending-properties-morphisms-fppf}.
Ce n'est pas le cas de la projectivité des morphismes.

\begin{lemma}
\label{lemma-non-descending-property-projective}
Les propriétés
\begin{enumerate}
\item[] $\mathcal{P}(f) =$ « $f$ est projectif », et
\item[] $\mathcal{P}(f) =$ « $f$ est quasi-projectif »
\end{enumerate}
ne sont pas locales sur la base pour la topologie de Zariski. A fortiori, elles ne sont pas locales
sur la base pour la topologie fpqc.
\end{lemma}

\begin{proof}
Suivant Hironaka \cite[exemple B.3.4.1]{H},
nous construisons un morphisme propre entre variétés complexes lisses de dimension 3
$f:V_Y\to Y$
qui est localement projectif pour la topologie de Zariski, mais non projectif. Comme $f$ est propre
et non projectif, il n'est pas non plus quasi-projectif.

\medskip\noindent
Soit $Y$ l'espace projectif de dimension 3 sur les nombres complexes ${\mathbf C}$.
Soient $C$ et $D$ des coniques lisses de $Y$ telles
que le sous-schéma fermé $C\cap D$ soit réduit et consiste
en deux points complexes $P$ et $Q$. (Par exemple,
prenons $C=\{ [x,y,z,w]: xy=z^2, w=0\}$, $D=\{ [x,y,z,w]:
xy=w^2, z=0\}$, $P=[1,0,0,0]$,
et $Q=[0,1,0,0]$.)
Sur $Y-Q$, éclatons d'abord la courbe $C$, puis la
transformée stricte de la courbe $D$ (Diviseurs, définition
\ref{divisors-definition-strict-transform}). Sur $Y-P$, éclatons d'abord
la courbe $D$, puis la transformée stricte de la courbe
$C$. Au-dessus de $Y-P-Q$, les deux variétés ainsi construites sont canoniquement
isomorphes; nous pouvons donc les recoller au-dessus de $Y-P-Q$. Le résultat
est une variété lisse propre $V_Y$ de dimension 3 sur ${\mathbf C}$. Le morphisme
$f:V_Y\to Y$ est propre et localement projectif pour la topologie de Zariski (car
c'est un éclatement au-dessus de $Y-P$ et de $Y-Q$), d'après Diviseurs,
lemme \ref{divisors-lemma-blowing-up-projective}. Nous allons montrer que
$V_Y$ n'est pas projectif sur ${\mathbf C}$. Il en résultera que
$f$ n'est pas projectif.

\medskip\noindent
Pour cela, soit $L$ l'image réciproque dans $V_Y$ d'un point complexe
de $C-P-Q$, et $M$ celle d'un point complexe de $D-P-Q$.
Alors $L$ et $M$ sont isomorphes à la droite projective
${\mathbf P}^1_{{\mathbf C}}$.
Soit ensuite $E$ le sous-schéma de $V_Y$ donné par l'image réciproque de $C\cup D\subset Y$ dans $V_Y$;
ainsi, $E\rightarrow C\cup D$ est un morphisme propre, dont les fibres
sont isomorphes à ${\mathbf P}^1$ au-dessus de $(C\cup D)-\{P,Q\}$.
L'image
réciproque de $P$ dans $E$ est la réunion de deux droites $L_0$ et $M_0$, et nous avons
les équivalences rationnelles de cycles $L\sim L_0+M_0$ et $M\sim M_0$ sur $E$
(puisque $C$ et $D$ sont isomorphes à ${\mathbf P}^1$).
Remarquons l'asymétrie qui résulte de l'ordre dans lequel nous avons éclaté
les deux courbes. Au voisinage de $Q$, la situation inverse se produit. L'image réciproque
de $Q$ est donc la réunion de deux droites $L_0'$ et $M_0'$, et nous avons
les équivalences rationnelles $L\sim L_0'$ et $M\sim L_0'+M_0'$ sur $E$.
En combinant ces équivalences, nous obtenons $L_0+M_0'\sim 0$
sur $E$, donc sur $V_Y$. Si $V_Y$ était projectif sur ${\mathbf C}$,
il posséderait
un fibré inversible ample $H$, de degré $> 0$ sur toute courbe
de $V_Y$. En particulier, $H$ serait de degré positif sur $L_0+M_0'$,
ce qui contredit le fait que le degré d'un fibré inversible est bien défini
sur les 1-cycles modulo équivalence rationnelle d'un schéma propre
sur un corps (Homologie de Chow,
lemme \ref{chow-lemma-proper-pushforward-rational-equivalence}
et lemme \ref{chow-lemma-factors}).
Ainsi, $V_Y$ n'est pas projectif sur ${\mathbf C}$.
\end{proof}

\noindent
Dans une autre terminologie, la variété de dimension 3 de Hironaka $V_Y$ est une petite
résolution de l'éclatement $Y'$ de $Y$ le long du sous-schéma réduit
$C\cup D$; ici, $Y'$ possède deux singularités nodales. Si l'on définit $Z$ en éclatant
$Y$ le long de $C$, puis le long de la transformée stricte
de $D$, alors $Z$ est une variété projective lisse de dimension 3, et la variété non projective
$V_Y$ diffère de $Z$ par un « flop » au-dessus de $Y-P$.


\section[Descente non effective pour des schémas projectifs]{Données de descente non effectives pour les schémas projectifs}
\label{section-non-effective-descent-projective}

\noindent
Dans le chapitre sur la descente, nous avons vu que les données de descente de schémas relativement
à un morphisme fpqc sont effectives pour plusieurs classes
de morphismes. En particulier, les morphismes affines et, plus généralement,
les morphismes quasi-affines satisfont à la descente pour les recouvrements fpqc
(Descente, lemme \ref{descent-lemma-quasi-affine}).
Ce n'est pas le cas des morphismes projectifs.

\begin{lemma}
\label{lemma-non-effective-descent-projective}
Il existe un recouvrement étale de schémas $X\to S$ et une donnée de descente
$(V/X,\varphi)$ relativement à $X\to S$ telle que
$V\to X$ soit projectif,
mais que la donnée de descente ne soit pas effective dans la catégorie des schémas.
\end{lemma}

\begin{proof}
Nous imitons l'exemple de Hironaka d'un espace algébrique complexe
lisse et séparé de dimension 3
qui n'est pas un schéma \cite[exemple B.3.4.2]{H}.

\medskip\noindent
Considérons l'action du groupe $G = \mathbf{Z}/2 = \{1, g\}$
sur l'espace projectif de dimension 3
$\mathbf{P}^3$ sur les nombres complexes, donnée par
$$
g[x,y,z,w] = [y,x,w,z].
$$
L'action est libre en dehors des deux droites disjointes
$L_1=\{ [x,x,z,z]\}$ et $L_2=\{ [x,-x,z,-z]\}$ de
${\mathbf P}^3$. Soit $Y={\mathbf P}^3-(L_1\cup L_2)$. Il existe un
schéma quasi-projectif lisse $S=Y/G$ sur ${\mathbf C}$ tel que
$Y\to S$ soit un torseur sous $G$ (Groupoïdes,
définition \ref{groupoids-definition-principal-homogeneous-space}).
Explicitement, nous pouvons définir $S$ comme l'image de l'ouvert $Y$
de ${\mathbf P}^3$ par le morphisme
\begin{align*}
{\mathbf P}^3 & \to \text{Proj } {\mathbf C}[x,y,z,w]^G\\
   & = \text{Proj } {\mathbf C}[u_0,u_1,v_0,v_1,v_2]/(v_0v_1=v_2^2),
\end{align*}
où $u_0=x+y$, $u_1=z+w$, $v_0=(x-y)^2$, $v_1=(z-w)^2$
et $v_2=(x-y)(z-w)$, l'anneau étant gradué avec $u_0,u_1$
en degré 1 et $v_0,v_1,v_2$ en degré 2.

\medskip\noindent
Soient $C=\{ [x,y,z,w]: xy=z^2, w=0\}$ et $D=\{ [x,y,z,w]:
xy=w^2, z=0\}$. Ce sont des coniques lisses de ${\mathbf P}^3$, contenues
dans l'ouvert $G$-invariant $Y$, avec $g(C)=D$. De plus,
$C\cap D$ consiste en deux points $P:=[1,0,0,0]$
et $Q:=[0,1,0,0]$, que l'action
de $G$ échange.

\medskip\noindent
Soit $V_Y\to Y$ le schéma défini au-dessus de $Y-P$
en éclatant $D$, puis la transformée stricte
de $C$, et au-dessus de $Y-Q$ en éclatant $C$, puis
la transformée stricte de $D$. (C'est la même construction
que dans la démonstration du lemme \ref{lemma-non-descending-property-projective},
à ceci près que $Y$ désigne ici un ouvert de ${\mathbf P}^3$
et non ${\mathbf P}^3$ tout entier.)
Alors l'action de $G$ sur $Y$ se relève en une action de $G$ sur $V_Y$
qui échange les images réciproques de $Y-P$ et de $Y-Q$. Cette action
de $G$ sur $V_Y$ fournit une donnée de descente $(V_Y/Y,\varphi_Y)$
sur $V_Y$ relativement au torseur sous $G$
$Y\to S$. Le morphisme $V_Y\to Y$
est propre mais non projectif, comme le montre la démonstration
du lemme \ref{lemma-non-descending-property-projective}.

\medskip\noindent
Soit $X$ la réunion disjointe des ouverts $Y-P$ et $Y-Q$;
nous avons alors des morphismes étales surjectifs $X\to Y\to S$.
Soit $V$ l'image réciproque de $V_Y\to Y$ sur $X$; alors le morphisme
$V\to X$ est projectif, puisque $V_Y\to Y$ est un éclatement au-dessus de chacun des ouverts
$Y-P$ et $Y-Q$. De plus, la donnée de descente $(V_Y/Y,\varphi_Y)$
induit par image réciproque une donnée de descente $(V/X,\varphi)$ relativement au
recouvrement étale $X\to S$.

\medskip\noindent
Supposons cette donnée de descente effective dans la catégorie
des schémas. Il existe alors un schéma $U\to S$
dont l'image réciproque est le morphisme $V\to X$, muni
de sa donnée de descente. Alors $U$ serait le quotient
de $V_Y$ par l'action de $G$.
$$
\xymatrix{
V \ar[r]\ar[d]&  X\ar[d] \\
V_Y \ar[r]\ar[d]& Y\ar[d] \\
U \ar[r]& S
}
$$

\medskip\noindent
Soit $E$ l'image réciproque de $C\cup D\subset Y$ dans $V_Y$;
ainsi, $E\rightarrow C\cup D$ est un morphisme propre, dont les fibres
sont isomorphes à ${\mathbf P}^1$ au-dessus de $(C\cup D)-\{P,Q\}$.
L'image réciproque de $P$ dans $E$ est la réunion de deux droites $L_0$
et $M_0$. Il en résulte que l'image réciproque de $Q=g(P)$ dans $E$
est la réunion des deux droites $L_0'=g(M_0)$ et $M_0'=g(L_0)$.
Comme le montre la démonstration
du lemme \ref{lemma-non-descending-property-projective},
nous avons une équivalence rationnelle $L_0+M_0'=L_0+g(L_0)\sim 0$ sur $E$.

\medskip\noindent
Par descente des sous-schémas fermés, il existe une courbe $L_1\subset U$
(isomorphe à ${\mathbf P}^1$)
dont l'image réciproque dans $V_Y$ est $L_0\cup g(L_0)$. (Utiliser Descente, lemme
\ref{descent-lemma-affine}, en remarquant qu'une immersion fermée est un morphisme
affine.)
Soit $R$ un point complexe de $L_1$. Puisque
nous avons supposé que $U$ est un schéma, nous pouvons choisir une fonction
$f$ dans l'anneau local $O_{U,R}$ qui s'annule en $R$, mais non
sur toute la courbe $L_1$. Soit $D_{\text{loc}}$ une composante irréductible
du fermé $\{f = 0\}$ de $\Spec O_{U,R}$; alors
$D_{\text{loc}}$ est de codimension 1.
L'adhérence de $D_{\text{loc}}$ dans $U$ est un diviseur irréductible $D_U$
de $U$ qui contient le point $R$, mais non toute la courbe $L_1$.
L'image réciproque de $D_U$ dans $V_Y$ est un diviseur effectif $D$
qui rencontre $L_0\cup g(L_0)$ mais ne contient aucune des deux
courbes $L_0$ et $g(L_0)$.

\medskip\noindent
Puisque la variété complexe $V_Y$ de dimension 3 est lisse, $O(D)$ est un fibré
inversible sur $V_Y$. Nous utilisons ici le fait qu'un anneau local régulier est à factorisation unique,
autrement dit factoriel; voir
Compléments sur l'algèbre, lemme \ref{more-algebra-lemma-regular-local-UFD}.
La restriction de $O(D)$ à la surface propre
$E\subset V_Y$ est un fibré inversible de degré positif sur le 1-cycle
$L_0+g(L_0)$, d'après les propriétés de $D$. Puisque
$L_0+g(L_0)\sim 0$ sur $E$, cela contredit
le fait que le degré d'un fibré inversible est bien défini
sur les 1-cycles modulo équivalence rationnelle d'un schéma propre
sur un corps (Homologie de Chow,
lemme \ref{chow-lemma-proper-pushforward-rational-equivalence}
et lemme \ref{chow-lemma-factors}). La donnée de descente
$(V/X,\varphi)$ n'est donc pas effective; autrement dit, $U$ n'existe pas
comme schéma.
\end{proof}

\noindent
Dans cet exemple, la donnée de descente {\it est }effective dans la catégorie
des espaces algébriques. Plus précisément,
$U$ existe comme espace algébrique lisse et séparé
de dimension 3 sur ${\mathbf C}$,
par exemple d'après Espaces algébriques, lemme \ref{spaces-lemma-quotient}.
La variété $U$ de Hironaka, de dimension 3, est une petite résolution de l'éclatement $S'$
de la variété quasi-projective lisse $S$ de dimension 3 le long de la courbe nodale irréductible
$(C\cup D)/G$; la variété $S'$ de dimension 3 possède une singularité nodale. L'autre
petite résolution de $S'$ (qui diffère de $U$ par un « flop »)
est encore un espace algébrique qui n'est pas un schéma.






\section{Une famille de courbes dont l'espace total n'est pas un schéma}
\label{section-family-of-curves}

\noindent
Dans Quot, section \ref{quot-section-curves}, nous définissons une famille de courbes sur
un schéma $S$ comme un morphisme propre, plat et de présentation finie
de dimension relative $\leq 1$ d'un espace algébrique $X$ vers $S$.
Si $S$ est le spectre d'un anneau local noethérien complet, alors
$X$ est un schéma ; voir
Compléments sur les morphismes d'espaces, lemme
\ref{spaces-more-morphisms-lemma-projective-over-complete}.
Dans cette section, nous montrons que ce résultat est faux en général.

\medskip\noindent
Soit $k$ un corps. Partons d'un morphisme propre et plat
$$
Y \longrightarrow \mathbf{A}^1_k
$$
et d'un point $y \in Y(k)$ situé au-dessus de $0 \in \mathbf{A}^1_k(k)$
possédant les propriétés suivantes :
\begin{enumerate}
\item la fibre $Y_0$ est une courbe lisse géométriquement irréductible sur $k$,
\item pour tout sous-schéma fermé propre $T \subset Y$ dominant
$\mathbf{A}^1_k$, l'intersection $T \cap Y_0$ contient au moins un point
distinct de $y$.
\end{enumerate}
À partir d'une telle surface, nous construisons notre exemple comme suit.
$$
\xymatrix{
Y \ar[rd] & Z \ar[d] \ar[l] \ar[r] & X \ar[ld] \\
& \mathbf{A}^1_k
}
$$
Ici, $Z \to Y$ est l'éclatement de $Y$ en $y$. Soit $E \subset Z$ le
diviseur exceptionnel et soit $C \subset Z$ la transformée stricte de $Y_0$.
On a $Z_0 = E \cup C$ au sens schématique (pour le voir, utiliser que
$Y$ est lisse en $y$ et que, de plus, $Y \to \mathbf{A}^1_k$ est lisse en $y$).
D'après les résultats d'Artin (\cite{ArtinII} ; utiliser
Réduction semi-stable, lemme \ref{models-lemma-properties-form}
pour voir que le fibré normal de $C$ est négatif),
nous pouvons contracter la courbe $C$ dans $Z$ et obtenir un espace algébrique $X$
comme dans le diagramme. Soit $x \in X(k)$ l'image de $C$.

\medskip\noindent
Nous affirmons que $X$ n'est pas un schéma. En effet, s'il en était un,
il existerait un voisinage ouvert affine $U \subset X$ de $x$.
Posons $T = X \setminus U$. Alors $T$ domine $\mathbf{A}^1_k$
(car les fibres de $X \to \mathbf{A}^1_k$ sont propres de dimension $1$
alors que les fibres de $U \to \mathbf{A}^1_k$ sont affines, donc différentes).
Soit $T' \subset Z$ le sous-schéma fermé qui s'envoie isomorphiquement
sur $T$ (car $x \not \in T$). Alors l'image de $T'$ dans $X$
contredit la condition (2) ci-dessus (car $T' \cap Z_0$ est contenu dans
le diviseur exceptionnel $E$ de l'éclatement $Z \to Y$).

\medskip\noindent
Pour achever la discussion, il nous reste à construire $Y$.
Nous supposerons que la caractéristique de $k$ n'est pas $3$.
Écrivons $\mathbf{A}^1_k = \Spec(k[t])$ et prenons
$$
Y \quad : \quad T_0^3 + T_1^3 + T_2^3 - tT_0T_1T_2 = 0
$$
dans $\mathbf{P}^2_{k[t]}$. La fibre pour $t = 0$
est une courbe projective lisse de genre $1$. Sur l'ouvert affine
$V_+(T_0)$, nous obtenons l'équation affine
$$
1 + x^3 + y^3 - txy = 0
$$
qui définit une surface lisse sur $k$. Par symétrie, il en va de même sur les
autres ouverts affines, et nous voyons donc que $Y$ est une surface lisse.
Enfin, l'équation affine montre aussi que le corps des fractions
est $k(x, y)$ ; la surface $Y$ est donc rationnelle.
Or le groupe de Picard d'une surface rationnelle est de type fini
(insérer ici une référence future).
Ainsi, pour choisir $y \in Y_0(k)$ possédant la propriété (2),
il suffit de choisir $y$ de telle sorte que
\begin{equation}
\label{equation-three}
\mathcal{O}_{Y_0}(ny) \not \in \Im(\Pic(Y) \to \Pic(Y_0))
\text{ pour tout }n > 0
\end{equation}
En effet, la somme des composantes irréductibles de dimension $1$ d'un $T$
contredisant (2) donnerait un diviseur de Cartier effectif dont l'intersection avec $Y_0$
serait le diviseur $ny$ pour un certain $n \geq 1$, et nous conclurions
que $\mathcal{O}_{Y_0}(ny)$ appartient à l'image de l'application de restriction.
Remarquons que, puisque $Y_0$ est de genre $\geq 1$, l'application
$$
Y_0(k) \to \Pic(Y_0),\quad y \mapsto \mathcal{O}_{Y_0}(y)
$$
est injective. Si maintenant $k$ est un corps algébriquement clos non dénombrable,
la dénombrabilité de $\Pic(Y)$ et la remarque
précédente permettent de trouver un $y \in Y_0(k)$ satisfaisant (\ref{equation-three}),
et donc (2).

\begin{lemma}
\label{lemma-family-of-curves-not-scheme}
Il existe un corps $k$ et une famille de courbes
$X \to \mathbf{A}^1_k$ telle que $X$ ne soit pas un schéma.
\end{lemma}

\begin{proof}
Voir la discussion ci-dessus.
\end{proof}






\section{Changement de base dérivé}
\label{section-derived-base-change}

\noindent
Soit $R \to R'$ un morphisme d'anneaux. Dans
Compléments d'algèbre, section \ref{more-algebra-section-derived-base-change},
nous construisons un foncteur de changement de base dérivé
$- \otimes_R^\mathbf{L} R' : D(R) \to D(R')$.
Soit ensuite $R \to A$ un second morphisme d'anneaux. Considérons le diagramme
$$
\xymatrix{
	A \ar[r] & A \otimes_R R' \ar@{=}[r] & A' \\
R \ar[u] \ar[r] & R' \ar[u] \ar[ur]
}
$$
Étant donné un $A$-module $M$, le produit tensoriel $M \otimes_R R'$ est un
$A \otimes_R R'$-module, c'est-à-dire un $A'$-module. Au morphisme d'anneaux
$A \to A'$ est associé un foncteur dérivé
$$
- \otimes_A^\mathbf{L} A' : D(A) \longrightarrow D(A')
$$
mais, en général, ce foncteur ne coïncide pas avec $- \otimes_R^\mathbf{L} R'$.
Plus précisément, pour $K \in D(A)$, le morphisme canonique
$$
K \otimes_R^{\mathbf{L}} R' \longrightarrow
K \otimes_A^{\mathbf{L}} A'
$$
dans $D(R')$, construit dans
Compléments d'algèbre, équation (\ref{more-algebra-equation-comparison-map}),
n'est pas un isomorphisme en général. On peut donc se demander s'il existe
un « foncteur de changement de base dérivé » $T : D(A) \to D(A')$, c'est-à-dire un foncteur
tel que l'image de $T(K)$ soit $K \otimes_R^\mathbf{L} R'$ dans $D(R')$.
Dans cette section, nous montrons qu'un tel foncteur n'existe pas en général.

\medskip\noindent
Soit $k$ un corps. Posons $R = k[x, y]$, $R' = R/(xy)$ et $A = R/(x^2)$.
L'objet $A \otimes_R^\mathbf{L} R'$ de $D(R')$ est représenté par
$$
x^2 : R' \longrightarrow R'
$$
et $H^0(A \otimes_R^\mathbf{L} R') = A \otimes_R R'$. Nous affirmons qu'il
n'existe aucun objet $E$ de $D(A \otimes_R R')$ dont l'image soit
$A \otimes_R^\mathbf{L} R'$ dans $D(R')$. En effet, pour un tel $E$, le module
$H^0(E)$ serait libre ; par conséquent, $E$ se décomposerait sous la forme
$H^0(E)[0] \oplus H^{-1}(E)[1]$. Or il est facile de voir que
$A \otimes_R^\mathbf{L} R'$ n'est pas isomorphe à la somme de ses
groupes de cohomologie dans $D(R')$.

\begin{lemma}
\label{lemma-no-derived-base-change}
Soient $R \to R'$ et $R \to A$ des morphismes d'anneaux. En général, il
n'existe pas de foncteur $T : D(A) \to D(A \otimes_R R')$
entre catégories triangulées tel qu'à tout $A$-module $M$ soit associé un
objet $T(M)$ de $D(A \otimes_R R')$ qui s'envoie sur
$M \otimes_R^\mathbf{L} R'$ par l'application $D(A \otimes_R R') \to D(R')$.
\end{lemma}

\begin{proof}
Voir la discussion ci-dessus.
\end{proof}




\section{Un objet compact intéressant}
\label{section-interesting-compact}


\noindent
Soit $R$ un anneau et soit $(A, \text{d})$ une $R$-algèbre différentielle graduée.
Si $A = R$, nous savons que tout objet compact de $D(A, \text{d}) = D(R)$
est représenté par un complexe fini de modules projectifs de type fini. En
d'autres termes, les objets compacts sont parfaits ; voir
Compléments d'algèbre, proposition \ref{more-algebra-proposition-perfect-is-compact}.
Dans le langage des modules différentiels gradués, la question analogue
serait : « Tout objet compact de $D(A, \text{d})$ est-il représenté
par un $A$-module différentiel gradué $P$, de type fini et
projectif comme module gradué ? »

\medskip\noindent
Pour une algèbre différentielle graduée générale, ce n'est pas vrai. En effet,
soit $k$ un corps de caractéristique $2$
(de sorte que nous n'ayons pas à nous soucier des signes).
Soit $A = k[x, y]/(y^2)$
avec
\begin{enumerate}
\item $x$ est de degré $0$,
\item $y$ est de degré $-1$,
\item $\text{d}(x) = 0$, et
\item $\text{d}(y) = x^2 + x$.
\end{enumerate}
Alors $x : A \to A$ est un projecteur dans $K(A, \text{d})$.
Par conséquent, on a
$$
A = \Ker(x) \oplus \Im(1 - x)
$$
dans $K(A, \text{d})$ ; voir
Algèbre différentielle graduée, lemme \ref{dga-lemma-homotopy-direct-sums}, et
Catégories dérivées, lemme
\ref{derived-lemma-projectors-have-images-triangulated}.
Il est clair que $A$ est un objet compact de $D(A, \text{d})$.
Il s'ensuit que $\Ker(x)$ est un objet compact de $D(A, \text{d})$,
comme il résulte de
Catégories dérivées, lemme \ref{derived-lemma-compact-objects-subcategory}.

\medskip\noindent
Supposons ensuite que $M$ soit un $A$-module différentiel gradué (à droite)
représentant $\Ker(x)$ et que $M$ soit de type fini et projectif
comme $A$-module gradué. Tout module gradué projectif de type fini
sur $k[x, y]/(y^2)$ étant gradué libre, nous voyons que $M$ est
libre de type fini comme $k[x, y]/(y^2)$-module gradué (autrement dit,
lorsque l'on oublie la différentielle). Posons $N = M/M(x^2 + x)$.
Considérons la suite exacte
$$
0 \to M \xrightarrow{x^2 + x} M \to N \to 0
$$
Comme $x^2 + x$ est de degré $0$, appartient au centre de $A$ et
$\text{d}(x^2 + x) = 0$, c'est une suite exacte courte de
$A$-modules différentiels gradués. De plus, puisque $\text{d}(y) = x^2 + x$,
la différentielle de $N$ est linéaire. Les applications
$$
H^{-1}(N) \to H^0(M)
\quad\text{et}\quad
H^0(M) \to H^0(N)
$$
sont des isomorphismes, car $H^*(M) = H^0(M) = k$ puisque $M \cong \Ker(x)$
dans $D(A, \text{d})$. Un calcul de l'application de connexion montre que
$H^*(N) = k[x, y]/(x, y^2)$ comme module gradué ; nous omettons
les détails. Comme $N$ est un $k[x, y]/(y^2, x^2 + x)$-module libre,
nous disposons d'une résolution
$$
\ldots \to N[2] \xrightarrow{y} N[1] \xrightarrow{y} N \to N/Ny \to 0
$$
compatible aux différentielles. Comme $N$ est borné et que
$H^0(N) = k[x,y]/(x, y^2)$, il résulte de
Homologie, lemme \ref{homology-lemma-first-quadrant-ss},
que $H^0(N/Ny) = k[x]/(x)$. Mais $N/Ny$ est un complexe fini de
$k[x]/(x^2 + x) = k \times k$-modules libres ; sa cohomologie
doit donc être de dimension paire, ce qui est une contradiction.

\begin{lemma}
\label{lemma-no-good-representatif-compact-object}
Il existe une algèbre différentielle graduée $(A, \text{d})$ et un objet
compact $E$ de $D(A, \text{d})$ tel que $E$ ne puisse pas
être représenté par un $A$-module différentiel gradué de type fini et
projectif comme module gradué.
\end{lemma}

\begin{proof}
Voir la discussion ci-dessus.
\end{proof}







\section{Deux catégories différentielles graduées}
\label{section-nongraded-differential-graded}

\noindent
Dans cette section, nous construisons deux catégories différentielles graduées
satisfaisant aux axiomes (A), (B) et (C) de
Algèbre différentielle graduée, situation \ref{dga-situation-ABC},
mais dont les objets ne sont pas munis d'une $\mathbf{Z}$-graduation.

\medskip\noindent
{\bf Exemple I.} Soit $X$ un espace topologique. Notons
$\underline{\mathbf{Z}}$ le faisceau constant de valeur $\mathbf{Z}$.
Soit $A$ un $\underline{\mathbf{Z}}$-torseur. Dans ce contexte, un faisceau
de groupes abéliens $\mathcal{F}$ est dit {\it $A$-gradué} si, pour toute
section locale $a \in A(U)$, il existe un projecteur
$p_a : \mathcal{F}|_U \to \mathcal{F}|_U$ tel que, pour tout isomorphisme
local $\underline{\mathbf{Z}}|_U \to A|_U$, on ait
$\mathcal{F}|_U = \bigoplus_{n \in \mathbf{Z}} p_n(\mathcal{F})$.
Autrement dit, localement sur $X$, le faisceau abélien $\mathcal{F}$ possède
une $\mathbf{Z}$-graduation, mais, sur les intersections, les différents choix
de graduation diffèrent d'un décalage de degré donné par les fonctions de
transition du torseur $A$.
Nous disons qu'un couple $(\mathcal{F}, \text{d})$ est un
{\it complexe $A$-gradué de faisceaux abéliens} si $\mathcal{F}$
est un faisceau abélien $A$-gradué et si $\text{d} : \mathcal{F} \to \mathcal{F}$ est une différentielle,
c'est-à-dire si $\text{d}^2 = 0$ et si
$p_{a + 1} \circ \text{d} = \text{d} \circ p_a$ pour toute section locale
$a$ de $A$. Autrement dit, $\text{d}(p_a(\mathcal{F}))$
est contenu dans $p_{a + 1}(\mathcal{F})$.

\medskip\noindent
Considérons ensuite la catégorie $\mathcal{A}$ définie comme suit :
\begin{enumerate}
\item ses objets sont les complexes $A$-gradués de faisceaux abéliens ;
\item pour deux objets $(\mathcal{F}, \text{d})$, $(\mathcal{G}, \text{d})$, nous posons
$$
\Hom_\mathcal{A}((\mathcal{F}, \text{d}), (\mathcal{G}, \text{d})) =
\bigoplus \Hom^n(\mathcal{F}, \mathcal{G})
$$
où $\Hom^n(\mathcal{F}, \mathcal{G})$ est le groupe des
morphismes de faisceaux abéliens $f$ tels que
$f(p_a(\mathcal{F})) \subset p_{a + n}(\mathcal{G})$ pour toute
section locale $a$ de $A$. Nous prenons pour différentielle
$\text{d}(f) = \text{d} \circ f - (-1)^n f \circ \text{d}$ ; voir
Algèbre différentielle graduée, exemple \ref{dga-example-category-complexes}.
\end{enumerate}
Nous omettons la vérification qu'il s'agit bien d'une catégorie différentielle
graduée satisfaisant à (A), (B) et (C). Toutes ces propriétés se vérifient
localement sur $X$, où l'on retrouve simplement la catégorie différentielle
graduée des complexes de faisceaux abéliens. Nous obtenons ainsi une catégorie
triangulée $K(\mathcal{A})$.

\medskip\noindent
Catégorie dérivée tordue de $X$. Remarquons qu'à tout objet
$(\mathcal{F}, \text{d})$ de $\mathcal{A}$ est associé un faisceau de
cohomologie $A$-gradué bien défini $H(\mathcal{F}, \text{d})$.
La notion de quasi-isomorphisme dans $K(\mathcal{A})$ est donc claire.
En inversant les quasi-isomorphismes, nous obtenons la
{\it catégorie dérivée $D(\mathcal{A})$ des complexes de faisceaux
$A$-gradués}. Si $A$ est le torseur trivial, alors $D(\mathcal{A})$
est égale à $D(X)$ ; pour un torseur non trivial, on obtient en revanche
une sorte de catégorie dérivée {\it tordue} de $X$.

\medskip\noindent
{\bf Exemple II.} Soit $C$ une courbe lisse sur un corps parfait $k$ de
caractéristique $2$. Alors $\Omega_{C/k}$ possède une racine carrée canonique.
Plus précisément, on peut écrire $\Omega_{C/k} = \mathcal{L}^{\otimes 2}$
de telle sorte que, pour toute fonction locale $f$ sur $C$, la section
$\text{d}(f)$ soit égale à $s^{\otimes 2}$ pour une certaine section locale
$s$ de $\mathcal{L}$. La « raison » est que
$$
\text{d}(a_0 + a_1t + \ldots +a_dt^d) =
(\sum\nolimits_{i\text{ impair}} a_i^{1/2} t^{(i - 1)/2})^2\text{d}t
$$
(insérer ici une référence future). Cela détermine en particulier une
connexion canonique
$$
\nabla_{can} :
\Omega_{C/k}
\longrightarrow
\Omega_{C/k} \otimes_{\mathcal{O}_C} \Omega_{C/k}
$$
dont la $2$-courbure est nulle (à savoir l'unique connexion pour laquelle
les carrés ont une dérivée nulle).
Remarquons que la catégorie des fibrés vectoriels
munis d'une connexion est une catégorie tensorielle ; nous obtenons donc aussi
des connexions canoniques
$\nabla_{can}$ sur les faisceaux inversibles $\Omega_{C/k}^{\otimes n}$ pour tous $n \in \mathbf{Z}$.

\medskip\noindent
Soit $\mathcal{A}$ la catégorie définie comme suit :
\begin{enumerate}
\item ses objets sont les couples $(\mathcal{F}, \nabla)$ formés d'un
faisceau localement libre de type fini $\mathcal{F}$ muni d'une connexion
$$
\nabla :
\mathcal{F}
\longrightarrow
\mathcal{F} \otimes_{\mathcal{O}_C} \Omega_{C/k}
$$
dont la $2$-courbure est nulle ;
\item les morphismes entre $(\mathcal{F}, \nabla_\mathcal{F})$ et
$(\mathcal{G}, \nabla_\mathcal{G})$
sont donnés par
$$
\Hom_\mathcal{A}((\mathcal{F}, \nabla_\mathcal{F}),
(\mathcal{G}, \nabla_\mathcal{G})) =
\bigoplus \Hom_{\mathcal{O}_C}(\mathcal{F},
\mathcal{G} \otimes_{\mathcal{O}_C} \Omega_{C/k}^{\otimes n})
$$
Pour un élément
$f : \mathcal{F} \to \mathcal{G} \otimes \Omega_{C/k}^{\otimes n}$
de degré $n$, nous posons
$$
\text{d}(f) =
\nabla_{\mathcal{G} \otimes \Omega_{C/k}^{\otimes n}} \circ f +
f \circ \nabla_\mathcal{F}
$$
moyennant les identifications convenables.
\end{enumerate}
Nous omettons la vérification que ceci définit une catégorie différentielle
graduée satisfaisant aux propriétés (A), (B) et (C). Nous obtenons ainsi une
catégorie d'homotopie triangulée $K(\mathcal{A})$.

\medskip\noindent
Si $C = \mathbf{P}^1_k$, alors $K(\mathcal{A})$ est la catégorie nulle.
En revanche, si $C$ est une courbe lisse propre de genre $> 1$, alors
$K(\mathcal{A})$ n'est pas nulle. En effet, supposons que $\mathcal{N}$
soit un faisceau inversible
de degré $0 \leq d < g - 1$
possédant une section non nulle $\sigma$. Posons alors
$(\mathcal{F}, \nabla_\mathcal{F}) = (\mathcal{O}_C, \text{d})$
et
$(\mathcal{G}, \nabla_\mathcal{G}) = (\mathcal{N}^{\otimes 2}, \nabla_{can})$.
On voit que
$$
\Hom_\mathcal{A}^n((\mathcal{F}, \nabla_\mathcal{F}),
(\mathcal{G}, \nabla_\mathcal{G})) =
\left\{
\begin{matrix}
0 & \text{si} & n < 0 \\
\Gamma(C, \mathcal{N}^{\otimes 2}) & \text{si} & n = 0 \\
\Gamma(C, \mathcal{N}^{\otimes 2} \otimes \Omega_{C/k}) & \text{si} & n = 1
\end{matrix}
\right.
$$
Le premier $0$ apparaît
parce que le degré de
$\mathcal{N}^{\otimes 2} \otimes \Omega_{C/k}^{\otimes -1}$
est négatif, par l'hypothèse $d < g - 1$. Or la dérivée de la section
$\sigma^{\otimes 2}$ est nulle ; par conséquent, le groupe d'homomorphismes
$$
\Hom_{K(\mathcal{A})}((\mathcal{F}, \nabla_\mathcal{F}),
(\mathcal{G}, \nabla_\mathcal{G}))
$$
est non nul.






\section{Le champ des espaces algébriques propres n'est pas algébrique}
\sectionmark{Le champ des espaces propres n'est pas algébrique}
\label{section-proper-spaces-not-algebraic}

\noindent
Dans Quot, section \ref{quot-section-stack-of-spaces}, nous avons
introduit et étudié le champ en groupoïdes
$$
p'_{fp, flat, proper} :
\Spacesstack'_{fp, flat, proper}
\longrightarrow
\Sch_{fppf}
$$
le champ dont la catégorie des sections sur un schéma $S$
est la catégorie des espaces algébriques plats, propres et de présentation
finie sur $S$. Nous avons démontré qu'il satisfait à de nombreux
axiomes d'Artin. Dans cette section, nous montrons que ce champ
n'est pas algébrique en établissant que l'effectivité formelle
échoue en général.

\medskip\noindent
L'exemple canonique repose sur le fait que l'espace de déformation universelle
d'une variété abélienne de dimension $g$ possède $g^2$ paramètres formels,
alors que toute déformation effective peut être définie sur un
anneau local complet de dimension $\leq g(g + 1)/2$.
Nous construirons notre exemple en explicitant une
déformation non effective convenable d'une surface K3.
Nous nous bornerons à esquisser l'argument, sans en donner tous les détails.

\medskip\noindent
Soit $k = \mathbf{C}$ le corps des nombres complexes.
Soit $X \subset \mathbf{P}^3_k$ une surface lisse de degré $4$
sur $k$. On a
$\omega_X \cong \Omega^2_{X/k} \cong \mathcal{O}_X$.
Enfin,
$\dim_k H^0(X, T_{X/k}) = 0$,
$\dim_k H^1(X, T_{X/k}) = 20$ et
$\dim_k H^2(X, T_{X/k}) = 0$.
Comme $L_{X/k} = \Omega_{X/k}$ puisque $X$ est lisse sur $k$,
comme $\Ext^i_{\mathcal{O}_X}(\Omega_{X/k}, \mathcal{O}_X) =
H^i(X, T_{X/k})$, et en vertu de Cotangent, lemme
\ref{cotangent-lemma-find-obstruction-ringed-topoi},
nous obtenons une déformation universelle de $X$
sur
$$
k[[x_1, \ldots, x_{20}]]
$$
Supposons que cette déformation universelle soit effective
(au sens de Axiomes d'Artin, section \ref{artin-section-formal-objects}).
Nous obtiendrions alors un morphisme plat et propre
$$
f : Y \longrightarrow \Spec(k[[x_1, \ldots, x_{20}]])
$$
où $Y$ est un espace algébrique qui réalise la déformation universelle.
Cela est impossible pour la raison suivante. Comme $Y$ est séparé,
nous pouvons trouver un sous-schéma ouvert affine $V \subset Y$. Puisque la
fibre spéciale $X$ de $Y$ est lisse, le morphisme $f$ est lisse. Ainsi $Y$
est régulier, puisqu'il est lisse sur une base régulière, et il s'ensuit
que le complémentaire $D$ de $V$ dans $Y$ est un diviseur de Cartier effectif.
Alors $\mathcal{O}_Y(D)$ est un élément non trivial de $\Pic(Y)$
(pour le démontrer, on montre que le complémentaire d'un ouvert affine non vide
d'un espace algébrique propre et lisse sur un corps donne toujours un élément
non trivial du groupe de Picard, puis on applique ce résultat à la fibre générique
de $f$). Enfin, afin d'obtenir une contradiction, montrons que $\Pic(Y) = 0$.
L'application $\Pic(Y) \to \Pic(X)$ est injective,
car $H^1(X, \mathcal{O}_X) = 0$ (toutes les déformations de
$\mathcal{O}_X$ à $Y \times \Spec(k[[x_i]]/\mathfrak m^n)$
sont donc triviales) et par le théorème d'existence de Grothendieck
(selon lequel deux modules cohérents induisant les mêmes faisceaux
sur les épaississements
sont isomorphes). Si $X$ est suffisamment générale, alors
$\Pic(X) = \mathbf{Z}$, engendré par $\mathcal{O}_X(1)$. Il suffit donc de montrer que
$\mathcal{O}_X(n)$, pour $n > 0$, ne se déforme pas au voisinage du
premier ordre\footnote{Cet argument fonctionne dès lors que l'application
$c_1 : \Pic(X) \to H^1(X, \Omega_{X/k})$ est injective, ce qui est vrai
pour tout corps $k$ de caractéristique zéro et toute hypersurface lisse
$X$ de degré $4$ dans $\mathbf{P}^3_k$.}.
Considérons le cup-produit
$$
H^1(X, \Omega_{X/k}) \times H^1(X, T_{X/k})
\longrightarrow H^2(X, \mathcal{O}_X)
$$
Il s'agit d'un accouplement non dégénéré par dualité cohérente.
Un calcul montre que la classe de Chern
$c_1(\mathcal{O}_X(n)) \in H^1(X, \Omega_{X/k})$
est non nulle en cohomologie de Hodge.
Il existe donc une déformation du premier ordre
dont le cup-produit avec $c_1(\mathcal{O}_X(n))$ est non nul.
Enfin, on montre que ce cup-produit est précisément la classe
d'obstruction au relèvement.

\begin{lemma}
\label{lemma-proper-spaces-not-algebraic}
Le champ en groupoïdes
$$
p'_{fp, flat, proper} :
\Spacesstack'_{fp, flat, proper}
\longrightarrow
\Sch_{fppf}
$$
dont la catégorie des sections sur un schéma $S$ est la catégorie des
espaces algébriques plats, propres et de présentation finie sur $S$
(voir Quot, section \ref{quot-section-stack-of-spaces})
n'est pas un champ algébrique.
\end{lemma}

\begin{proof}
S'il s'agissait d'un champ algébrique, tout objet formel serait
effectif ; voir Axiomes d'Artin, lemme \ref{artin-lemma-effective}.
La discussion précédente montre que ce n'est pas le cas
après changement de base à $\Spec(\mathbf{C})$.
D'où la conclusion.
\end{proof}






\section{Un exemple de champ des morphismes non algébrique}
\label{section-non-algebraic-hom-stack}

\noindent
Soient $\mathcal{Y}, \mathcal{Z}$ des champs algébriques sur un schéma $S$.
Le {\it champ des morphismes} $\underline{\Mor}_S(\mathcal{Y}, \mathcal{Z})$
est le champ en groupoïdes sur $S$ dont la catégorie des sections sur
un schéma $T$ est la catégorie
$$
\Mor_T(\mathcal{Y} \times_S T, \mathcal{Z} \times_S T)
$$
dont les objets sont les $1$-morphismes et les morphismes, les $2$-morphismes.
Nous omettons la preuve qu'il s'agit bien d'un champ en groupoïdes sur
$(\Sch/S)_{fppf}$ (insérer ici une référence future). Bien entendu,
le champ des morphismes n'est pas algébrique en général.
Dans cette section,
nous en donnons un exemple où $\mathcal{Y}$ est représentable par un schéma
propre et plat sur $S$ et où $\mathcal{Z}$ est lisse et propre sur $S$.

\medskip\noindent
Soit $k$ un corps algébriquement clos qui ne soit pas la clôture
algébrique d'un corps fini. Posons $S = \Spec(k[[t]])$
et $S_n = \Spec(k[t]/(t^n)) \subset S$. Soit $f : X \to S$ un morphisme
satisfaisant aux conditions suivantes :
\begin{enumerate}
\item $f$ est projectif et plat, et ses fibres sont des courbes
géométriquement connexes,
\item la fibre $X_0 = X \times_S S_0$ est une courbe nodale à composantes
irréductibles lisses, dont le graphe dual contient un lacet formé de
courbes rationnelles,
\item $X$ est un schéma régulier.
\end{enumerate}
Pour construire une telle surface $X$, nous pouvons prendre par exemple
$$
X\quad :\quad T_0T_1T_2 - t(T_0^3 + T_1^3 + T_2^3) = 0
$$
dans $\mathbf{P}^2_{k[[t]]}$. Soit $A_0$ une variété abélienne non nulle sur $k$,
par exemple une courbe elliptique. Soit $A = A_0 \times_{\Spec(k)} S$ le
schéma abélien constant\footnote{Un schéma abélien sur une base $S$ est
un schéma en groupes plat, propre et de présentation finie sur $S$, dont
toutes les fibres sont des variétés abéliennes.}
sur $S$ associé à $A_0$. Nous montrerons que le
champ $\mathcal{X} = \underline{\Mor}_S(X, [S/A])$ n'est pas algébrique.

\medskip\noindent
Rappelons que $[S/A]$ est, d'une part, le champ quotient de l'action
triviale de $A$ sur $S$ et,
d'autre part, le champ classifiant les $A$-torseurs fppf ; voir
Exemples de champs, proposition
\ref{examples-stacks-proposition-equal-quotient-stacks}.
Observons que $[S/A] = [\Spec(k)/A_0] \times_{\Spec(k)} S$. Cela nous
permet de décrire comme suit la catégorie fibre au-dessus d'un schéma $T$ :
\begin{align*}
\mathcal{X}_T
& =
\underline{\Mor}_S(X, [S/A])_T \\
& =
\Mor_T(X \times_S T, [S/A] \times_S T) \\
& =
\Mor_S(X \times_S T, [S/A]) \\
& =
\Mor_{\Spec(k)}(X \times_S T, [\Spec(k)/A_0])
\end{align*}
pour tout $S$-schéma $T$. Autrement dit, le groupoïde $\mathcal{X}_T$
est le groupoïde des $A_0$-torseurs fppf sur $X \times_S T$.
Avant d'expliquer pourquoi $\mathcal{X}$ n'est pas un champ algébrique,
il nous faut quelques lemmes.

\begin{lemma}
\label{lemma-torsors-over-two-dimensional-regular}
Soit $W$ un schéma noethérien intègre régulier de dimension deux,
de corps des fonctions $K$. Soit $G \to W$ un schéma abélien.
Alors l'application
$H^1_{fppf}(W, G) \to H^1_{fppf}(\Spec(K), G)$ est injective.
\end{lemma}

\begin{proof}[Esquisse de la démonstration]
Soit $P \to W$ un $G$-torseur fppf, trivial au point générique.
Nous disposons alors d'un morphisme $\Spec(K) \to P$ au-dessus de $W$,
et pouvons considérer son image schématique $Z \subset P$. Comme $P \to W$
est propre (puisqu'il s'agit d'un torseur sous un espace algébrique en
groupes propre sur $W$), le morphisme $Z \to W$ est propre et birationnel.
D'après Espaces sur un corps, lemme \ref{spaces-over-fields-lemma-finite-in-codim-1}, le morphisme $Z \to W$
est fini en dehors d'un nombre fini de points fermés de $W$. D'après
(insérer ici une référence future sur la résolution des indéterminations
des morphismes de surfaces par transformations quadratiques),
les composantes irréductibles des fibres géométriques de $Z \to W$
sont des courbes rationnelles. Par
Compléments sur les groupoïdes en espaces, lemme
\ref{spaces-more-groupoids-lemma-no-nonconstant-morphism-from-P1-to-group}
il n'existe aucun morphisme non constant d'une courbe rationnelle
vers un schéma en groupes, ni vers un torseur sous un tel schéma.
Ainsi $Z \to W$ est fini ; par conséquent, $Z$ est un schéma et
$Z \to W$ est un isomorphisme d'après
Morphismes, lemme \ref{morphisms-lemma-finite-birational-over-normal}.
Autrement dit, le torseur $P$ est trivial.
\end{proof}

\begin{remark}
\label{remark-torsors-over-regular}
On peut supprimer la condition de dimension $2$ imposée à $W$ dans le
lemme \ref{lemma-torsors-over-two-dimensional-regular}. En effet, si $W$
est noethérien, régulier et intègre, alors
(1) $P(W)\to P(K)$ est bijective pour tout $G$-torseur $P$, et
(2) $H^1(W,G)\to H^1(K,G)$ est injective.
Voici une esquisse de la démonstration.
D'abord, (2) découle clairement de (1). Pour (1), l'injectivité est immédiate,
et la surjectivité revient (comme dans la preuve précédente) à montrer que
l'adhérence schématique d'une section rationnelle est une section.
Cette question est locale pour la topologie fppf sur $W$ ; après un changement
de base \'etale, nous pouvons donc supposer $P=G$ (les recouvrements \'etales
préservent la régularité). Observons ensuite que
$G=\underline{\mathrm{Pic}}^0_{G'/W}$, où $G'$ est le schéma abélien dual.
La conclusion découle du fait que $G'$ est un schéma régulier : les faisceaux
inversibles définis génériquement se prolongent. En fait, on peut remplacer
l'hypothèse « $W$ est régulier » par « les $W$-schémas lisses sont localement
factoriels ». Laurent Moret-Bailly tient cette astuce de Picard de Raynaud ;
il doit exister une référence quelque part dans les travaux de Raynaud.
\end{remark}

\begin{lemma}
\label{lemma-torsors-over-field-torsion}
Soit $G$ un espace algébrique en groupes lisse et commutatif
sur un corps $K$. Alors $H^1_{fppf}(\Spec(K), G)$ est de torsion.
\end{lemma}

\begin{proof}
Tout $G$-torseur $P$ sur $\Spec(K)$ est lisse sur $K$, puisqu'il est une forme
de $G$. Par conséquent, $P$ possède un point sur une extension finie
séparable $L/K$. Posons $[L : K] = n$. Notons $[n](P)$ le $G$-torseur
dont la classe est $n$ fois celle de $P$ dans
$H^1_{fppf}(\Spec(K), G)$. Il existe un morphisme canonique
$$
P \times_{\Spec(K)} \ldots \times_{\Spec(K)} P \to [n](P)
$$
d'espaces algébriques sur $K$. Ce morphisme est symétrique puisque
$G$ est commutatif. Il se factorise donc par le quotient
$$
(P \times_{\Spec(K)} \ldots \times_{\Spec(K)} P)/S_n
$$
D'autre part, le morphisme $\Spec(L) \to P$ définit un morphisme
$$
(\Spec(L) \times_{\Spec(K)} \ldots \times_{\Spec(K)} \Spec(L))/S_n
\longrightarrow (P \times_{\Spec(K)} \ldots \times_{\Spec(K)} P)/S_n
$$
et le lecteur vérifiera que le schéma de gauche possède un point $K$-rationnel.
Il s'ensuit que $[n](P)$ est le torseur trivial.
\end{proof}

\noindent
Pour montrer que $\mathcal{X} = \underline{\Mor}_S(X, [S/A])$
n'est pas un champ algébrique, il suffit, d'après
Axiomes d'Artin, lemme \ref{artin-lemma-effective},
d'établir le résultat suivant.

\begin{lemma}
\label{lemma-not-essentially-surjective}
L'application canonique $\mathcal{X}(S) \to \lim \mathcal{X}(S_n)$ n'est pas
essentiellement surjective.
\end{lemma}

\begin{proof}[Esquisse de la démonstration]
En déroulant les définitions, il suffit de vérifier que
$H^1(X, A_0) \to \lim H^1(X_n, A_0)$ n'est pas surjective.
Comme $X$ est régulier et projectif, les
lemmes \ref{lemma-torsors-over-field-torsion} et
\ref{lemma-torsors-over-two-dimensional-regular}
montrent que tout $A_0$-torseur sur $X$ est
de torsion. En particulier, le groupe $H^1(X, A_0)$ est de torsion.
Il suffit donc de montrer :
(a) le groupe $H^1(X_0, A_0)$ n'est pas de torsion, et
(b) les applications $H^1(X_{n + 1}, A_0) \to H^1(X_n, A_0)$ sont surjectives pour tout $n$.

\medskip\noindent
Pour (a), on construit un $A_0$-torseur $P_0$ sur la courbe nodale $X_0$
qui n'est pas de torsion, en recollant les $A_0$-torseurs triviaux sur chaque
composante de $X_0$ au moyen de points d'ordre infini de $A_0$ comme
isomorphismes au-dessus des nœuds. Plus précisément, soit $x \in X_0$ un nœud
appartenant à un lacet formé de courbes rationnelles.
Soit $X'_0 \to X_0$ la normalisation de $X_0$ dans $X_0 \setminus \{x\}$.
Soient $x', x'' \in X'_0$ les deux points au-dessus de $x$.
Nous prenons $A_0 \times_{\Spec(k)} X'_0$ et identifions
$A_0 \times \{x'\}$ à $A_0 \times \{x''\}$ par la translation
$A_0 \to A_0$ définie par un point d'ordre infini $a_0 \in A_0(k)$
(un tel point existe puisque $k$ est algébriquement clos et n'est pas la
clôture algébrique d'un corps fini --- ce fait n'est pas trivial à démontrer).
On peut montrer que le recollement est un espace algébrique (et même
un schéma) et qu'il s'agit d'un $A_0$-torseur qui n'est pas de torsion sur $X_0$.
En effet, si $[n](P_0)$ possédait une section, celle-ci fournirait un morphisme
$s : X'_0 \to A_0$ tel que
$[n](a_0) = s(x') - s(x'')$ pour la loi de groupe de $A_0(k)$. Or,
les composantes irréductibles du lacet étant rationnelles, la section $s$
y est constante (
Compléments sur les groupoïdes en espaces, lemme
\ref{spaces-more-groupoids-lemma-no-nonconstant-morphism-from-P1-to-group}).
Ainsi $s(x') = s(x'')$, ce qui donne une contradiction.

\medskip\noindent
Pour (b), la théorie des déformations montre que l'obstruction à déformer un $A_0$-torseur $P_n \to X_n$ en un $A_0$-torseur
$P_{n + 1} \to X_{n + 1}$ appartient à $H^2(X_0, \omega)$
pour un fibré vectoriel convenable $\omega$ sur $X_0$.
Ce groupe s'annule puisque $X_0$ est une courbe, ce qui prouve l'assertion.
\end{proof}

\begin{proposition}
\label{proposition-nonalghomstack}
Le champ $\mathcal{X} = \underline{\Mor}_S(X, [S/A])$ n'est pas algébrique.
\end{proposition}

\begin{proof}
Voir la discussion ci-dessus.
\end{proof}

\begin{remark}
\label{remark-contradict-aoki}
La proposition \ref{proposition-nonalghomstack} contredit
\cite[théorème 1.1]{AokiHomStacks}. La difficulté vient de la non-effectivité
des objets formels de $\underline{\Mor}_S(X, [S/A])$.
Cette même difficulté est mentionnée dans l'erratum
\cite{AokiHomStacksErr} à \cite{AokiHomStacks}. Malheureusement,
l'erratum affirme ensuite que $\underline{\Mor}_S(\mathcal{Y}, \mathcal{Z})$
est algébrique lorsque $\mathcal{Z}$ est séparé, ce qui contredit également
la proposition \ref{proposition-nonalghomstack}, puisque $[S/A]$ est séparé.
\end{remark}








\section[Champ algébrique sans effectivité formelle forte]{Un champ algébrique ne satisfaisant pas à l'effectivité formelle forte}
\label{section-non-formal-effectiveness}

\noindent
Il s'agit de \cite[exemple 4.12]{Bhatt-Algebraize}.
Soit $k$ un corps algébriquement clos.
Soit $A$ une variété abélienne sur $k$.
Supposons que $A(k)$ ne soit pas de torsion (c'est toujours le cas si $k$
n'est pas la clôture algébrique d'un corps fini).
Soit $\mathcal{X} = [\Spec(k)/A]$.
Nous affirmons qu'il existe un idéal $I \subset k[x, y]$
tel que
$$
\mathcal{X}_{\Spec(k[x, y]^\wedge)}
\longrightarrow
\lim \mathcal{X}_{\Spec(k[x, y]/I^n)}
$$
ne soit pas essentiellement surjectif. En effet, soit $I$
l'idéal engendré par $xy(x + y - 1)$.
Alors $X_0 = V(I)$ est formé de trois copies de $\mathbf{A}^1_k$
recollées en triangle en trois points. Nous pouvons donc construire un torseur
$P_0$ sous $A$ sur $X_0$, d'ordre infini, en prenant le torseur trivial
sur les composantes irréductibles de $X_0$ et en les recollant
par des translations définies par des points d'ordre infini.
Exactement comme dans la preuve du lemme \ref{lemma-not-essentially-surjective},
nous pouvons relever $P_0$ en un torseur $P_n$ sur $X_n = \Spec(k[x, y]/I^n)$.
Comme $k[x, y]^\wedge$ est un domaine régulier de dimension deux,
tout torseur $P$ sous $A$ sur $\Spec(k[x, y]^\wedge)$
est de torsion (lemmes \ref{lemma-torsors-over-two-dimensional-regular}
et \ref{lemma-torsors-over-field-torsion}). Le système de
torseurs n'appartient donc pas à l'image du foncteur affiché.

\begin{lemma}
\label{lemma-non-formal-effectiveness}
Soit $k$ un corps algébriquement clos qui ne soit pas la clôture
algébrique d'un corps fini. Soit $A$ une variété abélienne sur $k$.
Soit $\mathcal{X} = [\Spec(k)/A]$.
Il existe un système projectif de $k$-algèbres $R_n$,
à morphismes de transition surjectifs dont les noyaux sont localement nilpotents,
et un système $(\xi_n)$ d'objets de $\mathcal{X}$ au-dessus du système
$(\Spec(R_n))$, tel que ce système ne soit pas effectif
au sens de Axiomes d'Artin, remarque \ref{artin-remark-strong-effectiveness}.
\end{lemma}

\begin{proof}
Voir la discussion ci-dessus.
\end{proof}








\section{Un contre-exemple au théorème d'existence de Grothendieck}
\label{section-Grothendieck-existence}

\noindent
Soit $k$ un corps et soit $A = k[[t]]$. Soit $X$ le recollement de
$U = \Spec(A[x])$ et $V = \Spec(A[y])$ au moyen de l'identification
$$
U \setminus \{0_U\} \longrightarrow V \setminus \{0_V\}
$$
qui envoie $x$ sur $y$, où $0_U \in U$ et $0_V \in V$ sont les points
correspondant aux idéaux maximaux $(x, t)$ et $(y, t)$. Posons
$A_n = A/(t^n)$ et $X_n = X \times_{\Spec(A)} \Spec(A_n)$.
Soit $\mathcal{F}_n$ le faisceau cohérent sur $X_n$
correspondant au $A_n[x]$-module $A_n[x]/(x) \cong A_n$
et au $A_n[y]$-module $0$, avec le recollement évident.
Soit $\mathcal{I} \subset \mathcal{O}_X$ le faisceau
d'idéaux engendré par $t$. Alors $(\mathcal{F}_n)$ est un objet de
la catégorie $\textit{Coh}_{\text{support propre sur } A}(X, \mathcal{I})$
définie dans
Cohomologie des schémas, section \ref{coherent-section-existence-proper-support}.
D'autre part, cet objet n'appartient pas à l'image du
foncteur
de Cohomologie des schémas,
équation (\ref{coherent-equation-completion-functor-proper-over-A}).
En effet, s'il y appartenait, il existerait un $A[x]$-module $M$ de type fini,
un $A[y]$-module $N$ de type fini et un isomorphisme $M[1/t] \cong N[1/t]$
tels que $M/t^nM \cong A_n[x]/(x)$ et $N/t^nN = 0$ pour tout $n$.
Il est facile de voir que c'est impossible.

\begin{lemma}
\label{lemma-counter-Grothendieck-existence}
Contre-exemples à l'algébrisation des faisceaux cohérents.
\begin{enumerate}
\item Le théorème d'existence de Grothendieck, tel qu'il est énoncé dans
Cohomologie des schémas, théorème \ref{coherent-theorem-grothendieck-existence},
est faux si l'on omet l'hypothèse que $X \to \Spec(A)$ est séparé.
\item Le champ des faisceaux cohérents $\Cohstack_{X/B}$ des
théorèmes \ref{quot-theorem-coherent-algebraic-general} et
\ref{quot-theorem-coherent-algebraic} de Quot n'est, en général,
pas algébrique si l'on omet l'hypothèse que $X \to S$ est séparé.
\item Le foncteur $\Quotfunctor_{\mathcal{F}/X/B}$ de
Quot, proposition \ref{quot-proposition-quot},
n'est pas un espace algébrique en général si l'on omet l'hypothèse
que $X \to B$ est séparé.
\end{enumerate}
\end{lemma}

\begin{proof}
Nous avons vu (1) ci-dessus. Cela montre que $\textit{Coh}_{X/A}$
ne satisfait pas à l'axiome [4] de Axiomes d'Artin,
section \ref{artin-section-axioms}. Il ne peut donc pas être un champ algébrique, d'après Axiomes d'Artin, lemme
\ref{artin-lemma-effective}. On obtient ainsi (2).
Pour (3), remarquons qu'il existe des surjections compatibles
$\mathcal{O}_{X_n} \to \mathcal{F}_n$ pour tout $n$.
Ainsi $\Quotfunctor_{\mathcal{O}_X/X/A}$ ne satisfait pas à
l'axiome [4], et (3) se démontre comme précédemment.
\end{proof}






\section{Espaces algébriques formels affines}
\label{section-affine-formal-algebraic-space}

\noindent
Soit $K$ un corps et soit $(V_i)_{i \in I}$ un système projectif filtrant
d'espaces vectoriels non nuls sur $K$, à morphismes de transition surjectifs,
tel que $\lim V_i = 0$ ; voir la section \ref{section-zero-limit}.
Munissons $R_i = K \oplus V_i$ de la structure de $K$-algèbre pour laquelle
$V_i$ est un idéal de carré nul. Alors les $R_i$ forment un système projectif
de $K$-algèbres, à morphismes de transition surjectifs de noyaux nilpotents,
et $\lim R_i = K$. L'espace algébrique formel affine
$X = \colim \Spec(R_i)$ est donc un exemple d'espace
algébrique formel affine qui n'est pas McQuillan.

\begin{lemma}
\label{lemma-affine-not-mcquillan}
Il existe un espace algébrique formel affine qui n'est pas McQuillan.
\end{lemma}

\begin{proof}
Voir la discussion ci-dessus.
\end{proof}

\noindent
Soit $0 \to W_i \to V_i \to K \to 0$ un système de suites exactes
comme dans la section \ref{section-zero-limit}. Posons
$A_i = K[V_i]/(ww'; w, w' \in W_i)$.
Il existe alors un système compatible de surjections $A_i \to K[t]$
à noyaux nilpotents, et les morphismes
de transition $A_i \to A_j$ sont eux aussi surjectifs à noyaux nilpotents.
Rappelons que $V_i$ est libre sur $K$, de base les éléments $s \in S_i$.
Si la caractéristique de $K$ est nulle, la composante de degré $d$ de
$A_i$ est libre sur $K$, de base les $s^d$, $s \in S_i$, qui s'envoient
tous sur $t^d$. Le système projectif des composantes de degré $d$ des
$A_i$ est donc isomorphe au système projectif des espaces vectoriels $V_i$.
Comme $\lim V_i = 0$, nous concluons que $\lim A_i = K$, au moins lorsque
la caractéristique de $K$ est nulle.
Cela donne un exemple d'espace algébrique formel affine
dont les fonctions régulières ne séparent pas les points.

\begin{lemma}
\label{lemma-affine-formal-functions-do-not-separate-points}
Il existe un espace algébrique formel affine $X$
dont les fonctions régulières ne séparent pas les points, au sens suivant :
si l'on écrit $X = \colim X_\lambda$ comme dans
Espaces formels, définition
\ref{formal-spaces-definition-affine-formal-algebraic-space},
alors $\lim \Gamma(X_\lambda, \mathcal{O}_{X_\lambda})$
est un corps, tandis que $X_{red}$ possède une infinité de points.
\end{lemma}

\begin{proof}
Voir la discussion ci-dessus.
\end{proof}

\noindent
Soient $K$, $I$ et $(V_i)$ comme ci-dessus. Considérons les systèmes
$$
\Phi = (\Lambda, J_i \subset \Lambda, (M_i) \to (V_i))
$$
où $\Lambda$ est une $K$-algèbre augmentée,
$J_i \subset \Lambda$ est, pour $i \in I$, un idéal de carré nul,
$(M_i) \to (V_i)$ est un morphisme de systèmes projectifs de
$K$-espaces vectoriels tel que $M_i \to V_i$ soit surjectif pour tout $i$,
chaque $M_i$ possède une structure de $\Lambda$-module, les morphismes
de transition $M_i \to M_j$, pour $i > j$, sont $\Lambda$-linéaires, et
$J_j M_i \subset \Ker(M_i \to M_j)$ pour $i > j$.
Affirmation : il existe un système comme ci-dessus tel que
$M_j = M_i/J_j M_i$ pour tout $i > j$.

\medskip\noindent
Si l'affirmation est vraie, nous obtenons un morphisme représentable
$$
\colim_{i \in I} \Spec(\Lambda/J_i \oplus M_i)
\longrightarrow
\text{Spf}(\lim \Lambda/J_i)
$$
d'espaces algébriques formels affines dont la source n'est pas McQuillan,
mais dont le but l'est. Ici, $\Lambda/J_i \oplus M_i$ est muni de sa
structure usuelle de $\Lambda/J_i$-algèbre, où $M_i$ est un idéal de carré nul.
La représentabilité équivaut exactement à la condition
$M_i/J_jM_i = M_j$ pour $i > j$. La source du morphisme n'est pas
McQuillan, car les projections
$\lim_{i \in I} M_i \to M_i$ ne sont pas surjectives. Cela résulte du fait
que les applications $\lim V_i \to V_i$ ne sont pas surjectives
et que nous disposons de la surjection $M_i \to V_i$. Nous omettons certains détails.

\medskip\noindent
Démonstration de l'affirmation. Remarquons d'abord qu'il existe au moins
un système, à savoir
$$
\Phi_0 = (K, J_i = (0), (V_i) \xrightarrow{\text{id}} (V_i))
$$
Étant donné un système $\Phi$, nous montrerons qu'il existe un morphisme
de systèmes $\Phi \to \Phi'$ (les morphismes de systèmes étant définis
de manière évidente) tel que $\Ker(M_i/J_j M_i \to M_j)$ s'envoie sur
zéro dans $M'_i/J'_j M'_i$. Cela fait, nous pouvons appliquer le procédé
usuel : poser par récurrence $\Phi_n = (\Phi_{n - 1})'$ pour $n \geq 1$,
puis prendre $\Phi = \colim \Phi_n$ afin d'obtenir un système possédant
les propriétés voulues. Nous omettons les détails.

\medskip\noindent
Construction de $\Phi'$ à partir de $\Phi$. Considérons l'ensemble $U$ des triplets $u = (i, j, \xi)$ tels que $i > j$ et
$\xi \in \Ker(M_i \to M_j)$. Notons $s, t : U \to I$ les applications
$s(i, j, \xi) = i$ et $t(i, j, \xi) = j$. Posons ensuite
$\xi_u \in M_{s(u)}$ égal à la troisième composante de $u$.
Prenons
$$
\Lambda' = \Lambda[x_u; u \in U]/(x_u x_{u'}; u, u' \in U)
$$
munie de l'augmentation $\Lambda' \to K$ induite par celle de $\Lambda$
et envoyant les $x_u$ sur zéro.
Posons $J'_k = J_k \Lambda' + (x_{u,\ t(u) \geq k})$, puis
$$
M'_i = M_i \oplus \bigoplus\nolimits_{s(u) \geq i} K\epsilon_{i, u}
$$
Comme morphismes de transition $M'_i \to M'_j$ pour $i > j$, nous prenons
le morphisme donné $M_i \to M_j$ et envoyons $\epsilon_{i, u}$ sur
$\epsilon_{j, u}$. Le morphisme $M'_i \to V_i$ induit le morphisme donné
$M_i \to V_i$ et envoie $\epsilon_{i, u}$ sur zéro.
Enfin, faisons agir $\Lambda'$ sur $M'_i$ comme suit :
un élément $\lambda \in \Lambda$ agit par la structure de $\Lambda$-module
sur $M_i$ et, par l'augmentation $\Lambda \to K$, sur les $\epsilon_{i, u}$.
L'élément $x_u$ agit par $0$ sur $M_i$ pour tout $i$.
Enfin, nous posons
$$
x_u \epsilon_{i, u} = \text{image de }\xi_u\text{ dans }M_i
$$
et tous les autres produits $x_{u'} \epsilon_{i, u}$ égaux à zéro.
La formule affichée a un sens puisque $s(u) \geq i$ et
$\xi_u \in M_{s(u)}$. Les principaux points à vérifier sont que
$J'_j M'_i \subset M'_i$ s'envoie sur zéro dans $M'_j$ pour $i > j$,
et que $\Ker(M_i \to M_j)$ s'envoie sur zéro dans
$M'_i/J_j M'_i$. Ce dernier fait vient de ce que
$\xi = x_{(i, j, \xi)} \epsilon_{i, (i, j, \xi)} \in J'_j M'_i$
pour tout $\xi \in \Ker(M_i \to M_j)$.
Nous omettons les détails.

\begin{lemma}
\label{lemma-representable-morphism-affine-formal-not-mcquillan-top}
Il existe un morphisme représentable $f : X \to Y$ d'espaces
algébriques formels affines tel que $Y$ soit McQuillan, mais que $X$
ne le soit pas.
\end{lemma}

\begin{proof}
Voir la discussion ci-dessus.
\end{proof}









\section[Morphismes plats et colimites de présentation finie]{Les morphismes plats ne sont pas des colimites filtrantes de morphismes plats de présentation finie}
\label{section-flat-not-colimit-flat-finitely-presented}

\noindent
Le but de cette section est de donner un exemple de morphisme plat d'anneaux
qui ne soit pas une colimite filtrante de morphismes plats d'anneaux de
présentation finie. Dans \cite{gabber-nonexcellent}, il est montré que si
$A$ est un anneau local non excellent, de dimension $1$ et de caractéristique
résiduelle nulle, alors le morphisme plat $A \to A^\wedge$ vers son complété
n'est pas une colimite filtrante de morphismes plats d'anneaux de type fini.
Dans notre exemple, l'anneau source sera excellent. Nous invitons le lecteur
à proposer d'autres exemples ;
si vous en connaissez un, écrivez à
\href{mailto:stacks.project@gmail.com}{stacks.project@gmail.com}.

\medskip\noindent
Pour la construction, fixons un nombre premier $p$ et posons
$A = \mathbf{F}_p[x_1, \ldots, x_n]$.
Choisissons une clôture intégrale absolue $A^+$ de $A$, c'est-à-dire que $A^+$ est la
clôture intégrale de $A$ dans une clôture algébrique de son corps des fractions.
Il est montré dans \cite[\S 6.7]{HHBigCM} que $A \to A^+$ est plat.

\medskip\noindent
Nous affirmons que la $A$-algèbre $A^+$ n'est pas une colimite filtrante de
$A$-algèbres plates de présentation finie lorsque $n \geq 3$.

\medskip\noindent
Nous esquissons l'argument pour $n = 3$ et laissons au lecteur la généralisation
aux valeurs supérieures de $n$. Il suffit de prouver l'énoncé analogue pour
le morphisme $R \to R^+$, où $R$ est l'hensélisé strict de $A$ à l'origine
et $R^+$ sa clôture intégrale absolue. Remarquons que $R$
est un anneau local régulier hensélien dont le corps résiduel $k$
est une clôture algébrique de $\mathbf{F}_p$.

\medskip\noindent
Choisissons une surface abélienne ordinaire $X$ sur $k$ et un fibré inversible très ample $L$ sur $X$. L'anneau des sections
$\Gamma_*(X, L) = \bigoplus_n H^0(X,L^n)$ est l'anneau de coordonnées du cône affine sur $X$ relativement à $L$. Il est normal dès que $L$
est suffisamment positif. Notons $S$ l'hensélisé de $\Gamma_*(X, L)$
au sommet du cône. Alors $S$ est un domaine local hensélien noethérien
normal de dimension $3$. Nous obtenons un morphisme fini injectif
$R \to S$ en hensélisant une normalisation de Noether de la $k$-algèbre de type fini $\Gamma_*(X, L)$. Puisque $R^+$ est une
clôture intégrale absolue de $R$, nous pouvons aussi fixer un plongement
$S \to R^+$. Ainsi $R^+$ est également la clôture intégrale absolue de $S$.
Pour montrer que $R^+$ n'est pas une colimite filtrante de $R$-algèbres plates, il suffit d'établir les deux assertions suivantes :
\begin{enumerate}
\item Si $S \to P$ est un morphisme d'anneaux et si $P$ est plate et
de type fini sur $R$, il existe un morphisme d'anneaux $S \to T$ avec
$T$ finie et plate sur $R$.
\item Pour tout morphisme fini d'anneaux $S \to T$, l'anneau $T$
n'est pas plat sur $R$.
\end{enumerate}
En effet, comme $S$ est de présentation finie sur $R$, si l'on pouvait écrire
$R^+ = \colim_i P_i$ comme une colimite filtrante de $R$-algèbres plates
de présentation finie $P_i$, le morphisme $S \to R^+$ se factoriserait
par $S \to P_i \to R^+$ pour $i \gg 0$, contrairement aux deux assertions
précédentes. L'assertion (1) découle de l'hensélianité de $R$ et d'un
argument de sectionnement ; voir Compléments sur les morphismes, lemme
\ref{more-morphisms-lemma-qf-fp-flat-neighbourhood-dominates-fppf}.
L'assertion (2) a été démontrée dans \cite{BhattSmallCMMod} ; pour la
commodité du lecteur, nous rappelons l'argument.

\medskip\noindent
Soit $U \subset \Spec(S)$ le spectre épointé ; nous avons des morphismes naturels $X \leftarrow U \subset \Spec(S)$. Le premier fournit une
identification $H^1(U, \mathcal{O}_U) \simeq H^1(X, \mathcal{O}_X)$.
En passant aux vecteurs de Witt du perfectionné et en utilisant la suite
d'Artin--Schreier\footnote{On utilise ici que $S$ est un anneau local
strictement hensélien de caractéristique $p$ et que, par conséquent,
$S \to S$, $f \mapsto f^p - f$, est surjective.
De plus, $S$ est un domaine normal, donc
$\Gamma(U, \mathcal{O}_U) = S$. Ainsi $H^1_\etale(U, \mathbf{Z}/p)$
est le noyau de l'application
$H^1(U, \mathcal{O}_U) \to H^1(U, \mathcal{O}_U)$
induite par $f \mapsto f^p - f$.}, on obtient une identification
$H^1_\etale(U, \mathbf{Z}_p) \simeq H^1_\etale(X, \mathbf{Z}_p)$.
En particulier, ce groupe est un $\mathbf{Z}_p$-module libre de type fini
de rang $2$ (puisque $X$ est ordinaire). Pour obtenir une contradiction, supposons qu'il existe une $R$-algèbre plate $T$ comme dans (2).
Notons $V \subset \Spec(T)$ l'image réciproque de $U$, et
$f : V \to U$ le morphisme fini surjectif induit. Comme $U$ est normal,
il existe une application trace $f_*\mathbf{Z}_p \to \mathbf{Z}_p$
sur $U_\etale$, dont la composée avec l'image inverse
$\mathbf{Z}_p \to f_*\mathbf{Z}_p$ est la multiplication par $d = \deg(f)$.
En passant à la cohomologie et en utilisant que
$H^1_\etale(U, \mathbf{Z}_p)$ n'est pas de torsion, on en déduit que
$H^1_\etale(V, \mathbf{Z}_p)$ est non nul. Comme
$H^1_\etale(V, \mathbf{Z}_p) \simeq \lim H^1_\etale(V, \mathbf{Z}/p^n)$
et qu'aucun terme $R^1\lim$ n'intervient, le groupe
$H^1(V_\etale,\mathbf{Z}/p)$ est nécessairement non nul.
Puisque $T$ est plate sur $R$, on a $\Gamma(V, \mathcal{O}_V) = T$ ; cet anneau est strictement hensélien,
et la suite d'Artin--Schreier montre que $H^1(V, \mathcal{O}_V) \neq 0$.
Cela équivaut à $H^2_\mathfrak m(T) \neq 0$, où
$\mathfrak m \subset R$ est l'idéal maximal. On obtient donc une contradiction,
car $T$ est finie et plate (c'est-à-dire libre de rang fini) comme
$R$-module, et $H^2_\mathfrak m(R) = 0$. Cette contradiction prouve (2).

\begin{lemma}
\label{lemma-weird-flat-map}
Il existe un anneau commutatif $A$ et une $A$-algèbre plate $B$
qui ne peut s'écrire comme colimite filtrante de $A$-algèbres plates
de présentation finie. En fait, on peut choisir $A$ soit comme
$\mathbf{F}_p$-algèbre de type fini, soit comme anneau local noethérien
de dimension $1$, à corps résiduel de caractéristique $0$.
\end{lemma}

\begin{proof}
Voir la discussion ci-dessus.
\end{proof}




\section{La catégorie des modules modulo les modules de torsion}
\label{section-serre-quotient-modulo-torsion-modules}

\noindent
La catégorie des groupes de torsion est une sous-catégorie de Serre
(Homologie, définition \ref{homology-definition-serre-subcategory})
de la catégorie de tous les groupes abéliens. Plus généralement, pour tout
anneau $A$, la catégorie des $A$-modules de torsion est une sous-catégorie
de Serre de la catégorie de tous les $A$-modules ; voir Compléments d'algèbre, section
\ref{more-algebra-section-abelian-categories-modules}.
Si $A$ est un domaine intègre, la catégorie quotient
(Homologie, lemme \ref{homology-lemma-serre-subcategory-is-kernel})
est équivalente à la catégorie des espaces vectoriels sur le corps des fractions.
Ce fait découle de la proposition plus générale suivante.

\begin{proposition}
\label{proposition-localization-and-serre-quotients}
Soit $A$ un anneau et soit $S$ une partie multiplicative de $A$.
Notons $\text{Mod}_A$ la catégorie des $A$-modules et $\mathcal{T}$ sa sous-catégorie de Serre formée des modules dont tout élément est annulé
par un certain élément de $S$.
Il existe alors une équivalence canonique
$\text{Mod}_A/\mathcal{T} \rightarrow \text{Mod}_{S^{-1}A}$.
\end{proposition}

\begin{proof}
Le foncteur $\text{Mod}_A \to \text{Mod}_{S^{-1}A}$ défini par $M
\mapsto M \otimes_A S^{-1}A$ est exact (Algèbre, proposition
\ref{algebra-proposition-localization-exact})
et envoie les modules de $\mathcal{T}$ sur zéro.
La propriété universelle donnée par Homologie, lemme \ref{homology-lemma-serre-subcategory-is-kernel},
permet donc de le factoriser en un foncteur
$\text{Mod}_A/\mathcal{T} \to \text{Mod}_{S^{-1}A}$.

\medskip\noindent
Réciproquement, tout $A$-module $M$ tel que
$M \otimes_A S^{-1}A = 0$ appartient à $\mathcal{T}$, puisque
$M \otimes_A S^{-1}A \cong S^{-1} M$
(Algèbre, lemme \ref{algebra-lemma-tensor-localization}).
Homologie, lemme
\ref{homology-lemma-quotient-by-kernel-exact-functor}, montre donc que
le foncteur $\text{Mod}_A/\mathcal{T} \to \text{Mod}_{S^{-1}A}$ est fidèle.

\medskip\noindent
De plus, ce plongement est essentiellement surjectif : un antécédent d'un
$S^{-1}A$-module $N$ est $N_A$, c'est-à-dire $N$ considéré comme $A$-module,
car l'application canonique $N_A \otimes_A S^{-1}A \to N$, qui envoie
$x \otimes a/s$ sur $(a/s) \cdot x$, est un isomorphisme de $S^{-1}A$-modules.
\end{proof}

\begin{proposition}
\label{proposition-quotient-by-torsion-modules}
Soit $A$ un anneau. Notons $Q(A)$ son anneau total des fractions
(comme dans Algèbre, exemple \ref{algebra-example-localize-at-prime}).
Notons $\text{Mod}_A$ la catégorie des $A$-modules et $\mathcal{T}$ sa
sous-catégorie de Serre des modules de torsion. Notons
$\text{Mod}_{Q(A)}$ la catégorie des $Q(A)$-modules.
Il existe alors une équivalence canonique
$\text{Mod}_A/\mathcal{T} \rightarrow \text{Mod}_{Q(A)}$.
\end{proposition}

\begin{proof}
Cela résulte immédiatement de la proposition \ref{proposition-localization-and-serre-quotients}, appliquée à la partie
multiplicative
$S = \{f \in A \mid f \text{ n'est pas un diviseur de zéro dans }A\}$,
puisqu'un module est de torsion si et seulement si chacun de ses éléments
est annulé par un certain élément de $S$.
\end{proof}

\begin{proposition}
\label{proposition-quotient-by-finitely-generated-torsion-modules}
Soit $A$ un domaine intègre noethérien et soit $K$ son corps des fractions.
Notons $\text{Mod}_A^{fg}$ la catégorie des $A$-modules de type fini
et $\mathcal{T}^{fg}$ sa sous-catégorie de Serre des modules de torsion
de type fini. Alors $\text{Mod}_A^{fg}/\mathcal{T}^{fg}$ est canoniquement
équivalente à la catégorie des $K$-espaces vectoriels de dimension finie.
\end{proposition}

\begin{proof}
L'équivalence de la proposition
\ref{proposition-quotient-by-torsion-modules} se restreint, le long du
plongement $\text{Mod}_A^{fg}/\mathcal{T}^{fg} \to \text{Mod}_A/\mathcal{T}$,
en une équivalence
$\text{Mod}_A^{fg}/\mathcal{T}^{fg} \to \text{Vect}_K^{fd}$.
L'hypothèse noethérienne garantit que $\text{Mod}_A^{fg}$ est une
catégorie abélienne (voir Compléments d'algèbre, section
\ref{more-algebra-section-abelian-categories-modules}) et que le foncteur canonique
$\text{Mod}_A^{fg}/\mathcal{T}^{fg} \to \text{Mod}_A/\mathcal{T}$
est plein (sinon, les sous-modules de torsion des modules de type fini
pourraient ne pas être des objets de $\mathcal{T}^{fg}$).
\end{proof}

\begin{proposition}
\label{proposition-quotient-abelian-groups-by-torsion-groups}
Le quotient de la catégorie des groupes abéliens par sa
sous-catégorie de Serre des groupes de torsion est la catégorie des
$\mathbf{Q}$-espaces vectoriels.
\end{proposition}

\begin{proof}
L'assertion découle directement de la
proposition \ref{proposition-quotient-by-torsion-modules}.
\end{proof}




\section{Différentes topologies de colimite}
\label{section-colimit-topology}

\noindent
Cet exemple est \cite[exemple 1.2, page 553]{TSH}. Soit
$G_n = \mathbf{Q} \times \mathbf{R}^n$, $n \geq 1$, considéré comme groupe topologique
additif muni de la topologie usuelle (euclidienne). Considérons les
plongements fermés $G_n \to G_{n + 1}$ qui envoient
$(x_0, \ldots, x_n)$ sur $(x_0, \ldots, x_n, 0)$. Nous affirmons que
$G = \colim G_n$ muni de la topologie
$$
U \subset G\text{ ouvert} \Leftrightarrow G_n \cap U\text{ ouvert }\forall n
$$
ne définit pas un groupe topologique.

\medskip\noindent
Pour le voir, considérons l'ensemble
$$
U = \{(x_0, x_1, x_2, \ldots)\text{ tels que }
|x_j| < |\cos(jx_0)| \text{ pour } j > 0\}
$$
Comme $jx_0$ n'est jamais un multiple entier de $\pi/2$, puisque $\pi$
n'est pas rationnel, il est facile de montrer que $U \cap G_n$ est ouvert.
Puisque $0 \in U$, si la topologie ci-dessus faisait de $G$ un groupe
topologique, il existerait un voisinage ouvert $V \subset G$ de $0$
tel que $V + V \subset U$. Pour tout $j \geq 0$, il existerait alors
$\epsilon_j > 0$ tel que $(0, \ldots, 0, x_j, 0, \ldots) \in V$
dès que $|x_j| < \epsilon_j$. Puisque $V + V \subset U$, on aurait
$$
(x_0, 0, \ldots, 0, x_j, 0, \ldots) \in U
$$
pour $|x_0| < \epsilon_0$ et $|x_j| < \epsilon_j$. Or, si l'on choisit
$j$ assez grand pour que $j \epsilon_0 > \pi/2$, on peut choisir
$x_0 \in \mathbf{Q}$ de telle sorte que $|\cos(jx_0)|$ soit inférieur à
$\epsilon_j$ ; il existe alors un $x_j$ tel que
$|\cos(jx_0)| < |x_j| < \epsilon_j$. Cette contradiction prouve l'affirmation.

\begin{lemma}
\label{lemma-colimit-topology}
Il existe un système $G_1 \to G_2 \to G_3 \to \ldots$ de groupes
topologiques (abéliens) tel que $\colim G_n$, pris dans la catégorie des
espaces topologiques, diffère de $\colim G_n$, pris dans la catégorie
des groupes topologiques.
\end{lemma}

\begin{proof}
Voir la discussion ci-dessus.
\end{proof}






\section{Universellement submersif sans être un recouvrement V}
\label{section-universally-submersive-not-V}

\noindent
Soit $A$ un anneau de valuation. Soit $\mathfrak p \subset A$ un idéal
premier qui n'est ni l'idéal premier minimal ni l'idéal maximal.
(Un cas utile à garder à l'esprit est celui où $A$ possède trois idéaux
premiers et où $\mathfrak p$ est celui « du milieu ».)
Considérons le morphisme de schémas affines
$$
\Spec(A_\mathfrak p) \amalg \Spec(A/\mathfrak p) \longrightarrow \Spec(A)
$$
Nous affirmons que ce morphisme est universellement submersif. Pour le prouver,
soit $\Spec(B) \to \Spec(A)$ un morphisme de schémas affines défini par
le morphisme d'anneaux $A \to B$. Il faut montrer que
$$
\Spec(B_\mathfrak p) \amalg \Spec(B/\mathfrak pB) \to \Spec(B)
$$
soit submersif. Il est tout d'abord surjectif. Supposons ensuite que
$T \subset \Spec(B)$ soit une partie telle que
$T_1 = \Spec(B_\mathfrak p) \cap T$ et
$T_2 = \Spec(B/\mathfrak p B) \cap T$ soient fermées. Alors $T$ est l'image
du spectre d'une $B$-algèbre, car $T_1$ et $T_2$ sont tous deux des spectres
de $B$-algèbres. Pour montrer que $T$ est fermé, il suffit donc d'établir
que $T$ est stable par spécialisation ; voir
Algèbre, lemme \ref{algebra-lemma-image-stable-specialization-closed}.
Supposons que $p \leadsto q$ soit une spécialisation de points de
$\Spec(B)$, avec $p \in T$. Soit $A'$ un anneau de valuation et soit
$\Spec(A') \to \Spec(B)$ un morphisme tel que le point générique $\eta$
de $\Spec(A')$ s'envoie sur $p$ et le point fermé $s$ de $\Spec(A')$ sur $q$ ; voir
Schémas, lemme \ref{schemes-lemma-points-specialize}.
Observons que l'image du composé $\gamma : \Spec(A') \to \Spec(A)$
est exactement
l'ensemble des points
$\xi \in \Spec(A)$ avec $\gamma(\eta) \leadsto \xi \leadsto \gamma(s)$ (les détails sont omis).
Si $\mathfrak p \not \in \Im(\gamma)$, alors soit
$p, q \in \Spec(B_\mathfrak p)$, soit
$p, q \in \Spec(B/\mathfrak pB)$. Dans ce cas, puisque
$T_1$, resp.\ $T_2$, est fermée,
on a $q \in T_1$, resp.\ $q \in T_2$, et donc $q \in T$.
Supposons enfin que $\mathfrak p \in \Im(\gamma)$, disons
$\mathfrak p = \gamma(r)$. Nous avons alors les spécialisations
$p \leadsto r$ et $r \leadsto q$. Dans ce cas,
$p, r \in \Spec(B_\mathfrak p)$ et
$r, q \in \Spec(B/\mathfrak pB)$. On conclut d'abord que
$r \in T_1 \subset T$, puis que $r \in T_2$ puisque $r$ s'envoie sur
$\mathfrak p$, et enfin que $q \in T_2 \subset T$, comme voulu.

\medskip\noindent
D'autre part, nous affirmons que la famille réduite à un seul morphisme
$$
\{\Spec(A_\mathfrak p) \amalg \Spec(A/\mathfrak p) \longrightarrow \Spec(A)\}
$$
n'est pas un recouvrement V. Voir
Topologies, définition \ref{topologies-definition-V-covering}.
En effet, si elle était un recouvrement V, il existerait une extension
d'anneaux de valuation $A \subset B$ telle que
$\Spec(B) \to \Spec(A)$ se factorise par
$\Spec(A_\mathfrak p) \amalg \Spec(A/\mathfrak p)$.
Cela impliquerait que $\Spec(B)$ est non connexe, ce qui est absurde.

\begin{lemma}
\label{lemma-universally-submersive-not-V}
Il existe un morphisme $X \to Y$ de schémas affines,
universellement submersif, tel que $\{X \to Y\}$
ne soit pas un recouvrement V.
\end{lemma}

\begin{proof}
Voir la discussion ci-dessus.
\end{proof}









\section{Le spectre des entiers n'est pas quasi-compact}
\label{section-canonical}

\noindent
Bien entendu, le titre de cette section ne concerne pas le spectre des entiers
considéré comme espace topologique, car tout spectre est quasi-compact
en tant qu'espace topologique
(Algèbre, lemme \ref{algebra-lemma-quasi-compact}).
Il concerne plutôt le spectre des entiers pour la topologie canonique
sur la catégorie des schémas et la définition d'un objet quasi-compact
dans un site
(Sites, définition \ref{sites-definition-quasi-compact}).

\medskip\noindent
Soit $U$ un ultrafiltre non principal sur l'ensemble $P$ des nombres premiers.
Pour une partie $T \subset P$, notons $T^c = P \setminus T$ son complément.
Pour $A \in U$, soit $S_A \subset \mathbf{Z}$ la partie multiplicative
engendrée par les $p \in A$. Posons
$$
\mathbf{Z}_A = S_A^{-1}\mathbf{Z}
$$
Observons que $\Spec(\mathbf{Z}_A) = \{(0)\} \cup A^c \subset \Spec(\mathbf{Z})$
si l'on identifie $P$ à l'ensemble des points fermés de $\Spec(\mathbf{Z})$.
Si $A, B \in U$, alors $A \cap B \in U$ et $A \cup B \in U$
et nous avons
$$
\mathbf{Z}_{A \cap B} =
\mathbf{Z}_A \times_{\mathbf{Z}_{A \cup B}} \mathbf{Z}_B
$$
(produit fibré d'anneaux). En particulier, pour tout
entier $n$ et tous $A_1, \ldots, A_n \in U$, le morphisme
$$
\Spec(\mathbf{Z}_{A_1}) \amalg \ldots \amalg \Spec(\mathbf{Z}_{A_n})
\longrightarrow \Spec(\mathbf{Z})
$$
se factorise par $\Spec(\mathbf{Z}[1/p])$ pour un certain $p$
(en fait, pour tout $p \in A_1 \cap \ldots \cap A_n$).
Nous concluons que la famille de morphismes plats
$\{\Spec(\mathbf{Z}_A) \to \Spec(\mathbf{Z})\}_{A \in U}$
est conjointement surjective, mais qu'aucune sous-famille finie ne l'est.

\medskip\noindent
Pour un $\mathbf{Z}$-module $M$, posons
$$
M_A = S_A^{-1}M = M \otimes_{\mathbf{Z}} \mathbf{Z}_A
$$
Affirmation I : pour tout $\mathbf{Z}$-module $M$, on a
$$
M =
\text{Égaliseur}\left(
\xymatrix{
\prod\nolimits_{A \in U} M_A \ar@<1ex>[r] \ar@<-1ex>[r] &
\prod\nolimits_{A, B \in U} M_{A \cup B}
}
\right)
$$
Supposons d'abord $M$ sans torsion. Alors $M_A \subset M_P$ pour tout
$A \in U$. Il suffit donc de montrer que
$$
M = \bigcap\nolimits_{A \in U} M_A\text{ dans }M_P = M \otimes \mathbf{Q}
$$
En effet, comme $U$ n'est pas principal, pour tout nombre premier $p$, on a
$\{p\}^c \in U$. De plus, $M_{\{p\}^c} = M_{(p)}$ est la localisation
en l'idéal premier $(p)$.
L'égalité ci-dessus est donc claire, puisque $M = \bigcap_{p \in P} M_{(p)}$.
Supposons ensuite que $M$ soit de torsion. Alors
$$
M = \bigoplus\nolimits_{p \in P} M[p^\infty]
$$
et, par conséquent,
$$
M_A = \bigoplus\nolimits_{p \not \in A} M[p^\infty]
$$
puisque nous localisons aux nombres premiers appartenant à $A$.
Supposons que $(x_A) \in \prod M_A$ appartienne à l'égalisateur.
Notons $x_p = x_{\{p\}^c} \in M[p^{\infty}]$.
La propriété d'égalisation donne
$$
x_A = (x_p)_{p \not \in A}
$$
et implique en particulier que $x_p$ est nul sauf pour un nombre fini de
$p \not \in A$. Pour achever la preuve dans le cas de torsion, il suffit
de montrer que $x_p$ est nul pour presque tout nombre premier $p$.
Sinon, écrivons
$\{p \in P \mid x_p \not = 0\} = T \amalg T'$ comme réunion disjointe
de deux ensembles infinis. Alors $T \not \in U$ ou $T' \not \in U$,
puisque $U$ est un ultrafiltre (si $T, T'$ appartenaient à $U$, alors $U$ contiendrait
$T \cap T' = \emptyset$, ce qui est impossible).
Supposons $T \not \in U$. Alors $T = A^c$,
ce qui contredit la finitude établie ci-dessus.
Enfin, soit $M$ un module arbitraire. Considérons la suite exacte courte
$$
0 \to M_{tors} \to M \to M/M_{tors} \to 0
$$
et le grand diagramme suivant :
$$
\xymatrix{
M_{tors} \ar[r] \ar[d] &
\prod\nolimits_{A \in U} M_{tors, A} \ar@<1ex>[r] \ar@<-1ex>[r] \ar[d] &
\prod\nolimits_{A, B \in U} M_{tors, A \cup B} \ar[d] \\
M \ar[r] \ar[d] &
\prod\nolimits_{A \in U} M_A \ar@<1ex>[r] \ar@<-1ex>[r] \ar[d] &
\prod\nolimits_{A, B \in U} M_{A \cup B} \ar[d] \\
M/M_{tors} \ar[r] &
\prod\nolimits_{A \in U} (M/M_{tors})_A \ar@<1ex>[r] \ar@<-1ex>[r] &
\prod\nolimits_{A, B \in U} (M/M_{tors})_{A \cup B} \\
}
$$
Une chasse au diagramme, utilisant l'exactitude des colonnes ainsi que
le résultat pour le module de torsion $M_{tors}$ et pour le module
sans torsion $M/M_{tors}$, démontre l'affirmation I pour $M$.
On obtient ainsi un exemple du phénomène décrit dans le lemme suivant.

\begin{lemma}
\label{lemma-non-fpqc-descent}
Il existe un anneau $A$ et une famille infinie de morphismes plats
d'anneaux $\{A \to A_i\}_{i \in I}$ telle que, pour tout $A$-module $M$,
$$
M =
\text{Égaliseur}\left(
\xymatrix{
\prod\nolimits_{i \in I} M \otimes_A A_i \ar@<1ex>[r] \ar@<-1ex>[r] &
\prod\nolimits_{i, j \in I} M \otimes_A A_i \otimes_A A_j
}
\right)
$$
mais qu'aucune sous-famille finie ne possède cette propriété.
\end{lemma}

\begin{proof}
Voir la discussion ci-dessus.
\end{proof}

\noindent
Nous continuons avec notre ultrafiltre non principal $U$ sur l'ensemble $P$
des nombres premiers.
Soit $R$ un anneau. Notons $R_A = S_A^{-1}R = R \otimes \mathbf{Z}_A$
pour $A \in U$.
Affirmation II : étant données des parties fermées
$T_A \subset \Spec(R_A)$, $A \in U$, telles que
$$
(\Spec(R_{A \cup B}) \to \Spec(R_A))^{-1}T_A =
(\Spec(R_{A \cup B}) \to \Spec(R_B))^{-1}T_B
$$
pour tous $A, B \in U$, il existe une partie fermée $T \subset \Spec(R)$ telle que
$T_A = (\Spec(R_A) \to \Spec(R))^{-1}(T)$ pour tout $A \in U$.
Soit $I_A \subset R_A$, pour $A \in U$, l'idéal radical définissant $T_A$.
La condition de recollement implique alors
$S_{A \cup B}^{-1}I_A = S_{A \cup B}^{-1}I_B$ dans $R_{A \cup B}$ pour
tous $A, B \in U$ (car la localisation préserve le caractère radical
d'un idéal). Soit $I' \subset R$ l'ensemble des éléments
dont l'image appartient à $I_P \subset R_P = R \otimes \mathbf{Q}$.
Pour $A \in U$, on a alors :
\begin{enumerate}
\item $I_A \subset I'_A = S_A^{-1}I'$, et
\item $M_A = I'_A/I_A$ est un module de torsion.
\end{enumerate}
On obtient bien sûr des identifications canoniques
$S_{A \cup B}^{-1}M_A = S_{A \cup B}^{-1}M_B$ pour $A, B \in U$.
En décomposant les modules de torsion $M_A$ en leurs composantes
$p$-primaires, le lecteur vérifie aisément qu'il existe, pour chaque $p$, des
$R$-modules $M_p$ tels que
$$
M_A = \bigoplus\nolimits_{p \not \in A} M_p
$$
de manière compatible avec les identifications canoniques ci-dessus.
En posant $M = \bigoplus_{p \in P} M_p$, on obtient des isomorphismes
canoniques $M_A = S_A^{-1}M$, compatibles avec ces identifications.
On obtient alors une application canonique
$$
I' \longrightarrow M
$$
de $R$-modules qui redonne l'application $I'_A \to M_A$ pour tout
$A \in U$. Cela résulte de toutes les compatibilités mentionnées
ci-dessus et de l'affirmation démontrée précédemment, selon laquelle
$M$ est l'égalisateur des deux applications de
$\prod_{A \in U} M_A$ vers $\prod_{A, B \in U} M_{A \cup B}$.
Soit $I = \Ker(I' \to M)$. Alors $I$ est un idéal, et
$T = V(I)$ est une partie fermée qui redonne les parties fermées $T_A$ pour tout $A \in U$. Cela démontre l'affirmation II.

\begin{lemma}
\label{lemma-Z-not-quasi-compact}
Le schéma $\Spec(\mathbf{Z})$ n'est pas quasi-compact pour
la topologie canonique sur la catégorie des schémas.
\end{lemma}

\begin{proof}
Avec les notations précédentes, considérons la famille de morphismes
$$
\mathcal{W} = \{\Spec(\mathbf{Z}_A) \to \Spec(\mathbf{Z})\}_{A \in U}
$$
D'après Descente, lemme \ref{descent-lemma-universal-effective-epimorphism}, et les deux
affirmations démontrées ci-dessus, c'est un épimorphisme effectif
universel. Dans toute catégorie $\mathcal{C}$ possédant des produits fibrés, les
épimorphismes effectifs universels y définissent une
structure de site (à quelques difficultés ensemblistes près, faciles
à résoudre) qui définit la topologie canonique.
Ainsi, $\mathcal{W}$ est un recouvrement pour la topologie canonique.
D'autre part, nous avons vu ci-dessus que toute sous-famille finie
$$
\{\Spec(\mathbf{Z}_{A_i}) \to \Spec(\mathbf{Z})\}_{i = 1, \ldots, n},\quad
n \in \mathbf{N}, A_1, \ldots, A_n \in U
$$
se factorise par $\Spec(\mathbf{Z}[1/p])$ pour un certain $p$.
Cette famille finie ne peut donc pas être un épimorphisme effectif
universel et, plus généralement, aucun épimorphisme effectif universel
$\{g_j : T_j \to \Spec(\mathbf{Z})\}$ ne peut raffiner
$\{\Spec(\mathbf{Z}_{A_i}) \to \Spec(\mathbf{Z})\}_{i = 1, \ldots, n}$.
D'après Sites, définition \ref{sites-definition-quasi-compact},
cela signifie que $\Spec(\mathbf{Z})$ n'est pas quasi-compact
pour la topologie canonique. Pour voir que notre notion de
quasi-compacité concorde avec la définition usuelle
en théorie des topos, voir Sites, lemme \ref{sites-lemma-quasi-compact}.
\end{proof}







\begin{multicols}{2}[\section{Autres chapitres}]
\noindent
Préliminaires
\begin{enumerate}
\item \hyperref[introduction-section-phantom]{Introduction}
\item \hyperref[conventions-section-phantom]{Conventions}
\item \hyperref[sets-section-phantom]{Théorie des ensembles}
\item \hyperref[categories-section-phantom]{Catégories}
\item \hyperref[topology-section-phantom]{Topologie}
\item \hyperref[sheaves-section-phantom]{Faisceaux sur les espaces}
\item \hyperref[sites-section-phantom]{Sites et faisceaux}
\item \hyperref[stacks-section-phantom]{Champs}
\item \hyperref[fields-section-phantom]{Corps}
\item \hyperref[algebra-section-phantom]{Algèbre commutative}
\item \hyperref[brauer-section-phantom]{Groupes de Brauer}
\item \hyperref[homology-section-phantom]{Algèbre homologique}
\item \hyperref[derived-section-phantom]{Catégories dérivées}
\item \hyperref[simplicial-section-phantom]{Méthodes simpliciales}
\item \hyperref[more-algebra-section-phantom]{Compléments d'algèbre}
\item \hyperref[smoothing-section-phantom]{Lissage des morphismes d'anneaux}
\item \hyperref[modules-section-phantom]{Faisceaux de modules}
\item \hyperref[sites-modules-section-phantom]{Modules sur les sites}
\item \hyperref[injectives-section-phantom]{Objets injectifs}
\item \hyperref[cohomology-section-phantom]{Cohomologie des faisceaux}
\item \hyperref[sites-cohomology-section-phantom]{Cohomologie sur les sites}
\item \hyperref[dga-section-phantom]{Algèbre différentielle graduée}
\item \hyperref[dpa-section-phantom]{Algèbre à puissances divisées}
\item \hyperref[sdga-section-phantom]{Faisceaux différentiels gradués}
\item \hyperref[hypercovering-section-phantom]{Hyperrecouvrements}
\end{enumerate}
Schémas
\begin{enumerate}
\setcounter{enumi}{25}
\item \hyperref[schemes-section-phantom]{Schémas}
\item \hyperref[constructions-section-phantom]{Constructions de schémas}
\item \hyperref[properties-section-phantom]{Propriétés des schémas}
\item \hyperref[morphisms-section-phantom]{Morphismes de schémas}
\item \hyperref[coherent-section-phantom]{Cohomologie des schémas}
\item \hyperref[divisors-section-phantom]{Diviseurs}
\item \hyperref[limits-section-phantom]{Limites de schémas}
\item \hyperref[varieties-section-phantom]{Variétés}
\item \hyperref[topologies-section-phantom]{Topologies sur les schémas}
\item \hyperref[descent-section-phantom]{Descente}
\item \hyperref[perfect-section-phantom]{Catégories dérivées de schémas}
\item \hyperref[more-morphisms-section-phantom]{Compléments sur les morphismes}
\item \hyperref[flat-section-phantom]{Compléments sur la platitude}
\item \hyperref[groupoids-section-phantom]{Schémas en groupoïdes}
\item \hyperref[more-groupoids-section-phantom]{Compléments sur les schémas en groupoïdes}
\item \hyperref[etale-section-phantom]{Morphismes étales de schémas}
\end{enumerate}
Thèmes de théorie des schémas
\begin{enumerate}
\setcounter{enumi}{41}
\item \hyperref[chow-section-phantom]{Homologie de Chow}
\item \hyperref[intersection-section-phantom]{Théorie de l'intersection}
\item \hyperref[pic-section-phantom]{Schémas de Picard des courbes}
\item \hyperref[weil-section-phantom]{Théories cohomologiques de Weil}
\item \hyperref[adequate-section-phantom]{Modules adéquats}
\item \hyperref[dualizing-section-phantom]{Complexes dualisants}
\item \hyperref[duality-section-phantom]{Dualité pour les schémas}
\item \hyperref[discriminant-section-phantom]{Discriminants et différentes}
\item \hyperref[derham-section-phantom]{Cohomologie de de Rham}
\item \hyperref[local-cohomology-section-phantom]{Cohomologie locale}
\item \hyperref[algebraization-section-phantom]{Géométrie algébrique et formelle}
\item \hyperref[curves-section-phantom]{Courbes algébriques}
\item \hyperref[resolve-section-phantom]{Résolution des surfaces}
\item \hyperref[models-section-phantom]{Réduction semi-stable}
\item \hyperref[functors-section-phantom]{Foncteurs et morphismes}
\item \hyperref[equiv-section-phantom]{Catégories dérivées de variétés}
\item \hyperref[pione-section-phantom]{Groupes fondamentaux de schémas}
\item \hyperref[etale-cohomology-section-phantom]{Cohomologie étale}
\item \hyperref[crystalline-section-phantom]{Cohomologie cristalline}
\item \hyperref[proetale-section-phantom]{Cohomologie pro-étale}
\item \hyperref[relative-cycles-section-phantom]{Cycles relatifs}
\item \hyperref[more-etale-section-phantom]{Compléments de cohomologie étale}
\item \hyperref[trace-section-phantom]{La formule des traces}
\end{enumerate}
Espaces algébriques
\begin{enumerate}
\setcounter{enumi}{64}
\item \hyperref[spaces-section-phantom]{Espaces algébriques}
\item \hyperref[spaces-properties-section-phantom]{Propriétés des espaces algébriques}
\item \hyperref[spaces-morphisms-section-phantom]{Morphismes d'espaces algébriques}
\item \hyperref[decent-spaces-section-phantom]{Espaces algébriques décents}
\item \hyperref[spaces-cohomology-section-phantom]{Cohomologie des espaces algébriques}
\item \hyperref[spaces-limits-section-phantom]{Limites d'espaces algébriques}
\item \hyperref[spaces-divisors-section-phantom]{Diviseurs sur les espaces algébriques}
\item \hyperref[spaces-over-fields-section-phantom]{Espaces algébriques sur des corps}
\item \hyperref[spaces-topologies-section-phantom]{Topologies sur les espaces algébriques}
\item \hyperref[spaces-descent-section-phantom]{Descente et espaces algébriques}
\item \hyperref[spaces-perfect-section-phantom]{Catégories dérivées d'espaces}
\item \hyperref[spaces-more-morphisms-section-phantom]{Compléments sur les morphismes d'espaces}
\item \hyperref[spaces-flat-section-phantom]{Platitude sur les espaces algébriques}
\item \hyperref[spaces-groupoids-section-phantom]{Groupoïdes dans les espaces algébriques}
\item \hyperref[spaces-more-groupoids-section-phantom]{Compléments sur les groupoïdes dans les espaces}
\item \hyperref[bootstrap-section-phantom]{Amorçage}
\item \hyperref[spaces-pushouts-section-phantom]{Sommes amalgamées d'espaces algébriques}
\end{enumerate}
Thèmes de géométrie
\begin{enumerate}
\setcounter{enumi}{81}
\item \hyperref[spaces-chow-section-phantom]{Groupes de Chow des espaces}
\item \hyperref[groupoids-quotients-section-phantom]{Quotients de groupoïdes}
\item \hyperref[spaces-more-cohomology-section-phantom]{Compléments sur la cohomologie des espaces}
\item \hyperref[spaces-simplicial-section-phantom]{Espaces simpliciaux}
\item \hyperref[spaces-duality-section-phantom]{Dualité pour les espaces}
\item \hyperref[formal-spaces-section-phantom]{Espaces algébriques formels}
\item \hyperref[restricted-section-phantom]{Algébrisation des espaces formels}
\item \hyperref[spaces-resolve-section-phantom]{Résolution des surfaces revisitée}
\end{enumerate}
Théorie des déformations
\begin{enumerate}
\setcounter{enumi}{89}
\item \hyperref[formal-defos-section-phantom]{Théorie formelle des déformations}
\item \hyperref[defos-section-phantom]{Théorie des déformations}
\item \hyperref[cotangent-section-phantom]{Le complexe cotangent}
\item \hyperref[examples-defos-section-phantom]{Problèmes de déformation}
\end{enumerate}
Champs algébriques
\begin{enumerate}
\setcounter{enumi}{93}
\item \hyperref[algebraic-section-phantom]{Champs algébriques}
\item \hyperref[examples-stacks-section-phantom]{Exemples de champs}
\item \hyperref[stacks-sheaves-section-phantom]{Faisceaux sur les champs algébriques}
\item \hyperref[criteria-section-phantom]{Critères de représentabilité}
\item \hyperref[artin-section-phantom]{Axiomes d'Artin}
\item \hyperref[quot-section-phantom]{Espaces de Quot et de Hilbert}
\item \hyperref[stacks-properties-section-phantom]{Propriétés des champs algébriques}
\item \hyperref[stacks-morphisms-section-phantom]{Morphismes de champs algébriques}
\item \hyperref[stacks-limits-section-phantom]{Limites de champs algébriques}
\item \hyperref[stacks-cohomology-section-phantom]{Cohomologie des champs algébriques}
\item \hyperref[stacks-perfect-section-phantom]{Catégories dérivées de champs}
\item \hyperref[stacks-introduction-section-phantom]{Introduction aux champs algébriques}
\item \hyperref[stacks-more-morphisms-section-phantom]{Compléments sur les morphismes de champs}
\item \hyperref[stacks-geometry-section-phantom]{La géométrie des champs}
\end{enumerate}
Thèmes de théorie des espaces de modules
\begin{enumerate}
\setcounter{enumi}{107}
\item \hyperref[moduli-section-phantom]{Champs de modules}
\item \hyperref[moduli-curves-section-phantom]{Espaces de modules de courbes}
\end{enumerate}
Divers
\begin{enumerate}
\setcounter{enumi}{109}
\item \hyperref[examples-section-phantom]{Exemples}
\item \hyperref[exercises-section-phantom]{Exercices}
\item \hyperref[guide-section-phantom]{Guide bibliographique}
\item \hyperref[desirables-section-phantom]{Desiderata}
\item \hyperref[coding-section-phantom]{Style de programmation}
\item \hyperref[obsolete-section-phantom]{Obsolète}
\item \hyperref[fdl-section-phantom]{Licence de documentation libre GNU}
\item \hyperref[index-section-phantom]{Index généré automatiquement}
\end{enumerate}
\end{multicols}


\bibliography{my}
\bibliographystyle{amsalpha}

\end{document}
